\documentclass[12pt, a4paper]{report}
\usepackage{graphicx} % For including images
\usepackage{svg} % For \includesvg
\usepackage{titlesec} % For customizing section titles
\usepackage{tocloft} % For customizing table of contents
\usepackage{acro} % For acronyms
\usepackage{array}
\usepackage{pgfgantt}
\usepackage{longtable}
\usepackage{makecell}
\usepackage{caption}


%% ____Bibliography____%%
\usepackage[numbers,sort&compress]{natbib}
\usepackage{chapterbib}
\usepackage{amsmath}
\usepackage{pgfgantt}
\usepackage{soul} % For \hl highlighting
\usepackage{xcolor}
\usepackage{float}
\usepackage{amssymb} % for \checkmark
%\hypersetup{colorlinks=true,citecolor=blue,linkcolor=blue,urlcolor=blue}
% Page margins
\usepackage[left=1.5in, right=1in, top=1in, bottom=1in]{geometry}
\usepackage[breaklinks]{hyperref}
\hypersetup{
    colorlinks=true,
    citecolor=blue,
    linkcolor=blue,
    urlcolor=blue
}
\usepackage{xcolor}     % for \rowcolor
 \usepackage{colortbl}   % for row shading
 \usepackage{makecell}   % for \makecell and \thead
\definecolor{gapgray}{gray}{0.92}

\usepackage{comment}
\usepackage{algorithm}
% \usepackage{algorithmic}
\usepackage{algpseudocode}
\usepackage{soul}    % in preamble
% Remove page number from the first page
\thispagestyle{empty}

\usepackage{longtable}
\usepackage{array}

% Customize table of contents, list of figures, and list of tables
\renewcommand{\cfttoctitlefont}{\hfill\Large\bfseries}
\renewcommand{\cftaftertoctitle}{\hfill}
\renewcommand{\cftloftitlefont}{\hfill\Large\bfseries}
\renewcommand{\cftafterloftitle}{\hfill}
\renewcommand{\cftlottitlefont}{\hfill\Large\bfseries}
\renewcommand{\cftafterlottitle}{\hfill}


% Define acronyms
\DeclareAcronym{AI}{
  short = AI,
  long  = Artificial Intelligence
}
\DeclareAcronym{ML}{
  short = ML,
  long  = Machine Learning
}
\DeclareAcronym{SDN}{
  short = SDN,
  long  = Software-Defined Networking
}

\DeclareAcronym{SDVN}{
  short = SDVN,
  long  = Software-Defined Vehicular Network
}

\DeclareAcronym{VANET}{
  short = VANET,
  long  = Vehicular Ad Hoc Network
}

\DeclareAcronym{TCAM}{
  short = TCAM,
  long  = Ternary Content-Addressable Memory
}

\DeclareAcronym{SIFT}{
  short = SIFT,
  long  = SelectIve DeFense for TCAM
}

\DeclareAcronym{TAP}{
  short = TAP,
  long  = Timing Attack Prevention
}

\DeclareAcronym{FSDM}{
  short = FSDM,
  long  = Fast recovery Saturation attack Detection and Mitigation
}

\DeclareAcronym{V2I}{
  short = {V2I},
  long = {vehicle-to-infrastructure}
}
\DeclareAcronym{OLS}{
  short = {OLS},
  long = {ordinary least squares}
}
  
\DeclareAcronym{UCR}{
  short = {UCR},
  long  = {Unauthorized Copy Rate}

}
 

\DeclareAcronym{RSU}{
  short = RSU,
  long  = Roadside Unit
}

\DeclareAcronym{TPM}{
  short = TPM,
  long  = Trusted Platform Modules
}

\DeclareAcronym{DAA}{
  short = DAA,
  long  = Direct Anonymous Attestation
}

\DeclareAcronym{FL}{
  short = FL,
  long  = Federated Learning
}

\DeclareAcronym{HSA}{
  short = HSA,
  long  = Header Space Analysis
}

\DeclareAcronym{LSTM}{
  short = LSTM,
  long  = Long Short-Term Memory
}

\DeclareAcronym{SD-NDN}{
  short = SD-NDN,
  long  = Software Defined Named Data Network 
}

\DeclareAcronym{MANET}{
  short = MANET,
  long  = Moblie Ad-hoc Networks 
}

\DeclareAcronym{OBU}{
  short = OBU,
  long  = On Board Units
}
% --- Missing from Acronym List ---

\DeclareAcronym{ZKP}{
  short = ZKP,
  long  = Zero-Knowledge Proof
}

\DeclareAcronym{SUMO}{
  short = SUMO,
  long  = Simulation of Urban Mobility
}

\DeclareAcronym{PBFT}{
  short = PBFT,
  long  = Practical Byzantine Fault Tolerance
}


\DeclareAcronym{DSCP}{
  short = DSCP,
  long  = Differentiated Services Code Point
}

\DeclareAcronym{BSM}{
  short = BSM,
  long  = Basic Safety Message
}

\DeclareAcronym{HMAC}{
  short = HMAC,
  long  = Hash-based Message Authentication Code
}

\DeclareAcronym{EWMA}{
  short = EWMA,
  long  = Exponentially Weighted Moving Average
}

\DeclareAcronym{V2V}{
  short = V2V,
  long  = Vehicle-to-Vehicle
}

\DeclareAcronym{V2X}{
  short = V2X,
  long  = Vehicle-to-Everything
}

\DeclareAcronym{FedAvg}{
  short = FedAvg,
  long  = Federated Averaging
}

\DeclareAcronym{LLDP}{
  short = LLDP,
  long  = Link Layer Discovery Protocol
}

\DeclareAcronym{NS3}{
  short = NS-3,
  long  = Network Simulator 3
}

\DeclareAcronym{MCC}{
  short = MCC,
  long  = Matthews Correlation Coefficient
}

\DeclareAcronym{PDR}{
  short = PDR,
  long  = Packet Delivery Ratio
}

\DeclareAcronym{FPR}{
  short = FPR,
  long  = False Positive Rate
}

\DeclareAcronym{PKI}{
  short = PKI,
  long  = Public Key Infrastructure
}

\DeclareAcronym{SHA}{
  short = SHA-256,
  long  = Secure Hash Algorithm 256-bit
}

\DeclareAcronym{DKG}{
  short = DKG,
  long  = Distributed Key Generation
}
\DeclareAcronym{MLDSA}{
  short = ML-DSA-87,
  long  = {Module Lattice-based Digital Signature Algorithm (FIPS~204, Security Level~5)}
}
\DeclareAcronym{STARK}{
  short = STARK,
  long  = Scalable Transparent ARgument of Knowledge
}
\DeclareAcronym{SHAKE}{
  short = SHA3-512,
  long  = Secure Hash Algorithm 3, 512-bit
}

\DeclareAcronym{FlowMod}{
  short = FlowMod,
  long  = Flow Modification
}

\DeclareAcronym{BFT}{
  short = BFT,
  long  = Byzantine Fault Tolerance
}

\DeclareAcronym{FTISCON}{
  short = FTISCON,
  long  = Flow Table Integrity Through Smart Contract
}
\DeclareAcronym{BlockREV}{
  short = BlockREV,
  long  =  Blockchain-Enabled Multi-Controller Rule Enforcement Verification
}
\DeclareAcronym{GPS}{
  short = GPS,
  long  = Global Positioning System,
}
\begin{document}

\include{cover}

\renewcommand{\thepage}{\roman{page}} % Start page numbering in roman


\chapter*{Abstract}



% Software defined vehicular networks (SDVNs) are central to next generation
% intelligent transportation systems, yet their programmable and logically
% centralized architecture introduces critical security vulnerabilities.
% Selective time delay and hidden forwarding attacks are uniquely dangerous
% in this context: rather than disrupting network operation overtly, these
% attacks selectively delay safety-critical packets such as collision
% avoidance warnings to cause system desynchronization, or secretly duplicate
% and reroute sensitive data through covert channels to enable eavesdropping
% and misinformation injection---all while mimicking legitimate traffic to
% bypass volume-based detection. 
The centralized programmability of software-defined vehicular networks
(SDVNs) introduces critical vulnerabilities to selective time delay and
hidden forwarding attacks: rather than disrupting network operation
overtly, these attacks selectively delay safety-critical packets such as
collision avoidance warnings to cause system desynchronization, or
secretly duplicate and reroute sensitive data through covert channels to
enable eavesdropping and misinformation injection---all while mimicking
legitimate traffic to bypass volume-based detection.
Existing defenses are fundamentally
inadequate: centralized frameworks such as SIFT and FSDM assume a trusted
controller and fail catastrophically when the controller itself is
compromised, creating a single point of failure with no immutable audit
trail; path-validation approaches such as HSA centralize raw packet
metadata collection, violating driver privacy by exposing location and
movement patterns; and timing-based protocols such as TAP rely on static
detection thresholds that are easily evaded and cannot detect hidden
forwarding. Critically, no existing work accounts for how vehicular
mobility---handoff-induced latency jitter, rapid topology changes, and
short RSU observation windows---masks attack signatures, making all
existing SDN defenses ineffective under realistic vehicular conditions.

To address these gaps, this paper presents MOBIGUARD, a mobility-aware,
zero-trust defense framework that assumes no inherent trustworthiness in
any network element, including controllers, RSUs, or vehicles. MOBIGUARD
formally models eight distinct attack variants across control-plane and
data-plane adversary positions and defends against all eight through three
integrated layers. First, a dual-mode detection engine pairs a lightweight
rule-based signature detector on resource-constrained on-board units with
a full federated LSTM anomaly detector at RSU edge nodes, enabling
privacy-preserving collaborative detection by sharing only model weights.
Second, a hybrid cryptographic layer combines per-packet ML-DSA-87 signatures,
aggregate signature batch verification, and
\ac{STARK} zero-knowledge proofs for timing
compliance and next-hop legitimacy---preventing
spoofing, timestamp manipulation, and covert
forwarding while preserving vehicle identity privacy.
Third, a two-tier permissioned blockchain with a flat
multi-controller zero-trust architecture and
distributed key generation enforces tamper-proof
audit, automated quarantine, and Byzantine-robust
federated model validation, ensuring integrity even
under controller compromise and eliminating the
single point of failure present in all existing
centralized frameworks. Passive hidden forwarding
variants undetectable by cryptographic proof alone
are further addressed by a broadcast-radio
witness-monitoring mechanism with Byzantine
fault-tolerant alert consensus.

% Evaluation under
% realistic SUMO mobility traces demonstrates ${\geq}96\%$ detection
% accuracy, ${\leq}1\%$ false positive rate, and end-to-end latency below
% 85\,ms across all eight attack variants, satisfying the strict
% requirements of safety-critical vehicular applications.
MOBIGUARD is evaluated across twelve performance
metrics under five benchmarking experiments---varying
attack penetration, mobility speed from 10 to
140\,km/h, network scale from 100 to 400 vehicles,
and attack stealthiness---alongside thirteen ablation
experiments isolating each security component's
individual contribution.



% Table of Contents
\tableofcontents
\newpage

% List of Figures
\listoffigures
\newpage

% List of Tables
\listoftables
\newpage

% Acronyms
\addcontentsline{toc}{chapter}{Acronyms} % Add to table of contents
\acuseall % Use all acronyms to ensure they appear in the list
\printacronyms
\newpage

\renewcommand{\thepage}{\arabic{page}} % Start page numbering in arabic 
\setcounter{page}{1} % start page numbering from 1

% Main Content
\chapter{Introduction}
\section{Background}

% --- Part 1: Evolution of Vehicular Networks (Context) ---
A Vehicular Ad-hoc Network (VANET) is a specialized subclass of Mobile Ad-hoc Networks (MANETs) where mobile nodes consist of vehicles equipped with On-Board Units (OBUs) and there are  stationary Roadside Units (RSUs)\cite{Cardona_2020}\cite{s23031600}. These networks are highly mobile, and they have fast-changing topologies, and a highly unstable communication environment\cite{Cardona_2020}\cite{s23031600}. Despite their potential for improving road safety and traffic efficiency, traditional VANETs struggle with inability to provide system programmability and efficent information distribution\cite{KALUPAHANALIYANAGE2018226}\cite{s23031600}.

Software Defined Networking (SDN) introduced as a paradigm that decouples the control plane (the decision-making logic) from the data plane (the forwarding hardware)\cite{PASCOAL2020107223}. This separation allows a logically centralized controller to manage network resources through software-based programming, facilitating better configuration, automation, and monitoring\cite{PASCOAL2020107223}.

A Software Defined Vehicular Network (SDVN) integrates Software Defined Network into the vehicular domain to overcome the inherent limitations of traditional VANETs\cite{KALUPAHANALIYANAGE2018226}\cite{s23031600}. It provides bird’s eye view of the complete network, the SDVN controller can gather information on network to improve routing, load balancing, and overall network awareness.

% --- Part 2: The Specific Threat Landscape (Problem) ---
Despite these operational advantages, switching to SDVN poses significant security risks. SDVN's centralized control and programmability provide novel exploit vectors, such as Selective Time Delay and Hidden Forwarding attacks\cite{9088232}. These threats are particularly dangerous as they are hard to detect by existing volume based defences as they do not drastically change network volume. They mimic legitimate traffic or operate through hidden channels, which may lead to fatal accidents in time-critical environments\cite{PASCOAL2020107223}.

\section{Motivation}

These vulnerabilities in SDVN are critical because SDVN relies on a logically centralized controller to push dynamic flow rules to the data plane, making the entire system inherently sensitive to timing\cite{8187644}\cite{7981537}. SDVN has strict latency requirements because it supports safety-critical applications (typically $<100\,\mathrm{ms}$), making the entire system highly sensitive to timing; any timing manipulation can be potentially lethal.\cite{Timing_Attack}\cite{KALUPAHANALIYANAGE2018226}
Unlike traditional DDoS attacks, which are easily detected by volume-based monitoring tools, selective time-delay attacks allow the network to appear functional while silently delaying safety-critical information \cite{PASCOAL2020107223}\cite{DBLP:journals/air/QammarKND23}.

Parallel to timing manipulations, the unauthorized redirection secure information can pose severe threat to the privacy and security of the SDVN data plane. Adversaries exploits the programmable nature of SDVN and manipulating flow rules, establishing covert channels to secretly duplicate and reroute sensitive packets, such as location history or collison warning. This capability facilitates eavesdropping, allowing unauthorized entities to intercept confidential data without disrupting the primary flow, often bypassing traditional detection systems\cite{mohammadi2016detecting}. Furthermore, such vulnerabilities enables attackers inject malicious payloads or alter cooperative perception streams and distribute faulty packets in the network. The acceptance of this corrupted data is critically dangerous as  in autonomous driving environments, as the misinformation distributions can compromise decision-making algorithms, potentially leading to severe road accidents\cite{Cardona_2020}\cite{8616999}.

\section{Problem Introduction}

Selective Time Delay attacks involve a compromised entity, such as a malicious SDN controller, RSU, or a vehicle-intentionally introducing variable delays into high-priority safety packets \cite{Timing_Attack} instead of dropping them\cite{Timing_Attack}. This leads to the network desynchronization, causing systems to react based on invalid data, which can result in accidents\cite{Timing_Attack}. Simultaneously, In Hidden Forwarding attacks duplicate packets will be sent to unauthorized destinations using hidden channels. This attack will send original packet to legitimate destination in orderly manner while sending manipulated or unmanipulated packets to malicious destination, to enable eavesdropping on critical and private data or distribute misinformation without triggering standard security monitors.

\section{Limitations of Exsisting Approches}
While there are various exsisting  mitigation strategies for SDN and VANET security, they may not precisely address the specific challenges of SDVNs. First, more reliable frameworks such as SIFT \cite{pascoal2017slow} and FSDM \cite{9356082} rely on logically centralized architectures that operate under an ``assumed trust'' model \cite{electronics12143077}. If the controller itself is malicious the results will be fatal as this creates single point of faliure. These systems lack the immutable auditing to detect malicious rule manipulation \cite{jaballah2019survey}. Second, existing path-validation and anomaly detection methods, such as HSA \cite{mohammadi2016detecting}, typically require the centralized collection of raw metadata. This mass data collection poses severe privacy risks, potentially exposing sensitive user information like geographic vehical location and driving patterns \cite{11059879}. Existing defense mechanisms are focused on either timing-related anomalies or routing attacks separately. There is no integrated, multi-layered framework that can jointly mitigate Selective Time Delay and Hidden Forwarding attacks while maintaining the low-latency requirements of highly mobile vehicular networks.

\section{Problem Statement}

In software-defined vehicular networks, a compromised entity---whether a
malicious SDN controller, RSU, or vehicle---can selectively delay
safety-critical packets such as collision avoidance warnings to cause
network desynchronization, forcing automated systems to react based on
stale and invalid data, or can simultaneously forward original packets
to their legitimate destinations while secretly duplicating and rerouting
them to unauthorized destinations to enable eavesdropping on critical
private data and misinformation distribution---all without alerting the
system. Under high-mobility vehicular conditions, where handoff-induced
latency jitter and rapid topology changes statistically mask these
manipulations, no existing security framework can robustly detect both
attack classes simultaneously without either assuming a trusted
centralized controller that itself may be the adversary, or resorting
to mass ingestion of raw vehicular data that violates driver privacy
and exceeds the strict latency constraints of safety-critical
vehicular applications.

\section{Objectives and Scope}

\subsection{Objectives}

This project aims to build multilayerd secure SDVN framework against selective time delay and hidden forwarding attacks.

Objective 1: Design and simulate realistic Selective Time Delay and Hidden Forwarding attack scenarios within an SDVN environment using NS-3.


Objective 2: Develop a multi-layered security architecture that integrates blockchain for tamper-proof logging, cryptographic protocols for identity verification, and Federated LSTM models for decentralized anomaly detection.


Objective 3: Evaluate MOBIGUARD against three
state-of-the-art baselines across twelve performance
metrics spanning detection quality, attack-specific
containment (TVR, UCR), overhead, and privacy, using
five benchmarking experiments and thirteen ablation
experiments that isolate each component's contribution.



\subsection{Research Questions}
\label{sec:rqs}

\begin{enumerate}
    
    \item \textbf{RQ1:}
    How does MOBIGUARD's detection quality, threshold
    violation containment, and mitigation responsiveness
    change as attack penetration increases from 0\% to
    100\% of network nodes across three attack intensity
    levels, compared against state-of-the-art defenses
    that assume a fixed or single-adversary threat model?
    
    \item \textbf{RQ2:}
    Does the mobility-adjusted threshold design maintain
    detection quality and safety-critical packet
    protection across the full vehicular speed range
    from 10\,km/h to 140\,km/h, where increasing speed
    progressively raises handoff-induced latency jitter
    that masks attack signatures and causes
    static-threshold baselines to fail?
    
    \item \textbf{RQ3:}
    Can the \ac{MLDSA} + randomized batch verification
    + \ac{STARK} cryptographic layer and blockchain
    consensus pipeline maintain detection quality,
    mitigation latency, and end-to-end latency within
    the 100\,ms safety-critical bound as the network
    scales from 100 to 400 vehicles?
    
    \item \textbf{RQ4:}
    How does MOBIGUARD's detection quality and
    unauthorized copy containment degrade as attack
    observable evidence decreases from maximum
    (high selectivity, adjacent copy destination)
    to minimum (10\% selectivity, near-maximum
    divergence copy destination), and does the
    witness-based monitoring mechanism maintain
    passive hidden forwarding detection at low
    evidence levels where cryptographic proof
    corroboration alone is insufficient?
    
    \item \textbf{RQ5:}
    Across each of the eight individual attack variants
    under default settings, does MOBIGUARD achieve
    superior detection quality and mitigation performance
    compared to state-of-the-art defenses, and what
    detection gaps remain in the coverage of existing
    baselines that a joint detection framework resolves?

\end{enumerate}


\vspace{0.5\baselineskip}

\subsection{Scope}

\noindent Network Environment: The research will be focused on Software-Defined Vehicular Networks (SDVN). It will not cover traditional VANETs or standard SDN architecture.\\[\baselineskip]
\noindent Attack Vectors: We have defined attacks and evaluations will be based on selected Hidden forwarding and Selective Time delay attacks.\\[\baselineskip]
\noindent Perfomance Analysis: Perfomance Analysis will be limited to simulations using NS-3.


\section{Contributions and Novelty}


\subsection{Contributions}
\label{sec:contributions}

The principal contributions of this paper are:
\begin{itemize}
    \item \textbf{Mobility-specific attack modeling:} We formally define 
    eight attack variants (four selective time delay, four hidden 
    forwarding) and analyze how vehicular mobility---handoff masking, 
    topology churn, short observation windows, flow-rule 
    camouflage---amplifies each variant beyond what is possible in 
    static SDN.
    
    \item \textbf{Dual-mode detection architecture:} We propose a 
    two-tier detection system comprising (i)~a lightweight, rule-based 
    detection engine deployable on resource-constrained OBUs that 
    identifies known attack signatures using mobility-adjusted timing 
    thresholds, and (ii)~a full federated LSTM-based anomaly detector 
    at RSU edge nodes for detecting novel and sophisticated attack 
    patterns, with adaptive mode switching based on node capability 
    and threat level.
    
    \item \textbf{Attack signature identification:} We derive and 
    validate distinguishing signatures for all eight attack variants, 
    including handoff-normalized delay deviation, asymmetric flow 
    conservation residuals, and TCAM installation rate anomalies 
    relative to predicted vehicle density.
    
    \item \textbf{Hybrid cryptographic integrity scheme:} We design 
    an ML-DSA-87 per-packet integrity tag combined with aggregate 
    signatures for efficient batch verification at RSUs, augmented 
    with zero-knowledge proofs for privacy-preserving flow-rule 
    audit verification---ensuring packet authenticity without 
    exposing vehicle identity or trajectory.
    
    \item \textbf{Federated LSTM for privacy-preserving detection:} 
    We deploy LSTM-based anomaly detection across distributed RSU 
    nodes using federated learning, enabling collaborative model 
    training without centralizing sensitive flow data while 
    maintaining detection accuracy comparable to centralized 
    approaches.
    
    \item \textbf{Blockchain-based audit trail:} We implement a 
    lightweight permissioned blockchain for tamper-proof logging of 
    flow-rule installations, detection events, and remediation 
    actions, providing decentralized accountability and eliminating 
    the single-point-of-failure inherent in centralized SDN 
    controllers.

    \item \textbf{Multi-controller zero-trust architecture:}
We design a flat registered controller set with
blockchain-enforced trust scoring, smart
contract-triggered failover (SC-Revoke), and a
\ac{DKG} ceremony among \ac{RSU}s that removes the
controller from the key plane entirely, ensuring
control-plane continuity and cryptographic integrity
even under controller compromise.

\item \textbf{Witness-based passive attack detection:}
We introduce a broadcast-radio witness monitoring
mechanism that detects passive hidden forwarding
variants (S7, S8) by collecting independently
observed packet hashes from neighboring vehicles,
triggering a smart contract trust penalty only after
a $2f{+}1$ Byzantine fault-tolerant threshold of
verified alerts is reached---providing detection that
cryptographic proof corroboration alone cannot.

\end{itemize}


\textbf{Novelty:}
To the best of our knowledge, this work makes five
individually novel contributions that are further
amplified in combination.
First, the formal modeling of eight selective time
delay and hidden forwarding attack variants that
explicitly characterize how vehicular mobility
amplifies each attack class is absent from prior
literature, which addresses these attacks only in
static SDN or treats them separately.
Second, the dual-mode detection architecture---pairing
lightweight rule-based attack signature detection with
federated LSTM inference and broadcast-radio
witness-based monitoring within a single adaptive
framework---is the first to accommodate the full
spectrum of vehicular node capabilities without
sacrificing detection fidelity at either extreme,
including against passive covert variants that produce
no observable network anomaly.
Third, the integration of \ac{MLDSA} per-packet
integrity, randomized batch verification, and \ac{STARK}
zero-knowledge proofs for both timing compliance and
next-hop legitimacy into a unified post-quantum
cryptographic layer designed specifically for \ac{SDVN}
data-plane and flow-rule integrity is novel in both
scope and target domain.
Fourth, the flat multi-controller architecture with
blockchain-enforced trust scoring, \ac{DKG} among
\ac{RSU} ring members, and smart contract-triggered
failover is the first controller-independent design
in which the controller is excluded from the key
plane entirely.
Collectively, this is the first framework to jointly
detect and mitigate both attack classes under a
zero-trust, controller-independent architecture that
simultaneously preserves driver privacy and satisfies
vehicular latency constraints.




\section{Implications}
The proposed framework demonstrates that decentralization can be done in high-mobility vehicular networks without violating strict latency requirements of SDVN, challenging the common assumption that such systems must rely on centralized and "assumed  trust" controllers. By integrating Federated LSTM with blockchain-based consensus, the study can show the existence of Zero-Trust security models that rely on behavioural verification rather than centralized authority. The privacy-preserving nature of federated learning combined with immutable blockchain auditing suggests higher privacy preserving regulatory and policy development for future, particularly in autonomous vehicle accountability, cyber-physical forensics, and liability assessment. The outcomes of this research may further guide future studies on scalable, privacy-aware security architectures for next-generation vehicular networks.

\section{Paper Organization}

The remainder of this research proposal is organized as follows to provide a comprehensive understanding of the proposed defense framework:
\\

Chapter 2 - Literature Review:
This chapter reviews the existing research related with hidden forwarding  attack, selective time delay attack and their solutions in SDN and VANET domains. The chapter concludes by highlighting key limitations in current centralized and non-privacy-preserving defense mechanisms, and showing  the need for a decentralized and multi-layered security approach.
\\ 

Chapter 3 - Methodology:
This chapter presents the formal system design:
eight attack variant signatures with mobility
amplification analysis, dual-mode detection
algorithms (rule engine and federated LSTM),
a hybrid post-quantum cryptographic integrity
pipeline, witness-based passive attack monitoring,
a two-tier permissioned blockchain with zero-trust
multi-controller architecture, and distributed key
generation.
\\

Chapter 4 - Performance Evaluation and Benchmarking:
This chapter defines twelve performance metrics with
formal equations, presents thirteen ablation
experiments that isolate each component's
contribution, and five system-level benchmarking
experiments against three state-of-the-art baselines
varying attack penetration, mobility speed, network
scale, and attack stealthiness. Full simulation
settings are specified for reproducibility.
\\

Chapter 5 - Timeline and Resources Required:
This chapter outlines the expected project schedule and the required resources.
\\

Chapter 6 - Conclusion:
The final chapter summarizes the importance of the proposed framework and its benefits for the future of SDVNs.


\chapter{Literature Review}
\section{Previous Work}

% ============================================================
% 2.1.1 — UNCHANGED (original)
% ============================================================
\subsection{Hardware and Control Plane Induced Selective Time Delay Attacks}

Earlier work has demonstrated TCAM as a high-speed but very limited resource in SDN switches, usually holding only 1,500 to 8,000 rules \cite{PASCOAL2020107223}. Literature shows that attackers can exploit this small capacity using the Slow TCAM Exhaustion attack, a stealthy variant that periodically refreshes malicious rules at rates as low as 3.2 packets per second to bypass volume-based detection, effectively maintaining a table full state with minimal traffic overhead \cite{PASCOAL2020107223} \cite{pascoal2017slow}. 

While traditional flooding attacks simply crash the network or drop packets, previous studies highlight a specific timing problem within the OpenFlow SDN protocol architecture. Specifically, whenever a packet fails to match an existing forwarding rule, a table-miss occurs, compelling the switch to encapsulate the packet into a PACKET\_IN message to query the controller\cite{PASCOAL2020107223}. The packets failing to match a rule ``experience higher delay'' as they are diverted from the hardware-based ``fast path'' to the software-mediated ``slow path'' (controller-level processing) \cite{9356082} \cite{7929710}. In high-mobility vehicular environments, even small delays of a few milliseconds can disrupt safety-critical coordination. Such delays may lead to serious performance and reliability issues.

Previous studies also show that when the controller processes a PACKET\_IN message through the slow path, it can respond using a PACKET\_OUT message. This message directly instructs the switch, such as an RSU, to send the packet to its destination. By using this method, the controller does not need to add a permanent flow rule to the switch's flow table. Studies show that this approach reduces delays caused by updating the flow table and by packet reordering. As a result, controller-initiated forwarding improves packet delivery performance in latency-sensitive environments. \cite{PASCOAL2020107223} \cite{7929710}.

Previous research has proposed several centralized frameworks to reduce the effects of TCAM exhaustion and the network delays it causes. A primary example is the SelectIve DeFense for TCAM (SIFT) \cite{pascoal2017slow}, which was designed to maintain continuous SDN operation when the TCAM becomes full. SIFT is activated when the switch receives a TABLE\_FULL message. The framework uses a probabilistic method to decide whether a new flow rule should be accepted. If the new flow rule is approved, SIFT randomly removes one existing rule to make space in the TCAM. This process helps ensure that legitimate network users continue to receive service. This study identifies specific hardware conditions in which limited memory causes millisecond-level delays. These delays form the basis of a hardware-based Selective Time Delay attack.

Recent research has examined network desynchronization caused by control plane saturation. In these situations, malicious hosts exploit table-miss packets to consume network resources. The studies propose a framework called Fast Recovery Saturation attack Detection and Mitigation (FSDM) \cite{9356082}. FSDM uses an entropy-based detection method to identify ports involved in the attack. The framework also deploys Mitigation Agents at the network edge to filter malicious traffic. FSDM includes a Force\_Checking module that actively removes attack-related flows from the controller's event queue. This mechanism is reported to reduce ping round-trip time by more than 81\%. The study mainly focuses on restoring system performance after delays caused by saturation. These delays are similar to those observed in Selective Time Delay attack scenarios.

% ============================================================
% 2.1.2 — UNCHANGED (original)
% ============================================================
\subsection{Protocol-Level and Authentication Timing Manipulation in Selective Time Delay Attacks}

Recent studies show that in Software Defined Named Data Networks (SD-NDN) with VANETs, attackers can delay important packets, like those about accidents, VIP messages, or traffic jams \cite{8616999}. To prevent this, researchers propose the Timing Attack Prevention (TAP) protocol, which allows vehicles to calculate signal speed and distance to detect differences between the message's sending time and its arrival time. This study explains how to detect such differences in packet arrival times, which is essential for finding Selective Time Delay attacks \cite{8616999}.

Another case is Infrastructure-based Authentication Timing Attacks, which target communication between vehicles and road infrastructure (V2I). In this attack, an attacker can slow down a vehicle's authentication message sent to a roadside unit (RSU). \cite{Timing_Attack}. To stop this, researchers recommend using a Trusted Platform Module (TPM) and the Direct Anonymous Attestation (DAA) protocol. These methods use secure hardware and strong cryptography to protect the timing and integrity of messages \cite{Timing_Attack}. This study shows that strong cryptographic hardware can maintain timing security in vehicle-to-infrastructure (V2I) systems.

% ============================================================
% 2.1.3 — UNCHANGED (original)
% ============================================================
\subsection{Path Validation and Routing Integrity in Hidden Forwarding Attacks}

Earlier research has identified hidden forwarding attacks, also called Copy attacks, where a compromised switch secretly duplicates packets and sends them to unauthorized destinations while pretending to follow normal routing rules \cite{bu2020unveiling} \cite{mohammadi2016detecting}. To prevent this, previous studies proposed methods that compare the network's forwarding state with the actual movement of packets. One main detection method uses HSA to model how the network should transfer packets and Trajectory Sampling to track the real movement of packets in real time. \cite{mohammadi2016detecting}. This method turns the network into a distributed detection system that can find switches that secretly send duplicate packets to illegal destinations. Although this works well for verifying paths, HSA-based monitoring is centralized, which slows down real-time vehicular security.

% ============================================================
%% SUPERVISOR COMMENT: "INSERT AFTER Section 2, existing literature
%% discussion (after HSA review)" — Subsections 2.1.4–2.1.7 added below.
%% Covers: (a) DA-DIS/ECSD, (b) FADE/LOFT, (c) FTISCON,
%% (d) TFMD-SDVN, (e) LSTM-Fuzzy/CNN-LSTM,
%% (f) FL+blockchain IDS, (g) post-quantum/ZKP schemes.
% ============================================================

% ============================================================
% 2.1.4 — NEW (inserted per supervisor comment)
% ============================================================
\subsection{Detection of Timing and Forwarding Attacks in SDN}
\label{sec:attack_detection_sdn}

Pascoal et al.~\cite{Pascoal2017} demonstrated that Slow-TCAM exhaustion attacks can deny service at rates as low as 3.2\,packets/s and proposed SIFT, a probabilistic rule installation strategy activated upon \texttt{TABLE\_FULL} events. While effective in static SDN, SIFT's random rule dropping can discard safety-critical vehicular flows. Cao et al.~\cite{Cao2023LOFT} extended this line with the LOFT attack, computing the minimum rate to overflow flow tables by fingerprinting TCAM parameters, and proposed LOFTGuard as a collaborative data-to-control-plane defense. Tang et al.~\cite{Tang2023SFTO} introduced SFTO-Guard for real-time slow-rate overflow detection using rule-count prediction and adaptive eviction, though it was validated only on static Mininet topologies.

Huang et al.~\cite{Huang2020} proposed FSDM for control-plane saturation detection using Control Channel Occupation Rate (CCOR) distributions and distributed mitigation agents. FSDM reduces HTTP RTT by 87\% under attack but does not address selective per-flow delays or hidden forwarding.

For forwarding anomaly detection, Zhang et al.~\cite{Zhang2021FADE} proposed FADE, which generates dedicated measurement flow rules to detect packet dropping, duplication, and path deviation using the flow conservation principle. However, FADE relies on flow statistics accuracy from switches and cannot detect selective per-packet delays. Chi et al.~\cite{Chi2020} developed ECSD, a stochastic RNN-based multivariate time-series anomaly detector for compromised SDN switches exhibiting abnormal behaviors including delaying and modifying traffic; however, it requires extensive training data and was not evaluated under vehicular mobility.

Mohammadi et al.~\cite{Mohammadi2016} used Header Space Analysis (HSA) to correlate intended versus actual forwarding paths, detecting packet drops and misroutings. Their approach does not address time delay attacks and assumes static topology. Duffield and Grossglauser~\cite{Duffield2001} proposed trajectory sampling using deterministic hash functions to reconstruct packet paths across network domains; while foundational, the method relies on immutable packet fields and was designed for backbone networks without mobility.

% ============================================================
% 2.1.5 — NEW (inserted per supervisor comment)
% ============================================================
\subsection{Misbehavior and Intrusion Detection in SDVN}
\label{sec:sdvn_detection}

Arsalan and Rehman~\cite{Arsalan2018} proposed TAP, the only timing attack prevention protocol specifically targeting SDN-enabled VANETs. TAP compares expected versus actual response times at the SDN controller to identify attacker vehicles in an NDN-VANET architecture. Despite being mobility-aware in principle, TAP requires a large number of connections for statistical reliability and has not been evaluated under high-density, high-speed scenarios.

Nayak et al.~\cite{Nayak2022} introduced TFMD-SDVN, a trust framework for misbehavior detection at the SDVN edge using rating, recommendation, and similarity metrics. Evaluated with SUMO mobility traces, TFMD-SDVN addresses packet dropping and data falsification but does not target timing manipulation or hidden forwarding. Sedjelmaci et al.~\cite{Sedjelmaci2014} designed ELIDV, a lightweight multi-attack IDS for VANETs using cooperative guard vehicles and reputation-based evaluation; while efficient, it predates SDN integration and does not address control-plane attacks.

Krishnan et al.~\cite{Krishnan2023} proposed FTISCON, preserving flow-table integrity through Ethereum smart contracts that verify flow rules before deployment. Li et al.~\cite{Li2022BlockREV} introduced BlockREV, using blockchain with aggregate signatures for multi-controller rule enforcement verification. Both approaches add blockchain overhead that may be prohibitive for latency-critical vehicular applications and verify only pre-deployment rule correctness—not runtime forwarding fidelity.

Wang et al.~\cite{Wang2020Topology} demonstrated topology poisoning in SDN-enabled vehicular edge networks, showing that vehicular mobility significantly lowers the attack barrier for topology manipulation across four mainstream SDN controllers.

% ============================================================
% 2.1.6 — NEW (inserted per supervisor comment)
% ============================================================
\subsection{AI-Based Anomaly Detection for SDN and Vehicular Networks}
\label{sec:ai_defense}

Novaes et al.~\cite{Novaes2020} combined LSTM with fuzzy logic (LSTM-FUZZY) for DDoS and portscan detection in SDN, achieving autonomous mitigation through OpenFlow rule dropping. While demonstrating LSTM's effectiveness for SDN traffic characterization, the system is centralized and targets only volumetric attacks. Elsayed et al.~\cite{Elsayed2021} proposed a hybrid CNN-LSTM architecture for SDN anomaly detection achieving 96.32\% accuracy on the InSDN dataset; however, generalization to zero-day attacks and vehicular contexts remains unvalidated.

For federated vehicular IDS, Liu et al.~\cite{Liu2021BFLVEC} proposed blockchain-based federated learning for collaborative intrusion detection in vehicular edge computing, offloading training to distributed edge devices while ensuring tamper-proof model aggregation. Abou El Houda et al.~\cite{AbouElHouda2024} advanced this with an edge-based FL+blockchain framework achieving high accuracy against zero-day threats in ITS. Mansouri et al.~\cite{Mansouri2024} implemented MLP, LSTM, GRU, MANN, and CW-RNN models under federated training with blockchain for VANET intrusion detection, achieving 68--94\% precision across attack types on the VeReMi dataset.

\textbf{Critical gap:} Existing FL+LSTM approaches target generic intrusion classes (DDoS, Sybil, position forging). No work applies federated LSTM specifically to detecting selective time delay patterns or hidden forwarding signatures in SDVN, nor do any existing works offer a dual-mode (lightweight rule-based + full AI) detection architecture suitable for resource-heterogeneous vehicular nodes.

\subsection{Cryptographic and Blockchain-Based Defenses for Vehicular Networks}
\label{sec:crypto_defense}

Sumra et al.~\cite{Sumra2011} proposed a TPM-based chain of trust with Direct Anonymous Attestation (DAA) for privacy-preserving VANET authentication. While pioneering, the approach requires hardware TPM modules in every vehicle and does not address SDN-specific threats. Ju\'{a}rez C\'{a}diz and Bordel S\'{a}nchez~\cite{JuarezCadiz2023} developed a dual-layer blockchain architecture with Bayesian reputation assessment, achieving 86\% efficacy against malicious behaviors—insufficient for safety-critical vehicular applications.

For post-quantum vehicular authentication, El-Dalahmeh and Alia~\cite{ElDalahmeh2025} proposed a lattice-based V2R authentication protocol integrating post-quantum cryptography (PQC), zero-knowledge proofs, and fog computing, achieving lower message delay and packet loss than ECC-based methods. Al-Mekhlafi et al.~\cite{AlMekhlafi2024} designed a lattice-based anonymous authentication scheme for 5G vehicular communications using fog computing. A post-quantum ZKP framework for decentralized connected autonomous vehicles was proposed in~\cite{PQZKP2025}, achieving cryptographic commitments in $<$5\,$\mu$s and Binius proof verification in $<$0.17\,s.

Shahrouz and Analouei~\cite{Shahrouz2024} combined ZKP with blockchain for anonymous VANET authentication with conditional privacy preservation. These zero-knowledge approaches enable privacy-preserving verification but introduce computational overhead that remains challenging for real-time V2V messaging.

\textbf{Critical gap:} No existing work integrates ML-DSA-87-based packet integrity, aggregate signatures for batch verification, and ZKP-based privacy preservation into a unified cryptographic layer specifically designed for SDVN flow-rule and data-plane integrity verification. Furthermore, no scheme provides a lightweight fallback mode for resource-constrained scenarios alongside a full cryptographic verification mode.

% ============================================================
\section{Gaps in Literature}
% ============================================================

\subsection{Limitations of Existing Approaches}

\subsubsection{Limitations of the SIFT Framework}

Previous studies \cite{pascoal2017slow} show that the SIFT framework assumes the SDN controller is trustworthy. It is a centralized software module, so if the controller is compromised, the defense fails and accountability is lost.

Furthermore, SIFT uses limited detection logic, removing rules randomly when the flow table is full instead of checking packet timing or content. This method depends on the likelihood that attacker rules will fill the table. SIFT only reacts after a TABLE\_FULL message appears and cannot prevent attacks beforehand.

This focus on availability over speed causes delays; experiments show that attacks make legitimate clients wait longer while rules are removed. Since SIFT relies on a central controller to manage flow rules, it cannot protect the privacy of the data it handles.

\subsubsection{Limitations of the FSDM Framework}

Previous studies \cite{9356082} show that the FSDM framework works under a centralized trust model as a standalone controller module. This creates a single point of failure and cannot audit flow rules in a decentralized way.

Furthermore, FSDM uses fixed thresholds for traffic volume and entropy to detect saturation, which cannot detect slow or subtle attacks like Slow TCAM Exhaustion. Attackers can study these thresholds and change traffic to avoid detection.

Consequently, FSDM reacts only after attack traffic fills the controller's queue, instead of preventing resource exhaustion. The framework collects and processes raw packet-in messages centrally to calculate entropy, exposing sensitive network information and lacking privacy protection.

% ============================================================
% UNCHANGED (original)
% ============================================================
\subsubsection{Limitations of the TAP Protocol}

Previous studies \cite{8616999} show that the TAP protocol uses a centralized trust model and a changeable Controller-Defaulter-List. If the controller is compromised, it could change the list without being noticed.

Additionally, TAP checks message arrival times using fixed physics-based calculations, but attackers can bypass this by changing coordinates. Therefore, TAP only detects delays and cannot find hidden forwarding attacks where packets are duplicated but arrive on time. TAP requires vehicles to send Attacker IDs to the controller, which may risk user privacy.

% ============================================================
% UNCHANGED (original)
% ============================================================
\subsubsection{Limitations of the TPM and DAA Framework}

Previous studies \cite{Timing_Attack} show that TPM and DAA frameworks face problems with scalability and deployment because hardware-based cryptography is expensive and not compatible with older vehicles. These frameworks only check messages and cannot predict future attacks or detect subtle insider threats.

The hardware can be attacked using side-channel techniques, like power analysis or voltage glitches, which can steal cryptographic keys. Therefore, these frameworks cannot detect logical flow attacks, such as hidden forwarding, where malicious nodes send valid packets to unauthorized destinations.

% ============================================================
% UNCHANGED (original)
% ============================================================
\subsubsection{Limitations of the HSA Framework}

Previous studies \cite{mohammadi2016detecting} show that HSA has limitations because it depends on a centralized model, creating a single point of failure and processing delays. If the controller is compromised, the defense does not work.

Furthermore, collecting packet headers and trajectory data in one place can expose sensitive information about user locations and driving patterns. Moreover, HSA cannot clearly tell whether problems are caused by faulty nodes or protocol errors, often marking large areas as suspicious. Malicious switches can escape detection by failing to report or by altering packet headers, so HSA cannot see duplicated or harmful packets.

\subsection{Research Gap}

Analysis of the existing literature reveals five critical gaps.

\textbf{Gap~1 (Mobility-induced vulnerability):} Existing SDN defenses (SIFT, FSDM, FADE, ECSD, LOFT) assume stable topology and well-characterized timing baselines. In SDVN, handoff delays of 50--300\,ms mask attack-induced delays, flow-rule churn from vehicle mobility camouflages slow TCAM exhaustion, and vehicles reside in each RSU zone for only 10--30\,s---too brief for statistical detection. No existing work models or defends against these mobility-amplified attack variants.

\textbf{Gap~2 (No joint detection):} All reviewed detection mechanisms address either timing attacks \emph{or} forwarding anomalies, never both simultaneously. Selective time delay and hidden forwarding attacks can be launched in tandem by a single compromised node; detecting them independently leaves exploitable blind spots.

\textbf{Gap~3 (Centralization and single point of failure):} SDN-based detectors (SPHINX, HSA, LSTM-Fuzzy) rely on a trusted, centralized controller as both the detection engine and the attack target. A compromised controller application can manipulate detection outcomes. Blockchain-based rule verification (FTISCON, BlockREV) partially addresses this but introduces latency unsuitable for real-time vehicular applications.

\textbf{Gap~4 (Privacy violation):} Centralized detection requires mass ingestion of flow statistics and packet metadata, exposing vehicle trajectory and communication patterns. While federated learning approaches (Liu~\cite{Liu2021BFLVEC}, Mansouri~\cite{Mansouri2024}) preserve training data locality, they do not protect individual flow-rule verification queries or audit trail privacy.

\textbf{Gap~5 (No dual-mode defense):} Vehicular nodes range from resource-constrained OBUs to powerful RSU edge servers. No existing framework offers both a lightweight rule-based detection mode for constrained nodes and a full AI/cryptographic verification mode for capable nodes, adaptively switching based on available resources and threat level.

% \hl{...}

\renewcommand{\arraystretch}{1.15}

% \begin{table}[H]
% \centering
% \caption{Comparison of Related Approaches for Selective Time Delay 
% and Hidden Forwarding Attack Detection and Mitigation in SDN/VANET/SDVN}
% \label{tab:lit_comparison}
% \renewcommand{\arraystretch}{1.15}
% \footnotesize
% \begin{tabular}{|p{1.53cm}|p{1.04cm}|p{2.1cm}|p{0.55cm}|p{0.55cm}|p{1.4cm}|p{2.3cm}|p{0.55cm}|p{2.6cm}|}
% % {|p{1.8cm}|p{0.9cm}|p{2.1cm}|p{0.55cm}|p{0.55cm}|p{1.4cm}|p{2.3cm}|p{0.55cm}|p{2.6cm}|}
% \hline
% \textbf{Approach} & \makecell{\textbf{Dom-}\\\textbf{ain}} & \textbf{Attack Addressed} & 
% \textbf{Mob.} & \textbf{Dec.} & \makecell{\textbf{Control-}\\\textbf{ler Trust}}  & 
% \textbf{Detection Method} & \textbf{Priv.} & \textbf{Key Gap} \\
% \hline
% SIFT~\cite{Pascoal2017} 
% & SDN 
% & Slow TCAM exhaustion 
% & N & N 
% & Assumed trusted 
% & Probabilistic rule eviction 
% & N 
% & Reactive only; no timing identification; no audit trail \\
% \hline
% FSDM~\cite{Huang2020} 
% & SDN 
% & Control-plane saturation delay 
% & N & N 
% & Assumed trusted 
% & CCOR entropy analysis 
% & N 
% & Fixed thresholds; no selective per-flow delay detection \\
% \hline
% TAP~\cite{Arsalan2018} 
% & SD-NDN/ VANET 
% & Selective packet timing delay 
% & Y & N 
% & Assumed trusted 
% & Physics-based arrival time check 
% & N 
% & Static thresholds; no hidden forwarding; tamperable list \\
% \hline
% HSA~\cite{Mohammadi2016} 
% & SDN 
% & Packet drops and misrouting 
% & N & N 
% & Assumed trusted 
% & Header Space Analysis 
% & N 
% & No delay detection; centralized raw data collection \\
% \hline
% LOFT~\cite{Cao2023LOFT} 
% & SDN 
% & Low-rate TCAM overflow 
% & N & N 
% & Assumed trusted 
% & TCAM fingerprinting + collaborative detection 
% & N 
% & Static topology only; no per-flow delay detection \\
% \hline
% SFTO-Guard~\cite{Tang2023SFTO} 
% & SDN 
% & Slow-rate flow table overflow 
% & N & N 
% & Assumed trusted 
% & Rule-count prediction + ML 
% & N 
% & Degrades under topology change; no forwarding detection \\
% \hline
% ECSD~\cite{Chi2020} 
% & SDN 
% & Compromised switch behavior 
% & N & N 
% & Assumed trusted 
% & Stochastic RNN time-series 
% & N 
% & No vehicular evaluation; high compute; no forwarding detection \\
% \hline
% FADE~\cite{Zhang2021FADE} 
% & SDN 
% & Packet drop, duplication, path deviation 
% & N & N 
% & Assumed trusted 
% & Flow conservation measurement rules 
% & N 
% & No per-packet delay; relies on switch-reported statistics \\
% \hline
% TFMD-SDVN~\cite{Nayak2022} 
% & SDVN 
% & General misbehavior (drop, falsification) 
% & Y & P 
% & Semi-trusted 
% & Trust rating and recommendation 
% & N 
% & No timing or hidden forwarding specificity \\
% \hline
% FTISCON ~\cite{Krishnan2023} 
% & SDN 
% & Flow-rule tampering 
% & N & Y 
% & Blockchain-verified 
% & Smart contract rule verification 
% & N 
% & Pre-deployment only; no runtime forwarding fidelity \\
% \hline
% LSTM-Fuzzy~\cite{Novaes2020} 
% & SDN 
% & DDoS and volumetric anomalies 
% & N & N 
% & Assumed trusted 
% & LSTM + fuzzy logic mitigation 
% & N 
% & Centralized; volumetric only; no timing or forwarding \\
% \hline
% CNN-LSTM ~\cite{Elsayed2021} 
% & SDN 
% & Multi-class SDN intrusion 
% & N & N 
% & Assumed trusted 
% & Hybrid CNN-LSTM classifier 
% & N 
% & No vehicular context; no timing or forwarding specificity \\
% \hline
% \textbf{Proposed} 
% & \textbf{SDVN} 
% & \textbf{Selective time delay + hidden forwarding (8 variants)} 
% & \textbf{Y} & \textbf{Y} 
% & \textbf{Zero-trust (BC+ZKP)} 
% & \textbf{Dual-mode: rule-based + federated LSTM; EdDSA + agg. sig. + ZKP} 
% & \textbf{Y} 
% & \textbf{---} \\
% \hline
% \end{tabular}
% \vspace{2pt}
% \begin{flushleft}
% \footnotesize
% Mob.\,=\,Mobility-Aware; Dec.\,=\,Decentralized; 
% Priv.\,=\,Privacy-Preserving; P\,=\,Partially; 
% BC\,=\,Blockchain; agg.\,sig.\,=\,Aggregate Signature.
% \end{flushleft}
% \end{table}

\renewcommand{\arraystretch}{1.2}
\captionsetup[table]{justification=raggedright,singlelinecheck=false}

\footnotesize

\begin{longtable}{
|>{\raggedright\arraybackslash}p{1.53cm}
|>{\centering\arraybackslash}p{1.04cm}
|>{\raggedright\arraybackslash}p{2.1cm}
|>{\centering\arraybackslash}p{0.55cm}
|>{\centering\arraybackslash}p{0.55cm}
|>{\raggedright\arraybackslash}p{1.4cm}
|>{\raggedright\arraybackslash}p{2.3cm}
|>{\centering\arraybackslash}p{0.55cm}
|>{\raggedright\arraybackslash}p{2.6cm}|}


\caption{Comparison of Related Approaches for Selective Time Delay 
and Hidden Forwarding Attack Detection and Mitigation in SDN/VANET/SDVN}
\label{tab:lit_comparison} \\

\hline
\textbf{Approach} & \makecell{\textbf{Dom-}\\\textbf{ain}} & \textbf{Attack Addressed} & 
\textbf{Mob.} & \textbf{Dec.} & \makecell{\textbf{Control-}\\\textbf{ler Trust}}  & 
\textbf{Detection Method} & \textbf{Priv.} & \textbf{Key Gap} \\
\hline
\endfirsthead

\hline
\textbf{Approach} & \makecell{\textbf{Dom-}\\\textbf{ain}} & \textbf{Attack Addressed} & 
\textbf{Mob.} & \textbf{Dec.} & \makecell{\textbf{Control-}\\\textbf{ler Trust}}  & 
\textbf{Detection Method} & \textbf{Priv.} & \textbf{Key Gap} \\
\hline
\endhead

SIFT~\cite{Pascoal2017} 
& SDN 
& Slow TCAM exhaustion 
& N & N 
& Assumed trusted 
& Probabilistic rule eviction 
& N 
& Reactive only; no timing identification; no audit trail \\
\hline

FSDM~\cite{Huang2020} 
& SDN 
& Control-plane saturation delay 
& N & N 
& Assumed trusted 
& CCOR entropy analysis 
& N 
& Fixed thresholds; no selective per-flow delay detection \\
\hline

TAP~\cite{Arsalan2018} 
& SD-NDN/ VANET 
& Selective packet timing delay 
& Y & N 
& Assumed trusted 
& Physics-based arrival time check 
& N 
& Static thresholds; no hidden forwarding; tamperable list \\
\hline

HSA~\cite{Mohammadi2016} 
& SDN 
& Packet drops and misrouting 
& N & N 
& Assumed trusted 
& Header Space Analysis 
& N 
& No delay detection; centralized raw data collection \\
\hline

LOFT~\cite{Cao2023LOFT} 
& SDN 
& Low-rate TCAM overflow 
& N & N 
& Assumed trusted 
& TCAM fingerprinting + collaborative detection 
& N 
& Static topology only; no per-flow delay detection \\
\hline

SFTO-Guard~\cite{Tang2023SFTO} 
& SDN 
& Slow-rate flow table overflow 
& N & N 
& Assumed trusted 
& Rule-count prediction + ML 
& N 
& Degrades under topology change; no forwarding detection \\
\hline

ECSD~\cite{Chi2020} 
& SDN 
& Compromised switch behavior 
& N & N 
& Assumed trusted 
& Stochastic RNN time-series 
& N 
& No vehicular evaluation; high compute; no forwarding detection \\
\hline

FADE~\cite{Zhang2021FADE} 
& SDN 
& Packet drop, duplication, path deviation 
& N & N 
& Assumed trusted 
& Flow conservation measurement rules 
& N 
& No per-packet delay; relies on switch-reported statistics \\
\hline

TFMD-SDVN~\cite{Nayak2022} 
& SDVN 
& General misbehavior (drop, falsification) 
& Y & P 
& Semi-trusted 
& Trust rating and recommendation 
& N 
& No timing or hidden forwarding specificity \\
\hline

FTISCON~\cite{Krishnan2023} 
& SDN 
& Flow-rule tampering 
& N & Y 
& Blockchain-verified 
& Smart contract rule verification 
& N 
& Pre-deployment only; no runtime forwarding fidelity \\
\hline

LSTM-Fuzzy~\cite{Novaes2020} 
& SDN 
& DDoS and volumetric anomalies 
& N & N 
& Assumed trusted 
& LSTM + fuzzy logic mitigation 
& N 
& Centralized; volumetric only; no timing or forwarding \\
\hline

CNN-LSTM~\cite{Elsayed2021} 
& SDN 
& Multi-class SDN intrusion 
& N & N 
& Assumed trusted 
& Hybrid CNN-LSTM classifier 
& N 
& No vehicular context; no timing or forwarding specificity \\
\hline

\textbf{Proposed} 
& \textbf{SDVN} 
& \textbf{Selective time delay + hidden forwarding (8 variants)} 
& \textbf{Y} & \textbf{Y} 
% & \textbf{Zero-trust (BC+ZKP)} 
& \makecell[l]{\textbf{Zero}\\
\textbf{-trust}\\
\textbf{(BC+} \\
\textbf{+ZKP}}
& \textbf{Dual-mode: rule-based + federated LSTM; ML-DSA-87 + agg. sig. + ZKP} 
& \textbf{Y} 
& \textbf{---} \\
\hline

\end{longtable}

\vspace{2pt}

\begin{flushleft}
\footnotesize
Mob.\,=\,Mobility-Aware; Dec.\,=\,Decentralized; 
Priv.\,=\,Privacy-Preserving; P\,=\,Partially; 
BC\,=\,Blockchain; agg.\,sig.\,=\,Aggregate Signature.
\end{flushleft}

\normalsize

\chapter{Methodology}

% The methodology for this research adopts a multi-layered approach to design, implement, and validate a decentralized security framework specifically for SDVNs. The process moves from initial behavioral analysis of specific threats to the implementation of a defense system that integrates cryptographically secure blockchain and federated LSTM-based defense mechanism to identify and mitigate stealthy network attacks.

The methodology for this research adopts a multi-layered, zero-trust
approach to design, implement, and validate a decentralized security
framework specifically for SDVNs. The process moves from formal threat
modeling and mobility-specific attack characterization to the
implementation of a defense system that integrates a dual-mode detection
engine, a hybrid cryptographic integrity layer, and a two-tier
permissioned blockchain to jointly detect and mitigate selective time
delay and hidden forwarding attacks while preserving driver privacy
and satisfying vehicular latency constraints.

\section{System Architecture}



Figure \ref{fig:System Architecture} illustrates proposed defense framework, organized into a four-layer hierarchical architecture, designed to provide end-to-end security for SDVNs \cite{Cardona_2020}\cite{s23031600}. This structured approach ensures that intelligence and security tasks are decoupled from the high-speed forwarding requirements of the vehicular data plane \cite{PASCOAL2020107223}.


% Bullet list

\begin{figure}[H]
    \centering
    \includegraphics[width=0.95\textwidth]{System_Architecture.eps} % Replace with your image file
    \caption{System Architecture}
    \label{fig:System Architecture}
\end{figure}


As shown in Figure~\ref{fig:System Architecture}, the proposed defense framework 
is organized into a four-layer hierarchical architecture, designed to provide 
end-to-end security for SDVNs~\cite{Cardona_2020, s23031600}. This structured 
approach ensures that intelligence and security tasks are decoupled from the 
high-speed forwarding requirements of the vehicular data 
plane~\cite{PASCOAL2020107223}.

\begin{itemize}

\item \textbf{Data Plane:} This layer consists of mobile vehicles equipped with 
On-Board Units (OBUs) and stationary Roadside Units (RSUs), which handle local 
Vehicle-to-Vehicle (V2V) and Vehicle-to-Infrastructure (V2I) communications. OBUs 
operate in \textit{Lightweight Mode}, performing HMAC-based packet authentication 
and rule-based signature evaluation locally. When a suspicious event is detected, 
the OBU escalates the flagged traffic to the nearest RSU, which operates in 
\textit{Full Mode}, applying the complete cryptographic integrity pipeline and 
federated LSTM anomaly detection. The RSU is the sole entity responsible for 
writing to the blockchain, submitting model weight updates and detection events 
via PBFT-based endorsement.

\item \textbf{Control Plane:} A registered set of \ac{SDVN} controllers
$\mathcal{C}_{\mathrm{trusted}}$ (read-only, zero-trust) manages the
network by pushing dynamic flow rules to the data plane via Southbound 
APIs~\cite{8187644, 7981537}, such as OpenFlow. Under the zero-trust 
architecture, each controller is granted \textit{read-only} access to the 
blockchain and receives only verified defense policies distributed by smart 
contracts. No controller does not participate in model aggregation, blockchain 
writes, or any trust management operations; its behavior is independently audited 
by RSU-initiated blockchain logging rather than by self-reporting.

\item \textbf{Intelligence Layer:} Local LSTM models are trained exclusively at 
RSU edge nodes on sequential packet data, including timing features, PACKET\_IN 
flood rates, TCAM utilization, and ZKP failure indicators. To preserve privacy, 
only model weights---accompanied by cryptographic hash commitments---are submitted 
to the federated aggregation server rather than raw traffic data. The global 
federated model is produced after Byzantine-robust aggregation and broadcast back 
to all RSUs as a policy output. The SDVN controllers receives the resulting global 
weights in a read-only capacity and has no write access to the aggregation process.

\item \textbf{Security and Consensus Layer:} This is the top layer, where an 
immutable ledger of validated model hashes, endorsed flow rules, ZKP verification 
outcomes, and detection events is maintained using a permissioned 
blockchain~\cite{DBLP:journals/air/QammarKND23, Zhai_2019}. All blockchain writes 
originate exclusively from RSUs via PBFT consensus endorsement; no controller 
has write access. Smart contracts verify submitted model hashes against 
pre-committed values, apply Byzantine-robust Krum filtering during federated 
aggregation, manage per-node trust scores, and automatically issue quarantine 
instructions when a node's trust score falls below the minimum threshold. Verified 
defense policies are distributed from this layer back down to the SDVN controllers, 
eliminating the single point of failure inherent in traditional centralized 
frameworks.

\end{itemize}

\begin{table}[H]
        \scriptsize
        \begin{center}
                \caption{A summary of the notations.}
                %\vspace{-10pt}
                \label{tab_notations}
                \begin{tabular}{| m{2.50cm} | m{8.5cm} |}
                        \hline
                        \textbf{Notation} & \textbf{Description} \\
                        \hline
                        $p, v, r, c_i$ & Packet, vehicle node, Roadside Unit (RSU), and SDVN controller identifiers, respectively. \\
                        \hline
                        $\rho, \bar{v}, t$ & Vehicle density, mean vehicle speed, and discrete time/timestep, respectively. \\
                        \hline
                        $\delta_p, \bar{\delta}_r, \sigma_r$ & Packet forwarding delay, mobility-adjusted baseline delay, and running delay variance, respectively. \\
                        \hline
                        $\lambda_{obs}, \hat{\lambda}_a, \mathbb{E}[\lambda_l|\rho]$ & Observed TCAM rule installation rate, density-normalized anomalous rate, and expected legitimate rate, respectively. \\
                        \hline
                        $\Delta_{max}, T_{fwd}$ & Maximum allowable forwarding delay threshold and safety-critical forwarding deadline, respectively. \\
                        \hline
                        $sk_i, pk_i, k_i$ & Private key, public key, and pre-shared symmetric key of node $i$, respectively. \\
                        \hline
                        $\sigma_i, \boldsymbol{\sigma}$ & Individual ML-DSA-87 signature
                        and accumulated per-hop signature set $\{\sigma_i\}_{i=1}^n$,
                        respectively.
                        \\
                        \hline
                        $\tau_i$ & Lightweight
                        HMAC-SHA3-512 authentication
                        tag (64\,bytes) generated by
                        resource-constrained \ac{OBU}s
                        in place of \ac{MLDSA}
                        signing. \\

                        \hline
                        $\mathbf{r} = (r_1, \ldots, r_n)$ & Fiat-Shamir random challenge vector 
                        for randomized batch verification. \\
                        \hline
                        $\pi_{delay}, \pi_{hop}$ & STARK proofs for timing compliance and next-hop legitimacy at hop i 
                        \\
                         \hline
                         $b_{delay,i}, b_{hop,i}$ &
                        Binary STARK verification outcomes (1=pass, 0=fail) for timing and hop proofs \\
                        \hline
                        $b_\pi$ & Combined binary STARK verification outcome;
                        $b_\pi = \textsc{STARK.Verify}(\pi_{delay},\ldots)
                        \land \textsc{STARK.Verify}(\pi_{hop},\ldots)$;
                        used in trust score update
                        (Eq.~\ref{eq:trust_update}). \\

                        \hline
                                             
                        $\delta_{best}(r,t)$ &
                        Mean per-hop forwarding delay of best-effort
                        (non-high-priority) packets at \ac{RSU} $r$
                        during observation window; used as the selectivity
                        condition in Signature~S1
                        (Eq.~\ref{eq:rule_s1}). \\
                        \hline
                        $\mathcal{E}_{delay}(c_i,t)$ &
                        Delay evidence count for controller $c_i$: number
                        of \ac{RSU}s currently assigned to $c_i$ that have
                        detected a timing anomaly (S1) attributable to
                        $c_i$'s \texttt{FlowMod}s
                        (Eq.~\ref{eq:delay_evidence}). \\
                        \hline
                        $T_{hold}$ &
                        Maximum duration for \ac{OBU} local forwarding
                        suspension pending \ac{RSU} confirmation
                        (Eq.~\ref{eq:local_quarantine}). \\
                        
                       \hline
                        $\mathcal{L}_w, W, H(p)$ & Witness short-term cache, observation time window, and cryptographic packet hash, respectively. \\

                        \hline
                       $\mathcal{P}_{crit}$ &
                        Set of safety-critical (high-priority) packets
                        used in TVR computation
                          (Eq.~\ref{eq:tvr}). \\
                        \hline
                        $\text{TVR}, \text{UCR}$ &
                        Safety-Critical Threshold Violation Rate and
                        Unauthorized Copy Rate; attack-class-specific
                        evaluation metrics for selective time delay
                        and hidden forwarding respectively
                        (Eqs.~\ref{eq:tvr}, \ref{eq:ucr}). \\   
                        
                        \hline
                        $\alpha_w, \beta_w, f$ & Signed duplication alert, signed non-forwarding alert, and maximum number of local Byzantine nodes, respectively. \\
                        \hline
                        $\mathcal{C}, \mathcal{C}_{trusted}$ & Entire registered controller set and active blockchain-trusted controller set, respectively. \\
                        \hline
                        $T_v, T_{c_i}, T_{min}$ & Trust score of node $v$, trust score of controller $c_i$, and minimum operational trust threshold, respectively. \\
                        \hline
                        $\Delta_r, \Delta_p$ & Trust reward increment and trust penalty decrement, respectively. \\
                        \hline
                        $\mathbf{x}_t^{(f)}, \mathbf{h}_t, \mathcal{A}_t^{(k)}$ & Input feature vector, LSTM hidden state, and autoencoder reconstruction anomaly score, respectively. \\
                        \hline
                        $\mathbf{W}_{local}^{(k)}, \mathbf{W}_{global}, \tilde{\mathbf{W}}$ & Local RSU model weight matrix, global aggregated weights, and coordinate-wise weight median, respectively. \\
                        \hline
                        $\gamma, \theta^{(k)}$ & Outlier rejection threshold for Krum filtering and RSU-specific local detection threshold, respectively. \\
                        \hline
                        $T_{ref}(t)$ & Network-wide distributed trusted time reference based on coordinate-wise RSU clock medians. \\
                          \hline
                        $T_{min}^{ctrl}$ &
                        Minimum operational trust
                        threshold for controllers;
                        controllers below this are
                        excluded from
                        $\mathcal{C}_{trusted}$. \\
                        \hline
                        $\Delta_r^{ctrl},
                        \Delta_p^{ctrl}$ &
                        Controller trust reward and
                        penalty increments, swept
                        independently of vehicle/RSU
                        increments $\Delta_r$,
                        $\Delta_p$. \\
                        \hline
                        $T_{sync}$ &
                        Blockchain commit interval for
                        the distributed time reference
                        $T_{ref}(t)$; determines how
                        frequently the consensus clock
                        is anchored on-chain. \\
                        \hline
                        $B$ &
                        Batch size in packets for
                        randomized batch verification
                        at the \ac{RSU}
                        (Eq.~\ref{eq:overhead_batch}).\\

                        \hline
                        $H(p)$ &
                        SHA3-512 packet hash used as the universal packet
                        identity across \ac{UCR} (Eq.~\ref{eq:ucr}),
                        witness cache, and Signature~S6
                        (Eq.~\ref{eq:sig_s6}); replaces \textit{msg\_id}
                        for notational consistency. \\
                        \hline
                        $\textsc{BC.QueryReceipt}(H(p),d',W)$ &
                        Blockchain query returning the reception log entry
                        (Eq.~\ref{eq:packet_receipt_log}) for packet $H(p)$
                        at destination $d'$ within window $W$; enables
                        cross-\ac{RSU} duplication detection for
                        Signatures~S6--S8. \\
                        \hline
                        $v_{atk,copy}$ &
                        Forwarding node identified by per-signature
                        fallback verification (Eq.~\ref{eq:batch_fallback})
                        after \textsc{BatchVerify} failure; the specific
                        node whose signature fails in the multi-hop chain. \\
                          \hline
                          $U_{thresh}$ &
                          TCAM utilisation detection threshold for
                          Signature~S4 (Eq.~\ref{eq:rule_s4}); calibrated
                          from benign validation traces; separates benign
                          utilisation ($\leq 29\%$) from attack saturation
                          ($\approx 100\%$). \\


                \end{tabular}
        \end{center}
        \vspace{-10pt}
\end{table}
\section{Threat Model}
The threat model assumes an internal or external adversary capable of manipulating SDVN components to disrupt communication integrity and timing without triggering immediate volumetric alarms \cite{PASCOAL2020107223}\cite{DBLP:journals/air/QammarKND23}.

The attacker can:
% Bullet list
\begin{itemize}
  \item Delay, duplicate, fabricate, or covertly forward packets.
  \item Exploit SDVN flow rule installation and forwarding behavior.
  \item Participate in both V2I and V2V communication paths\cite{app10093217}.
\end{itemize}

\subsection{Assumptions}

\begin{itemize}
  \item The threat model assumes the attacker will use ``low-rate'' techniques to stay hidden, focusing on duplicating or delaying specific packets rather than crashing the entire network.
  \item The threat model assumes that an adversary couldn't exist in both data plane and control plane at once.
  \item The zero-trust assumption over the SDN controller imposes a strict
architectural constraint: the controller is granted \emph{read-only}
blockchain access and its behavior is audited through independent
RSU-initiated logging rather than self-reporting. All blockchain
write operations originate exclusively from RSU nodes, which are
assumed to be individually corruptible but collectively honest under
the $f < n/3$ Byzantine fault tolerance bound of the PBFT consensus
protocol.

  \item The threat model assumes a flat set of $m$ registered \ac{SDVN}
  controllers $\mathcal{C}=\{c_1,\ldots,c_m\}$. No controller is
  permanently active or inherently trusted, and all controllers are
  granted read-only blockchain access only. The blockchain is the sole
  monitoring authority for controller behavior.


  
\end{itemize}



\section{Attack Scenarios}
The research specifically analyzes two primary stealthy attack scenarios that exploit the centralized logic of the SDVN controller.
% Bulleted List
\begin{itemize}
  \item Selective Time Delay Attack:
  \begin{itemize}
    \item Mechanism: A compromised node (RSU or vehicle) identifies safety-critical packets and introduces intentional, variable lags before forwarding them.\cite{PASCOAL2020107223}
    \item Impact: Packets arrive late rather than being dropped, causing automated systems to react based on ``stale data,'' which is potentially lethal in dynamic environments.
  \end{itemize}

  \subsection{Overview of the four Selective Time Delay Attacks}
\begin{figure}[H]
  \centering
  \includegraphics[width=0.95\textwidth]{selective_time_delay.png} 
  \caption{Selective Time Delay Attack Versions}
  \label{fig:Selective Time Delay Attack Versions}
\end{figure}
  \subsubsection{Attack 1: Selective Time Delay Attack: Control Plane}As illustrated in Figure \ref{fig:Selective Time Delay Attack Versions} (a) an adversary acting as a malicious controller, intentionally transmits manipulated flowMod packets \textcircled{1} to the RSU.When a legitimate packet \textcircled{2} arrives at the RSU, it is processed according to these compromised instructions, causing the RSU to forward the packet with an intentional, variable delay \textcircled{3}.Due to this sensitive operations like collision avoidance warnings which usually reqires a latency of less than 100ms are desyncronized.\cite{Timing_Attack}\cite{KALUPAHANALIYANAGE2018226}.As a result, reliance on ``stale data'' can lead to accidents in dynamic vehicular networks.

  \subsubsection{Attack 2: Selective Time Delay Attack: Data Plane}
 Figure \ref{fig:Selective Time Delay Attack Versions} (b) shows how an adversary within the data plane could manipulate the timing of data flows. In this attack scenario, a malicious RSU or an attacker vehicle acts as the adversary by modifying flowMod data \textcircled{1}. As a result when a legitimate vehicle transmits packets (indicated by \textcircled{2} and \textcircled{4}), the malicious RSU and attacker vehicle intentionally intercept and buffer these packets before forwarding them. Therefore, the forwarded packets (indicated by  \textcircled{3} and \textcircled{5}) arrive at the destination with delays, which violate the strict latency requirements of the vehicular network.

 \subsubsection{Attack 3: Slow TCAM Exhaustion Attack : Contol Plane}
  As shown in the Figure \ref{fig:Selective Time Delay Attack Versions} (c) an adversary in the control plane could delay a packet by exhausting or saturating the TCAM table of the RSU. In this attack scenario, the malicious controller continuously sends manipulated FlowMod packets \textcircled{1} to the RSU, which then will consume the limited TCAM resources. As a result, when a safety-critical packet \textcircled{3} is sent by vehicle A, it fails to find a matching flow rule in the TCAM table. As a result the packet takes the slow path and the packet is forwarded to vehicle B with a significant time delay \textcircled{4}.

  \subsubsection{Attack 4: Slow TCAM Exhaustion Attack: Data Plane}
   Figure \ref{fig:Selective Time Delay Attack Versions} (d) illustrates how an adversary is in the data plane can cause a selective time delay attack by exhausting the RSU's TCAM table \cite{PASCOAL2020107223}\cite{pascoal2017slow}. In this scenario, an attacker vehicle sends a continuous stream of low-rate, unique packets \textcircled{1} to the RSU. These packets are created in a way that they will not match with any existing flow rules resulting the RSU to repeatedly request new instructions from the controller \textcircled{2}, which  saturates the TCAM table eventually. As a result, when vehicle A sends a legitimate, safety-critical packet \textcircled{4}, it fails to find an matching flow rule \textcircled{5} and is forced to ``slow path''\cite{9356082}\cite{7929710}. As a result, the packet \textcircled{6} is forwarded to its destination with an intentional time delay.
   
  
  \item Hidden Forwarding Attack:
  \begin{itemize}
    \item Mechanism: This involves man-in-the-middle techniques where packets are secretly duplicated or forwarded through unauthorized channels or ``hidden tunnels''\cite{app10093217}.
    \item Impact: Original packets follow legitimate paths while duplicates are secretly distributed, enabling eavesdropping.\cite{kaur2024vanetsdn}.
  \end{itemize}

  \subsection{Overview of the four Hidden Forwarding Attacks}
    \begin{figure}[H]
  \centering
  \includegraphics[width=0.95\textwidth]{Hidden Forwarding Attack Versions1.png} % Replace with your image file
  \caption{Hidden Forwarding Attack Versions}
  \label{fig:Hidden Forwarding Attack Versions}
  \end{figure}
  

  \subsubsection{Attack 5: Active Hidden Forward Attack : Contol Plane}
   As shown in Figure \ref{fig:Hidden Forwarding Attack Versions} (e) scenario the adversary is the malicious SDVN controller. The malicious controller starts the attack by sending manipulated FlowMod packets \textcircled{1} to the RSU. Then after the RSU installs these malicious flow rules, it gains the ability to secretly add different rounting paths from original routing paths. Therefore, when Vehicle A attempts to transmit a legitimate packet \textcircled{2} to Vehicle B, the RSU forwards the packet to its original destination \textcircled{3} at the same time duplicating and fabricating the data \textcircled{4} to be sent to an unauthorized vehicle (Vehicle C) within the network. This allows the attacker to steal data or distribute false information while maintaining the appearance of a legitimate communication path.

  \subsubsection{Attack 6: Active Hidden Forward Attack : Data Plane}
   As illustrated in Figure \ref{fig:Hidden Forwarding Attack Versions} (f),in this attack scenario the adverseray is in the data plane it is either an attacker vehicle or an malicious RSU. The attackers starts the attack by modifying flow rules\textcircled{1} and \textcircled{5} to change the forwarding behavior of the RSU and the attacker vehicle. Once these unauthorized rules are active, when Vehicle A sends legitimate packets \textcircled{2} and \textcircled{6} intended for their original destinations they forward the packets to their intended destinations, Vehicle B \textcircled{3} and Vehicle D \textcircled{7}. At the same time, the adversaries create and send fabricated packets \textcircled{4} and \textcircled{8} to an unauthorized Vehicle C. This allows the attackers to secretly forward or duplicate packets while keeping the appearance of a normal communication path to the rest of vehicular network.

  \subsubsection{Attack 7: Passive Hidden Forward Attack: Control Plane}
  As illustrated in Figure \ref{fig:Hidden Forwarding Attack Versions} (g), in this scenario a malicious controller act as the adversary. The attacker initiates the attack by sending manipulated FlowMod packets \textcircled{1} to the RSU. Once the RSU installs these malicious flow rules, it has the ability to duplicate packets. When Vehicle A sends a legitimate packet \textcircled{2} through the RSU, the RSU follows the previously installed instructions to forward the packets. It forwards the original packet \textcircled{3} to the intended destination Vehicle C while secretly duplicating packet \textcircled{4} and sending it to an unauthorized listener vehicle (Vehicle B). Unlike in active forwarding attacks, passive attacks focuses on silent eavesdropping, allowing the controller to steal sensitive data while the original path remains completely unchanged \cite{kaur2024vanetsdn}.

  \subsubsection{Attack 8: Passive Hidden Forward Attack: Data Plane}
    Figure \ref{fig:Hidden Forwarding Attack Versions} (h), shows an passive hidden forwarding attack where data plane nodes either a malicious RSU or an attacker vehicle act as the adversaries. The attack is launched by altering flow rules at the RSU \textcircled{1} and the attacker vehicle \textcircled{5} to override standard flow rules. When Vehicle A forward legitimate packets \textcircled{2} and \textcircled{6} intended for Vehicle C and Vehicle D, the malicious RSU and attacker vehicle forward the duplicated packet through a hidden tunnel. To keep the appearance of a normal connection and avoid detection, the malicious RSU and attacker vehicle forward the original packets \textcircled{3} and \textcircled{7} to their original destinations, Vehicle C and Vehicle D. At the same time, these malicious nodes generate and send additional duplicated packets \textcircled{4} and \textcircled{8} to an unauthorized Vehicle B through a hidden channel. This shows how adversaries within the data plane can secretly duplicate or forward packets while the original communication path remains unaffected.


\end{itemize}
...

\subsection{Mobility-Induced Threat Amplification in SDVN}
\label{sec:mobility_amplification}

The eight attack variants defined in Sections{sec:Overview of the four Selective Time Delay Attacks}
and {sec:Overview of the four Hidden Forwarding Attacks} are not merely transposed from static SDN to a
vehicular setting; vehicular mobility structurally amplifies each
attack class through four mechanisms that must be explicitly modeled
to derive accurate detection thresholds and cryptographic verification
scopes.

\textbf{Mechanism 1: Handoff masking of delay signatures.}
At highway speeds of 80--120\,km/h, a vehicle traverses a typical
RSU coverage zone of 300--500\,m in approximately 9--22\,s, inducing
legitimate handoff latencies of 50--300\,ms during flow-rule
recomputation and reinstallation~\cite{Islam2021SDVN}. An adversary
executing Attack Variants~1 and~2 (Selective Time Delay,
control-plane and data-plane) can inject delays within this range,
rendering attack-induced delay statistically indistinguishable from
legitimate handoff jitter. Let $\delta_h$ denote the handoff latency
distribution and $\delta_a$ the attacker-injected delay; detection
is feasible only when the two distributions are separable, i.e.,
when their Bhattacharyya distance exceeds a detection threshold
$\epsilon$, as formalised in Equation~\eqref{eq:bhattacharyya}:
\begin{equation}
D_B(\delta_h, \delta_a) = -\ln\left(\int \sqrt{p_h(x)\,p_a(x)}\,dx\right) > \epsilon
\label{eq:bhattacharyya}
\end{equation}
As shown in Equation~\eqref{eq:bhattacharyya}, the
Bhattacharyya distance $D_B$ quantifies the overlap between the
handoff latency distribution $p_h(x)$ and the attack-induced delay
distribution $p_a(x)$; a larger $D_B$ indicates greater separability
and thus higher detection feasibility, whereas $D_B \approx 0$
indicates near-complete overlap rendering the two distributions
statistically indistinguishable. Existing static-SDN detectors
implicitly assume $\delta_h = 0$, which forces $D_B \to \infty$ and
makes them blind to the overlap that arises under realistic vehicular
mobility.
for some detection threshold $\epsilon$. Existing static-SDN detectors
implicitly assume $\delta_h = 0$, making them blind to this overlap.

\textbf{Mechanism 2: Topology change exploiting control-plane
staleness.}
The SDVN controller maintains its network view through periodic LLDP
topology discovery messages, typically every 1--5\,s~\cite{Wang2020Topology}.
Between discovery cycles, the controller's topology view is stale by
definition. In static SDN this staleness is negligible; in SDVN,
vehicles enter and leave RSU coverage zones continuously, meaning the
controller's forwarding graph can diverge substantially from the true
network state within a single discovery interval. Attackers exploit
this window in two ways. First, for selective time delay attacks,
a compromised RSU can exploit stale routing decisions that direct
traffic through a suboptimal multi-hop path, adding propagation delay
that is indistinguishable from attack-induced delay since the
controller itself expects elevated latency due to its outdated topology
view. Second, for hidden forwarding attacks, a malicious node can
install covert forwarding rules during a topology update cycle; these
rules are not present in the controller's stale view and therefore
escape flow-rule consistency verification. The effective detection
window narrows with the topology update interval $T_{topo}$; the
probability that a malicious rule survives undetected for $k$
consecutive discovery cycles is formally defined in
Equation~\eqref{eq:evasion_probability}:
\begin{equation}
P_{evade}(k) = \left(1 - P_{detect}\right)^k
\label{eq:evasion_probability}
\end{equation}
As shown in Equation~\eqref{eq:evasion_probability}, the
evasion probability $P_{evade}(k)$ grows exponentially with the
number of consecutive undetected discovery cycles $k$; in
high-mobility SDVN, frequent topology changes reduce $P_{detect}$
per cycle by continuously invalidating the reference forwarding model
used for comparison, causing $P_{evade}(k)$ to approach 1 even for
moderate $k$ and confirming that static-topology detection strategies
are structurally insufficient in vehicular environments.
where $P_{detect}$ is the per-cycle detection probability. In high-mobility
SDVN, frequent topology changes reduce $P_{detect}$ per cycle by
invalidating the reference forwarding model used for comparison,
making $P_{evade}(k)$ unacceptably large even for moderate $k$.

\textbf{Mechanism 3: Flow-rule churn camouflage of TCAM exhaustion.}
In SDVN, legitimate mobility-driven flow installations at busy
intersections reach rates of hundreds of \texttt{FlowMod} messages
per minute as vehicles continuously enter and leave RSU coverage
zones. The slow TCAM exhaustion variants (Attacks~3 and~4) operate
at rates of 3.2--40\,packets/s~\cite{Pascoal2017}, which fall within
the lower tail of the legitimate installation rate distribution.
Even after mobility-conditioning, rate-based detection
alone is insufficient. The mobility-conditioned anomalous
rate estimator is formalised in
Equation~\eqref{eq:density_normalized_rate} to
characterise the residual separation problem.
\begin{equation}
\hat{\lambda}_a(t) = \lambda_{obs}(t) -
\mathbb{E}\!\left[
\lambda_l(t)\,|\,\rho(t),\bar{v}(t)
\right]
\label{eq:density_normalized_rate}
\end{equation}
As given in Equation~\eqref{eq:density_normalized_rate},
even conditioning on both density $\rho(t)$ and speed
$\bar{v}(t)$, the estimated residual $\hat{\lambda}_a(t)$
does not reliably separate attack from benign traffic.
Under realistic \ac{SDVN} mobility, a busy intersection
generates legitimate \texttt{FlowMod} installations at
the same rate as the slow attacker, causing the attack
and benign rate distributions to overlap. For S4, the
rate collapses to zero precisely when the attack is
strongest---when the \ac{TCAM} table is saturated and
new rule installations are refused---inverting the
expected detection signal. Simulation confirms
$\approx 3\%$ detection rate for S3 and $\approx 2\%$
for S4 at the empirical 99th-percentile rate threshold,
validating the theoretical overlap argument. Rate-based
monitoring therefore serves as a formal characterisation
of why mobility camouflage makes slow \ac{TCAM}
exhaustion fundamentally hard to detect by rate alone,
motivating the signal choices in the detection rules
below.







\textbf{Mechanism 4: Short observation windows for hidden forwarding
detection.}
Flow conservation-based hidden forwarding detectors such as FADE and
HSA require sufficient packet samples over a flow's lifetime to detect
duplication or path deviation. In SDVN, a vehicle resides within a
single RSU's observation window for only 9--22\,s at highway speed,
and each topology change effectively resets the observation window
for multi-hop path validation since the reference header space must
be recomputed. Let $N_{obs}$ denote the number of packets observed
within a single RSU zone crossing and $N_{min}$ the minimum sample
size for reliable flow conservation testing at significance level
$\alpha$. The condition that determines whether flow-conservation
detection is statistically reliable is formalised in
Equation~\eqref{eq:observation_window}:
\begin{equation}
N_{obs} = r \cdot t_{zone} < N_{min}(\alpha)
\label{eq:observation_window}
\end{equation}
As shown in Equation~\eqref{eq:observation_window}, the number
of observable packets $N_{obs}$ is the product of the per-flow packet
rate $r$ and the RSU zone residence time $t_{zone}$; when this
product falls below the minimum sample size $N_{min}(\alpha)$ required
for reliable flow conservation testing at significance level $\alpha$,
the detection condition is violated and flow-conservation approaches
become statistically unreliable. This condition is frequently violated
under highway mobility, which fundamentally motivates the use of
cryptographic per-packet ZKP-based behavioral verification as the
primary hidden forwarding mitigation mechanism, since ZKP verification
is per-packet and requires no observation window.
where $r$ is the per-flow packet rate and $t_{zone}$ is the zone
residence time, is frequently violated under highway mobility,
rendering flow-conservation approaches statistically unreliable.
This fundamentally motivates the use of cryptographic per-packet
ZKP-based behavioral verification as the primary hidden forwarding
mitigation mechanism, since ZKP verification is per-packet and
requires no observation window, complemented by federated LSTM
sequence modeling that can detect anomalous patterns over short
windows without requiring flow-level aggregation.



\clearpage
\section{Proposed Solution}
% \begin{enumerate}
%   \item Layer 1: Packet Monitoring: This involves capturing sequential packet data, including timing and routing histories \cite{DBLP:journals/air/QammarKND23}
%   \item Layer 2: Cryptographic Verification: Uses digital signatures and secure hashing to ensure the integrity and authenticity of every packet \cite{Zhai_2019}
%   \item Layer 3: To find odd patterns or anomalies, local long short-term memory (LSTM) models examine time-dependent data. \cite{DBLP:journals/air/QammarKND23}
%   \item Layer 4:Blockchain Ledger:  To remove single points of failure, securely record network events and reach agreement on a tamper-proof distributed ledger. \cite{9519843} \cite{engproc2025092048}
%   \item Layer 5: To improve the global detection model while maintaining privacy, local model updates are combined via a blockchain-protected server.
% \end{enumerate}

\subsection{Role of Cryptography: Pre-Detection or Post-Detection?}
\label{sec:crypto_role}

A critical architectural question is whether the proposed cryptographic
layer independently mitigates both attack classes
or whether detection remains necessary. We analyze each attack class
separately.

For hidden forwarding attacks(figure \ref{fig:Hidden Forwarding Attack Versions}), the cryptographic layer
provides strong mitigation: a forwarding node cannot produce a valid
ZKP demonstrating legitimate next-hop forwarding unless it actually
forwarded the packet to an authorized destination and obtained a valid
receipt signature from that destination. Since the controller and
blockchain verify this proof before updating trust, hidden forwarding
is \emph{cryptographically prevented} for nodes participating honestly
in the signing protocol. However, a Sybil attacker generating colluding
fake next-hop identities could potentially produce false receipt
signatures; detection therefore remains necessary as a second line of
defense.

For selective time delay attacks (figure \ref{fig:Selective Time Delay Attack Versions}), the cryptographic
layer does \emph{not} independently prevent the attack. An attacker can
delay a packet by the desired interval $\delta_a$ and then sign it with
the correct timestamp \emph{after} the delay, producing a valid
signature that reflects the manipulated timestamp. The ZKP demonstrating
forwarding within delay threshold $\Delta_{max}$ would fail in this
case, but only if the receiving next-hop records an independent
timestamp at reception, which is the mechanism that enables detection.
Therefore, for timing attacks, cryptography acts as a \emph{detection
enabler} applied \emph{before} the detection logic: the ML-DSA-87 timestamp binding and STARK proof generation provide the tamper-evident evidence on which
the rule-based and LSTM detectors operate.

Consequently, the architecture applies cryptography \emph{before}
detection for timing attacks (as evidence generation) and \emph{as
primary mitigation} for forwarding attacks (with detection as fallback).
Blockchain serves \emph{both} roles: it immutably logs detection events
(detection support) and enforces trust penalties and node isolation via
smart contracts (mitigation).
\\ \\




% \hl{...}
\subsection{Framework Architecture: Detection and Mitigation Separation}
\label{sec:framework_arch}


The proposed framework cleanly separates detection logic from mitigation
logic across two operational modes as summarized in Table~\ref{tab:dm_separation}.

\begin{table}[H]
\centering
\caption{Detection and Mitigation Logic Separation Across Dual Modes}
\label{tab:dm_separation}
\renewcommand{\arraystretch}{1.1}
\footnotesize
\begin{tabular}{|p{2.8cm}|p{2.8cm}|p{2.8cm}|}
\hline
\textbf{Component} & \textbf{Lightweight Mode} & \textbf{Full Mode} \\
\hline
\textbf{Detection} 
& Rule-based detection of S1 and S2-partial using
mobility-adjusted thresholds and \ac{HMAC} timestamps;
local forwarding suspension on detection.

& Federated LSTM anomaly detection at RSU edge nodes \\
\hline
\textbf{Mitigation} 
& \ac{HMAC}-based lightweight packet authentication;
local forwarding suspension (Eq.~\ref{eq:local_quarantine})
pending \ac{RSU} confirmation.

& ML-DSA-87 + aggregate signatures + ZKP behavioral proofs; blockchain-enforced trust penalties \\
\hline
\textbf{Blockchain} 
& Lightweight event logging only (detection support) 
& Smart contract-enforced trust scoring, node isolation, and model validation (detection + mitigation) \\
\hline
\end{tabular}
\end{table}



The lightweight mode in ~\ref{tab:dm_separation} is designed for resource-constrained OBUs operating
under computational budgets of $<$100\,MHz and $<$256\,MB RAM; the full
mode targets RSU edge nodes with GPU-capable edge computing. Both modes
are simultaneously active in the network; OBUs run lightweight detection
and escalate suspicious events to the nearest RSU, which applies full-mode
analysis. This escalation mechanism ensures that no attack variant evades
detection due to resource constraints at the detecting node.

\subsection{Dual mode Architecture}
\label{sec:dual_mode}
The end-to-end processing pipeline of the dual mode framework, illustrating the two
parallel operational modes and their convergence at the blockchain layer,
is presented in Figure~\ref{fig:dual-mode-framework}.

\begin{figure}[H]
    \centering
    \includegraphics[width=0.69\textwidth]{Figure_1_Final.eps}
    \caption{Dual-Mode Framework Overview}
    \label{fig:dual-mode-framework}
\end{figure}
As shown in Figure~\ref{fig:dual-mode-framework}, the
end-to-end processing pipeline of the MOBIGUARD framework, showing the
two parallel operational modes. The Lightweight Mode (left, blue) runs
on resource-constrained OBUs, where incoming packets undergo HMAC-$\tau$
authentication followed by rule-based signature evaluation, with
suspicious events escalated to the nearest RSU. The Full Mode (right,
green) runs on RSU edge nodes, where packets are signed with ML-DSA-87 and STARK timing proofs, aggregated via randomized batch verification, and verified for
next-hop legitimacy, before feeding into the Federated LSTM anomaly
detector. Both modes converge at the Blockchain Smart Contract (top,
orange), which performs trust score updates, immutable detection logging,
and automated quarantine/mitigation. The SDN Controller operates in a
read-only, zero-trust capacity, receiving only verified policies from
the blockchain rather than serving as a trusted authority.




% As shown in Figure~\ref{fig:dual-mode-framework}, the Lightweight Mode (left) runs on resource-constrained OBUs, where incoming packets undergo HMAC-authentication followed by rule-based signature evaluation, with suspicious events escalated to the nearest RSU. The Full Mode (right) runs on RSU edge nodes where packets are signed with EdDSA and ZKP timing proofs aggregated via batch signatures and verified for next-hop legitimacy, before feeding into the Federated LSTM anomaly detector.Both modes converge at the Blockchain Smart Contract, which performs trust score
% updates, immutable detection logging, and automated
% quarantine/mitigation. The SDN Controller operates in
% a read-only, zero-trust capacity, receiving only verified policies from the blockchain rather than serving as a trusted authority.
% ...
% REQUIRED IN PREAMBLE:
% \usepackage{soul}
% \usepackage{verbatim}   (only if the comment block is kept; better to delete it)

\subsection{Attack Signature Identification}   % ONE heading only
\label{sec:attack_signatures}

The rule-based lightweight detection engine operates on distinguishing
signatures---observable, measurable behavioral patterns that characterize
each attack variant and differentiate it from benign activity. We derive
one primary signature per variant, accounting for mobility-induced
masking identified in Section~\ref{sec:mobility_amplification}.

%------------------------------------------------------------------
% S1
%------------------------------------------------------------------
\textbf{Signature~S1 (Selective Time Delay, Control Plane):}
Safety-critical packets (identified by DSCP priority marking or BSM
message type) originating from the same source vehicle $v$ experience
per-hop forwarding delays at RSU $r$ that consistently exceed the
mobility-adjusted baseline $\bar{\delta}_r(t)$ by more than $k$
standard deviations, while best-effort traffic from the same RSU
remains within baseline. The selective nature (only high-priority
packets are delayed) is the distinguishing feature separating Attack~1
from generic congestion.
The condition formalising Signature~S1 is given in
Equation
\begin{equation}
\text{S1:}\quad
\delta_p(v,r,t) > \bar{\delta}_r(t) + k\sigma_r(t)
\;\land\;
\text{Priority}(p) = \textsc{High}
\;\land\;
\delta_{best}(r,t) \leq \bar{\delta}_r(t) + k\sigma_r(t)
\label{eq:sig_s1}
\end{equation}


As given in Equation~\eqref{eq:sig_s1}, the third
conjunct $\delta_{best}(r,t)\leq\bar{\delta}_r(t)+k\sigma_r(t)$
is the selectivity condition: Attack~1 delays only
high-priority packets while best-effort traffic at
the same \ac{RSU} remains within baseline. Under
genuine congestion, both traffic classes experience
elevated delay simultaneously, so the selectivity
condition evaluates to false and S1 does not fire,
preventing congestion-induced false positives

\textit{Limitation:} The \ac{EWMA} variance estimator
(Eq.~\ref{eq:ewma_variance}) is susceptible to gradual
inflation under a sustained slow attack, since the
attacker's induced delays contribute to $\sigma_r^2(t)$
over time, progressively widening the detection
threshold and reducing sensitivity. In the vehicular
setting this effect is bounded by the zone residence
time of 9--22\,s --- after which the vehicle hands
off to a new \ac{RSU} and the \ac{EWMA} estimate
resets on the fresh \ac{RSU}'s local traffic
distribution. For the slowest simulated attacker
(3.2\,packets/s), the inflation effect within a single
zone crossing is limited and is further mitigated by
the selectivity condition (Eq.~\ref{eq:rule_s1}),
which remains sensitive even under moderate variance
inflation.

%------------------------------------------------------------------
% S2
%------------------------------------------------------------------
\textbf{Signature~S2 (Selective Time Delay, Data Plane):}
The inter-node forwarding delay at an intermediate vehicle node $u$
exceeds the expected single-hop propagation and processing delay
$\delta_{prop} + \delta_{proc}$ by more than $\Delta_{max}$, as
evidenced by a ZKP failure at the next hop (next hop's independently
recorded reception timestamp contradicts the forwarding node's claimed
timestamp).
This condition is formalised in Equation~\eqref{eq:sig_s2}.
\begin{equation}
\text{S2:}\quad t^{recv}_{u+1} - t^{fwd}_{u} > \Delta_{max}
\quad \land \quad \pi_{delay}(u) = \bot
\label{eq:sig_s2}
\end{equation}
As shown in Equation~\eqref{eq:sig_s2}, both the timestamp
discrepancy and the ZKP failure must hold simultaneously, preventing
false positives caused by transient legitimate propagation delays that
do not produce a corresponding proof failure.

%------------------------------------------------------------------
% S3
%------------------------------------------------------------------
\textbf{Signature~S3 (Slow TCAM Exhaustion, Control Plane):}
The \texttt{FlowMod} installation rate $\lambda_{FM}(r,t)$ at RSU $r$
from the controller exceeds the density-normalized expected rate by
threshold $\lambda_{thresh}$, while the installed rules do not
correspond to any active vehicle flow in the network topology.
Detection of this signature does not rely on controller reporting;
the RSU independently cross-checks every received \texttt{FlowMod}
against the blockchain-committed endorsed policy set via
Eq.~\eqref{eq:unauth_flowmod} before installation.
Note that Eq.~\eqref{eq:unauth_flowmod}, referenced here, is
defined in Section~\eqref{sec:blockchain_redesign} as part of the blockchain
architecture; a \texttt{FlowMod} arriving without a corresponding
endorsed commitment is flagged as unauthorised and used as
tamper-proof ground truth for evaluating~S3.
The condition formalising Signature~S3 is given in
Equation~\eqref{eq:sig_s3}.
\begin{equation}
\text{S3:}\quad
f_{unauth}(r_k, t) = 1
\;\land\;
\nexists\, v : \text{flow}(r) \in
\mathcal{F}_{active}(v)
\label{eq:sig_s3}
\end{equation}


As given in Equation~\eqref{eq:sig_s3}, S3 is detected
when a received \texttt{FlowMod} fails the blockchain
endorsement check ($f_{unauth}{=}1$,
Eq.~\ref{eq:unauth_flowmod}) and no corresponding
active vehicle flow exists. The endorsement check is
the primary signal: every legitimate \texttt{FlowMod}
has a valid blockchain commitment by design, so an
unauthorized one is detected at 0\% \ac{FPR}
regardless of its rate. The rate-based estimator
(Eq.~\ref{eq:density_normalized_rate}) is insufficient
as a primary detector because the slow attacker
deliberately operates within the legitimate rate
distribution; its role is characterising the
camouflage problem, not detecting it.
Importantly, detection correctly ceases once the
compromised controller is revoked via
\textsc{SC.Revoke} (Eq.~\ref{eq:sc_revoke}), since
no further unauthorized \texttt{FlowMod}s can be
issued by the revoked controller.

 

%------------------------------------------------------------------
% S4
%------------------------------------------------------------------
\textbf{Signature~S4 (Slow TCAM Exhaustion, Data Plane):}
A single source vehicle $v_{atk}$ generates packets with
unique 5-tuples causing repeated \texttt{PACKET\_IN}
floods to the controller, exhausting the \ac{TCAM} table.
The attack effect is \ac{TCAM} saturation: legitimate
\ac{RSU} \ac{TCAM} utilisation remains below $\approx
30\%$ under normal vehicular mobility, while exhaustion
attacks drive utilisation to $\approx 100\%$, creating
a reliable separation that rate monitoring cannot provide.
The per-source \texttt{PACKET\_IN} rate $\lambda_{PI}$
identifies which vehicle is the highest-volume flooder
and serves as the attacker attribution signal. The S4
detection condition is formalised in
Equation~\eqref{eq:sig_s4}.


\begin{equation}
\text{S4:}\quad
U_{TCAM}(r,t) > U_{thresh}
\;\land\;
v_{atk} = \arg\max_{v'}\,\lambda_{PI}(v',r,t)
\label{eq:sig_s4}
\end{equation}

As given in Equation~\eqref{eq:sig_s4}, $U_{TCAM}>
U_{thresh}$ is the primary detection trigger, exploiting
the empirically confirmed separation between benign
utilisation ($\leq 30\%$) and attack utilisation
($\approx 100\%$). Note that the \texttt{PACKET\_IN}
rate collapses to zero at full saturation since new
rule installations are refused --- using rate as the
primary signal would therefore produce near-zero
detection precisely when the attack is most severe.
$v_{atk}$ is identified as the highest \texttt{PACKET\_IN}
source vehicle at the saturated \ac{RSU} and passed
to \texttt{BTMM} for attribution and quarantine.
The condition $v_{atk}$ has no legitimate active session
requiring such flow diversity is verified against the
blockchain-committed authorised flow policy
(Eq.~\ref{eq:policy_commit}).


\textbf{Signature~S5 (Active Hidden Forwarding, Control Plane):}

The fabrication of packet content is confirmed by an ML-DSA-87 verification
failure on the copy received at $d'$, since the modified content no longer
matches the original sender's signature.
Additionally, the intermediate RSU node $u$ cannot produce a valid
Zero-Knowledge Proof of hop legitimacy $\pi_{hop}(u)$, as the unauthorized
destination $d'$ is absent from the authorized next-hop policy set
$\mathcal{P}(s,d)$.
The conjunction of an explicitly installed unauthorized \texttt{FlowMod}
rule, a failed ML-DSA-87 verification on the fabricated copy, and a failed
hop-legitimacy proof distinguishes Attack~5 from passive forwarding variants and from benign routing updates. A key distinction governs proof behaviour across active
and passive hidden forwarding variants. For active attacks
(S5, S6), the compromised node injects a fabricated copy
through the normal protocol stack at $d'$, where the
copy must carry a signature attributable to the original
sender -- since the attacker does not possess the
original sender's \ac{MLDSA} private key, the copy
fails \ac{MLDSA} verification and consequently fails
aggregate \textsc{BatchVerify}. For passive attacks
(S7, S8), the original already-signed packet is
retransmitted verbatim over the shared broadcast radio
medium, bypassing the protocol stack entirely -- all
cryptographic verifications pass because the
signatures are authentic; only the unauthorized
destination and the witness observations reveal the
attack.
The condition formalising Signature~S5 is given in
Equation~\eqref{eq:sig_s5}.
\begin{equation}
\begin{aligned}
\text{S5:}\quad
  & d' \notin \mathcal{P}(s,d) \\
  & \wedge\; \texttt{FlowMod}(r)
    \text{ installed } (s \to d') \\
  & \wedge\; \textsc{ML-DSA-87.Verify}(
    \sigma_{copy}, pk_s, m_{copy}) = 0 \\
  & \wedge\; \textsc{BatchVerify}(
    \boldsymbol{\sigma},\{pk_i\},\{m_i\},
    \mathbf{r}) = 0 \\
  & \wedge\; b_{hop}(u) = 0
\end{aligned}
\label{eq:sig_s5}
\end{equation}

As given in Equation~\eqref{eq:sig_s5},
$\textsc{ML-DSA-87.Verify}=0$ on the copy at $d'$
(retrieved via Eq.~\ref{eq:packet_receipt_log})
confirms content fabrication; since the attacker lacks
the original sender's private key, the fabricated copy
also fails \textsc{BatchVerify} at the aggregate level.
$b_{hop}(u) = 0$ is the \ac{STARK} next-hop proof
failure. The \textsc{ML-DSA-87.Verify} result at $d'$
is evaluated from the blockchain receipt log
(Eq.~\ref{eq:packet_receipt_log}) rather than locally,
since $d'$ may lie under a different \ac{RSU}.


% ------------------------------------------------------------

\textbf{Signature~S6 (Active Hidden Forwarding, Data Plane):}

The fabrication of packet content is evidenced by an
\ac{MLDSA} verification failure on the copy received
at $d'$, retrieved from the blockchain receipt log
(Eq.~\ref{eq:packet_receipt_log}). Duplication is
confirmed when the same packet hash $H(p)$ is
observed locally at $d$ and queried via
$\textsc{BC.QueryReceipt}(H(p),d',W)=1$ at the
off-path destination $d'$ within window $W$.
The absence of a controller-issued \texttt{FlowMod}
authorizing flows toward $d'$ confirms the attack
originates in the data plane and distinguishes S6
from the control-plane variant S5. The compromised
node $u$ additionally fails to produce a valid
hop-legitimacy \ac{ZKP} $\pi_{hop}(u)$.

The condition formalising Signature~S6 is given in
Equation~\eqref{eq:sig_s6}.
\begin{equation}
\begin{aligned}
\text{S6:}\quad
  & \exists\, d, d' : d \neq d' \\
  & \wedge\; H(p) \in \mathcal{R}(d, W) \\
  & \wedge\; \textsc{BC.QueryReceipt}(H(p), d', W)
    = 1 \\
  & \wedge\; \nexists\,\texttt{FM}(r){:}
    \text{dst}{=}d' \\
  & \wedge\; \textsc{ML-DSA-87.Verify}(
    \sigma_{copy}, pk_s, m_{copy}) = 0 \\
  & \wedge\; b_{hop}(u) = 0
\end{aligned}
\label{eq:sig_s6}
\end{equation}

As given in Equation~\eqref{eq:sig_s6}, the
observation of $H(p)$ at both $d$ (locally) and $d'$
(via blockchain receipt query,
Eq.~\ref{eq:packet_receipt_log}) within window $W$
confirms duplication. The \ac{FlowMod}-absence
conjunct $\nexists\,\texttt{FM}(r){:}\text{dst}{=}d'$
distinguishes S6 (data-plane origin, no controller
\texttt{FlowMod} authorizing $d'$) from S5
(control-plane origin, controller-installed rule).
Without this condition, S6 would also fire under S5
events where the same packet hash arrives at $d'$.
The \ac{MLDSA} copy verification result is retrieved
from the blockchain receipt log at $d'$'s \ac{RSU}
(Eq.~\ref{eq:packet_receipt_log})



% ------------------------------------------------------------

\textbf{Signature~S7 (Passive Hidden Forwarding, Control Plane):}

The distinguishing signature for Passive Hidden
Forwarding, Control Plane (Attack~7) is given in
Equation~\eqref{eq:sig_s7}.
\begin{equation}
\text{S7:}\;
\tfrac{d}{dt}\text{Vol}(d',t) > \epsilon_{vol}
\;\land\; \nexists\,\texttt{FM}(r){:}\text{dst}{=}d'
\;\land\; \textsc{ML-DSA-87.Verify}(
\sigma_c, pk_s, m_c) = 1
\;\land\; b_{hop}(u) = 0
\label{eq:sig_s7}
\end{equation}

As given in Equation~\eqref{eq:sig_s7},
$\textsc{ML-DSA-87.Verify}(\sigma_c, pk_s, m_c){=}1$
confirms the content is unmodified (passive duplication),
distinguishing Attack~7 from active variants where
verification fails; $\pi_{hop}(u){=}\bot$ provides
corroborating cryptographic evidence. The primary
detection mechanism for this attack is witness-based
monitoring via Equations~\eqref{eq:dup_alert_cond}
and~\eqref{eq:bft_penalty}; $b_{hop}(u){=}0$ is a
corroborating indicator only, since a compromised node
controls its own proof generation and can produce a
valid primary-channel proof while covertly duplicating
over the radio medium.



% ------------------------------------------------------------

\textbf{Signature~S8 (Passive Hidden Forwarding, Data Plane):}


The distinguishing signature for Passive Hidden
Forwarding, Data Plane (Attack~8) is given in
Equation~\eqref{eq:sig_s8}.
\begin{equation}
\text{S8:}\;
\textsc{BatchVerify}(\boldsymbol{\sigma},
\{pk_i\}, \{m_i\}, \mathbf{r}) = 1
\;\land\; \textsc{ML-DSA-87.Verify}(
\sigma_c, pk_s, m_c) = 1
\;\land\; b_{hop}(u) = 0
\label{eq:sig_s8}
\end{equation}
As given in Equation~\eqref{eq:sig_s8},
$\textsc{BatchVerify} = 1$ confirms correct
primary-path delivery across all hops;
$\textsc{ML-DSA-87.Verify} = 1$ on the copy
confirms no content modification; and
$b_{hop}(u) = 0$ is the \ac{STARK} proof failure
that is the corroborating indicator of covert
duplication, with primary detection via witness
monitoring (Equations~\eqref{eq:dup_alert_cond}
and~\eqref{eq:bft_penalty}).


\subsection{Distributed Trusted Time Reference}
\label{sec:time_ref}

The network-wide trusted time reference, computed as the
coordinate-wise median across all \ac{RSU} clocks to
eliminate single-authority dependency, is given in
Equation~\eqref{eq:time_consensus}.
\begin{equation}
T_{ref}(t)=\underset{j}{\text{median}}\;\tau_j(t),
\quad j=1,\ldots,n_{RSU}
\label{eq:time_consensus}
\end{equation}
As given in Equation~\eqref{eq:time_consensus}, the
median is robust to up to $f{<}n_{RSU}/2$ faulty or
malicious clocks. No single entity is the authoritative
time source.
The updated per-packet delay measurement anchored to
the consensus reference is given in
Equation~\eqref{eq:delay_updated}.
\begin{equation}
\delta_p(v,r,t)=t_{recv}-t_{send},\quad
t_{send},t_{recv}\text{ anchored to }T_{ref}(t)
\label{eq:delay_updated}
\end{equation}
As given in Equation~\eqref{eq:delay_updated}, both
send and receive timestamps are expressed relative to
$T_{ref}(t)$, preventing a compromised local clock from
inflating or deflating measured delays. $T_{ref}(t)$ is
committed to the blockchain by \ac{RSU} consensus at
regular intervals to provide a tamper-evident audit
trail for all detection decisions.

\subsection{the anomalous Rule-Based Detection Engine}
\label{sec:rule_based}

The lightweight detection engine evaluates signatures
S1, S2-partial, S3, and S4 at each \ac{OBU} node using
mobility-adjusted parameters computed locally without
requiring controller coordination. Signatures S2-full and
S5--S8 are evaluated exclusively at the \ac{RSU} following
escalation (Algorithm~\ref{alg:lrad_rsu}). For timing attacks, the
mobility-adjusted baseline delay at RSU $r$ at time $t$ is estimated
as a function of current vehicle density $\rho(t)$ and mean vehicle
speed $\bar{v}(t)$:
The mobility-adjusted baseline delay is formally defined in
Equation~\eqref{eq:mobility_baseline}.
\begin{equation}
\bar{\delta}_r(t) = \delta_0 + \alpha_\rho \cdot \rho(t) + \alpha_v \cdot \bar{v}(t)^{-1}
\label{eq:mobility_baseline}
\end{equation}
As shown in Equation~\eqref{eq:mobility_baseline}, the baseline
$\bar{\delta}_r(t)$ increases linearly with vehicle density $\rho(t)$
through density sensitivity coefficient $\alpha_\rho$ and decreases with mean speed
$\bar{v}(t)$ through speed sensitivity coefficient $\alpha_v$, with both coefficients
fitted from offline SUMO mobility traces to capture the dual effect of
traffic load and handoff frequency on nominal forwarding latency.
where $\delta_0$ is the static propagation baseline and $\alpha_\rho$,
$\alpha_v$ are density and speed sensitivity coefficients fitted from
offline mobility traces. The standard deviation $\sigma_r(t)$ is
estimated using an exponentially weighted moving average,
where $\beta \in (0,1)$ is the EWMA forgetting factor controlling
the relative weight of recent versus historical delay observations.
The variance estimate is given in Equation~\eqref{eq:ewma_variance}.
\begin{equation}
\sigma_r^2(t) = \beta \cdot \sigma_r^2(t-1)
+ (1-\beta)\cdot\left(\delta_r(t) - \bar{\delta}_r(t)\right)^2
\label{eq:ewma_variance}
\end{equation}
As shown in Equation~\eqref{eq:ewma_variance}, a larger $\beta$
retains more historical information, producing a smoother but less
responsive estimate, while a smaller $\beta$ prioritises recent
observations and adapts more quickly to rapid changes in vehicular
density.
A packet $p$ is flagged by the S1 rule as formalised in
Equation

\begin{equation}
f_{S1}(p,r,t) =
\begin{cases}
1 & \delta_p > \bar{\delta}_r(t){+}k\sigma_r(t)
    \;\land\; \text{Priority}(p){=}\textsc{High}
    \;\land\; \delta_{best}(r,t){\leq}
    \bar{\delta}_r(t){+}k\sigma_r(t) \\
0 & \text{otherwise}
\end{cases}
\label{eq:rule_s1}
\end{equation}

As given in Equation~\eqref{eq:rule_s1}, all three
conditions must hold simultaneously: the delay
threshold, the high-priority class, and the selectivity
condition confirming best-effort traffic is within
baseline. This ensures that genuine congestion, which
elevates delay across all traffic classes, does not
satisfy the selectivity conjunct and does not trigger
the flag.

For \ac{TCAM} exhaustion detection, the rate-based
estimator (Eq.~\ref{eq:density_normalized_rate}) was
found insufficient as a detection signal under
realistic vehicular mobility (see
Eq.~\ref{eq:density_normalized_rate} discussion in
Section~\ref{sec:mobility_amplification}). The S3
detection flag uses the blockchain endorsement check
as the primary signal
(Eq.~\ref{eq:rule_s3}) and the S4 detection flag
uses \ac{TCAM} utilisation saturation
(Eq.~\ref{eq:rule_s4}).

The S3 detection flag is given in
Equation~\eqref{eq:rule_s3}.
\begin{equation}
f_{S3}(r,t) =
\begin{cases}
1 & f_{unauth}(r_k,t) = 1
    \;\land\; \nexists\,v :
    \text{flow}(r) \in \mathcal{F}_{active}(v) \\
0 & \text{otherwise}
\end{cases}
\label{eq:rule_s3}
\end{equation}

As given in Equation~\eqref{eq:rule_s3}, S3 fires on
the blockchain endorsement failure alone ---
$f_{unauth}{=}1$ is a binary deterministic signal with
0\% \ac{FPR} by construction, so no rate threshold is
needed. The absence-of-active-flow condition cross-checks
that the \texttt{FlowMod} has no legitimate vehicle
session behind it, confirming exhaustion intent rather
than a legitimate but unendorsed update.


\begin{equation}
f_{S4}(r,t) =
\begin{cases}
1 & U_{TCAM}(r,t) > U_{thresh} \\
0 & \text{otherwise}
\end{cases}
\;\text{with}\;
v_{atk} = \arg\max_{v'}\,\lambda_{PI}(v',r,t)
\label{eq:rule_s4}
\end{equation}

As given in Equation~\eqref{eq:rule_s4}, $U_{TCAM}>
U_{thresh}$ is the sole detection condition, calibrated
at $U_{thresh}$ [tbd: $\{0.5, 0.6, 0.7, 0.8\}$]
selected for best \ac{MCC} at $\leq 1\%$ \ac{FPR} on
the benign validation split. $v_{atk}$ is computed
concurrently as the highest \texttt{PACKET\_IN} source
vehicle and passed to \texttt{BTMM}
(Eq.~\ref{eq:batch_fallback}) for quarantine; it does
not gate the detection decision.



The composite \ac{OBU}-side detection decision for
node $v$ at time $t$ is given in
Equation~\eqref{eq:composite_light}.
\begin{equation}
D_{OBU}(v,t) = f_{S1}(v,t) \lor f_{S2p}(v,t)
\label{eq:composite_light}
\end{equation}

As given in Equation~\eqref{eq:composite_light},
$D_{OBU}$ evaluates to 1 if either of the two
\ac{OBU}-computable signature rules fires,
immediately triggering local flow-rule quarantine
(Eq.~\ref{eq:local_quarantine}) and escalation to
the nearest \ac{RSU}. Signatures S3 and S4 require
\ac{RSU}-side infrastructure observables (FlowMod
rate, TCAM utilisation, PACKET\_IN rate per source)
that are inaccessible to resource-constrained \ac{OBU}s
and are evaluated exclusively in \texttt{LRAD-RSU}.
Signatures S2-full and S5--S8 require \ac{ZKP}
inputs unavailable on \ac{OBU}s.


The lightweight detection algorithm deployed on
resource-constrained \ac{OBU}s evaluates only signatures
computable from local data and \ac{HMAC}-authenticated
timestamps without any \ac{ZKP} dependency.
The \ac{OBU}-side detection procedure is given in
Algorithm~\ref{alg:lrad_obu}.


\begin{algorithm}[H]
\caption{\ac{OBU} Lightweight Rule-Based Detection
(\texttt{LRAD-OBU})}
\begin{algorithmic}[1]
\small
\Procedure{\texttt{LRAD-OBU}}{$p,v,r,t,\rho,\bar{v},\tau$}
  \State $\bar{\delta} \gets \delta_0
         +\alpha_\rho\cdot\rho
         +\alpha_v\cdot\bar{v}^{-1}$
         \hfill\Comment{Mobility-adjusted baseline
         (Eq.~\ref{eq:mobility_baseline})}
  \State $\sigma_r \gets \texttt{EWMA}(\delta_r,\beta)$
         \hfill\Comment{Running variance
         (Eq.~\ref{eq:ewma_variance})}
  \State $flag_{S1} \gets
         [\delta_p{>}\bar{\delta}{+}k\sigma_r]
         \land[\text{Priority}(p){=}\textsc{High}]
         \land[\delta_{best}(r,t){\leq}
         \bar{\delta}{+}k\sigma_r]$
         \hfill\Comment{S1: selective delay,
         congestion-filtered
         (Eq.~\ref{eq:rule_s1})}
  \State $ts_{recv} \gets \texttt{HMAC.Verify}(\tau,p)$
         \hfill\Comment{Recover authenticated timestamp}
  \State $flag_{S2p} \gets
         (t_{now}-ts_{recv})>\Delta_{max}$
         \hfill\Comment{S2-partial: threshold only,
         no ZKP}
  \State $D_{OBU} \gets flag_{S1}\lor flag_{S2p}$
         \hfill\Comment{S3, S4 evaluated at RSU only}
  \If{$D_{OBU}$}
    \State \texttt{HOLD\_FORWARD}$(v,r)$
           \hfill\Comment{Local quarantine: suspend
           forwarding pending RSU confirmation
           (Eq.~\ref{eq:local_quarantine})}
    \State \texttt{ESCALATE}$(p,v,r,
           \{flag_{S1},flag_{S2p}\})$
           \hfill\Comment{S2-full, S3--S8 at RSU only}
  \EndIf
  \State \textbf{return} $D_{OBU}$
\EndProcedure
\end{algorithmic}
\label{alg:lrad_obu}
\end{algorithm}

The local flow-rule quarantine action executed by the
\ac{OBU} upon detection is formalised in
Equation~\eqref{eq:local_quarantine}.
\begin{equation}
\texttt{HOLD\_FORWARD}(v,r) \Rightarrow
\text{fwd\_state}(v) = \textsc{Hold}
\quad\text{until}\quad
\texttt{RSU.Confirm}(v,r)\;\lor\;
t > t_{detect} + T_{hold}
\label{eq:local_quarantine}
\end{equation}
As given in Equation~\eqref{eq:local_quarantine},
upon $D_{OBU}{=}1$ the \ac{OBU} suspends forwarding
of flows matching the flagged signature for a maximum
hold window $T_{hold}$, awaiting either \ac{RSU}
confirmation or timeout. This prevents an \ac{OBU}
from continuing to forward potentially attack-affected
traffic before the \ac{RSU} completes full-mode
analysis.


The \ac{RSU}-side full detection algorithm receives escalated
events from \ac{OBU}s and evaluates all remaining
\ac{ZKP}-dependent signatures using the full cryptographic
pipeline. The \ac{RSU}-side full detection procedure is given
in Algorithm~\ref{alg:lrad_rsu}.

\begin{algorithm}[H]
\caption{\ac{RSU} Full Detection (\texttt{LRAD-RSU})}
\begin{algorithmic}[1]
\small
\Procedure{\texttt{LRAD-RSU}}{$p, v, r, \pi_{delay},
  \pi_{hop}, \boldsymbol{\sigma},
  \{pk_i\}, \{m_i\}, D_{LSTM}^{(k)}$}
  % --- S2-full ---
  \State $flag_{S2f} \gets
         [\textsc{STARK.Verify}(\pi_{delay},
         C_{delay}, H_{\text{SHA3-512}}) = 0]$
         \hfill\Comment{S2-full: timing proof failure}
  % --- Batch and hop verification ---
  \State $\mathbf{r} \gets
         H_{\text{SHA3-512}}(\boldsymbol{\sigma}
         \| m_1 \| \cdots \| m_n)$
         \hfill\Comment{Eq.~\ref{eq:batch_challenge}}
  \State $b_{batch} \gets
         \textsc{BatchVerify}(\boldsymbol{\sigma},
         \{pk_i\}, \{m_i\}, \mathbf{r})$
         \hfill\Comment{Eq.~\ref{eq:batch_verify}}
  \State $b_{hop} \gets
         \textsc{STARK.Verify}(\pi_{hop},
         C_{hop}, H_{\text{SHA3-512}})$
         \hfill\Comment{Eq.~\ref{eq:stark_hop_verify}}
 \State $flag_{S3} \gets
         [f_{unauth}(r,t){=}1]
         \land[\nexists\,v{:}\text{flow}(r)
         {\in}\mathcal{F}_{active}(v)]$
         \hfill\Comment{S3: Eq.~\ref{eq:rule_s3};
         blockchain endorsement primary}
  % --- S4: Data-plane TCAM (utilization-based) ---
  \State $flag_{S4} \gets
         [U_{TCAM}(r,t){>}U_{thresh}]$
         \hfill\Comment{S4: Eq.~\ref{eq:rule_s4};
         TCAM saturation primary}
  \State $v_{atk} \gets
         \arg\max_{v'}\,\lambda_{PI}(v',r,t)$
         \hfill\Comment{Attribution only; does not
         gate $flag_{S4}$}
  % --- S5--S8 (existing) ---
   % --- Attacker identification after batch failure ---
  \If{$\neg b_{batch}$}
    \State $v_{atk,copy} \gets
           \arg\min_{i}\{i :
           \textsc{ML-DSA-87.Verify}
           (\sigma_i,pk_i,m_i)=0\}$
           \hfill\Comment{Eq.~\ref{eq:batch_fallback}:
           isolate failing node}
  \EndIf
  % --- S5: Active HF Control Plane ---
  \State $\text{CopyVerify}_{d'} \gets
         \textsc{BC.QueryReceipt}(H(p),d',W).\text{verify}$
         \hfill\Comment{Eq.~\ref{eq:packet_receipt_log}}
  \State $flag_{S5} \gets
         [\neg b_{batch}]
         \land[f_{unauth}(r,t){=}1]
         \land[\text{CopyVerify}_{d'}{=}0]
         \land[b_{hop}{=}0]$
         \hfill\Comment{S5: Eq.~\ref{eq:sig_s5};
         all four conjuncts}
  % --- S6: Active HF Data Plane ---
  \State $\text{RecvAt}_{d'} \gets
         \textsc{BC.QueryReceipt}(H(p),d',W)$
         \hfill\Comment{Cross-RSU via blockchain}
  \State $\text{NoFM}_{d'} \gets
         \nexists\,\texttt{FM}(r){:}\text{dst}{=}d'$
         \hfill\Comment{FlowMod-absence check}
  \State $flag_{S6} \gets
         [\text{RecvAt}_{d'}{=}1]
         \land[\text{NoFM}_{d'}]
         \land[\text{CopyVerify}_{d'}{=}0]
         \land[b_{hop}{=}0]$
         \hfill\Comment{S6: Eq.~\ref{eq:sig_s6};
         gated on FlowMod-absence to exclude S5}
  % --- S7: Passive HF Control Plane ---
  \State $\text{CopyVerify}_{d',S7} \gets
         \textsc{BC.QueryReceipt}(H(p),d',W).\text{verify}$
  \State $flag_{S7} \gets
         [\tfrac{d}{dt}\text{Vol}(d',t){>}\epsilon_{vol}]
         \land[\nexists\,\texttt{FM}(r){:}\text{dst}{=}d']
         \land[\text{CopyVerify}_{d',S7}{=}1]
         \land[b_{hop}{=}0]$
         \hfill\Comment{S7: Eq.~\ref{eq:sig_s7};
         all four conjuncts; primary via witness}
  % --- S8: Passive HF Data Plane ---
  \State $\text{CopyVerify}_{d',S8} \gets
         \textsc{BC.QueryReceipt}(H(p),d',W).\text{verify}$
  \State $flag_{S8} \gets
         [b_{batch}]
         \land[\text{CopyVerify}_{d',S8}{=}1]
         \land[b_{hop}{=}0]$
         \hfill\Comment{S8: Eq.~\ref{eq:sig_s8};
         all three conjuncts; primary via witness}

  % --- LSTM path ---
  \State $flag_{LSTM} \gets D_{LSTM}^{(k)}$
         \hfill\Comment{Eq.~\ref{eq:lstm_detection};
         covers residual anomalies}
  % --- Composite RSU decision ---
  \State $D_{RSU} \gets flag_{S2f}\lor flag_{S3}
         \lor flag_{S4}\lor flag_{S5}\lor flag_{S6}
         \lor flag_{S7}\lor flag_{S8}\lor flag_{LSTM}$
  \If{$D_{RSU}$}
    % --- Plane-based attribution and routing ---
    \If{$flag_{S3}\lor flag_{S5}\lor flag_{S7}$}
      \State $c_{atk} \gets c^*(r,t)$
             \hfill\Comment{Eq.~\ref{eq:rsu_ctrl_assign}}
      \State \texttt{BC.Write}$(r,c_{atk},D_{RSU},ts)$
      \State \texttt{BTMM}$(c_{atk},
             \textsc{ctrl-plane})$
             \hfill\Comment{Controller trust update path}
      \If{$flag_{S5}$}
        \Comment{S5: dual attribution ---
        also penalize forwarding node}
        \State \texttt{BTMM}$(v_{atk,copy},
               \textsc{data-plane})$
               \hfill\Comment{Node identified by
               Eq.~\ref{eq:batch_fallback}}
      \EndIf
      \If{$flag_{S7}\lor flag_{S8}$}
        \Comment{Witness path invoked independently
        by SC.PenalizeTrust via Eq.~\ref{eq:bft_penalty}}
        \State \textit{// Witness threshold checked
               on-chain; no direct BTMM call needed}
      \EndIf
    \Else
      \hfill\Comment{Data-plane or LSTM: attacker is
      vehicle/RSU}
      \If{$flag_{S4}$}
        \State $v_{report} \gets v_{atk}$
      \Else
        \State $v_{report} \gets v$
      \EndIf
      \State \texttt{BC.Write}$(r,v_{report},
             D_{RSU},ts)$
      \State \texttt{BTMM}$(v_{report},
             \textsc{data-plane})$
    \EndIf
    \State \texttt{RSU.Confirm}$(v,r)$
           \hfill\Comment{Release OBU hold if
           warranted}
  \EndIf
  \State \textbf{return} $D_{RSU}$
\EndProcedure
\end{algorithmic}
\label{alg:lrad_rsu}
\end{algorithm}


As outlined in Algorithms~\ref{alg:lrad_obu}
and~\ref{alg:lrad_rsu}, \texttt{LRAD} splits
detection across a lightweight \ac{OBU} stage and a
full \ac{RSU} stage. The \ac{OBU} stage evaluates
only S1 and S2-partial --- the two signatures
computable from local data and \ac{HMAC}-authenticated
timestamps --- using mobility-adjusted baselines, and
immediately suspends forwarding (Eq.~\ref{eq:local_quarantine})
upon detection. The \ac{RSU} stage evaluates all
remaining signatures: S2-full via \ac{STARK} timing
proof, S3 and S4 via \ac{RSU}-observable TCAM metrics
and mobility-conditioned rate estimation, and S5--S8
via the aggregate signature and hop-legitimacy
pipeline. The \ac{RSU} stage further incorporates the
federated \ac{LSTM} detection output $D_{LSTM}^{(k)}$
as a complementary detection signal. Critically,
mitigation is routed by attack plane: control-plane
attacks (S3, S5, S7) attribute to the assigned
controller and invoke the controller trust update and
failover path; data-plane attacks (S2, S4, S6, S8)
attribute to the source vehicle and invoke node
quarantine.


\subsection{Witness-Based Forwarding Verification}
\label{sec:witness}

Zero-knowledge proofs demonstrate compliance on the primary
forwarding channel but cannot prove absence of covert
secondary transmissions. A compromised node can produce a
valid $\pi_{hop}$ for the authorised next-hop while
simultaneously transmitting a duplicate packet to an
unauthorised listener over the shared radio medium, or
suppressing forwarding entirely beyond the safety deadline.
A complementary witness-based mechanism exploits the
broadcast nature of vehicular radio to provide external
verification independent of the forwarding node.

\subsubsection{Witness Observation Model}
Every vehicle $v_w$ within radio range operates in
promiscuous monitoring mode, logging overheard transmissions
in short-term cache $\mathcal{L}_w$ with entries
$(H(p), \text{dst}, ts)$ where $H(p)$ is the packet hash,
$\text{dst}$ is the observed destination identity, and
$ts$ the observation timestamp.


\subsubsection{Duplication Alert}
The condition under which witness $v_w$ raises a duplication
alert is given in Equation~\eqref{eq:dup_alert_cond}.
\begin{equation}
\text{DA}_w(p){=}1\iff
\exists\,dst{\neq}dst':
H(p)\in\mathcal{L}_w(dst,W)
\land H(p)\in\mathcal{L}_w(dst',W)
\label{eq:dup_alert_cond}
\end{equation}
As given in Equation~\eqref{eq:dup_alert_cond},
$\mathcal{L}_w(dst,W)$ denotes hashes observed directed
to $dst$ within window $W$. Under single-path routing,
the same hash at two distinct destinations is impossible
for legitimate traffic.

\subsubsection{Delay-Equivalent Non-Delivery Detection}
A forwarding delay exceeding safety deadline $T_{fwd}$
renders a packet functionally equivalent to non-delivery,
constituting a severe case of existing Selective Time Delay
Attack Variants~1--4 without introducing a new attack type.
The delay-equivalent non-delivery condition is given in
Equation~\eqref{eq:nfwd_detect}.
\begin{equation}
\text{NFA}_w(p){=}1\iff
H(p)\in\mathcal{L}_w^{recv}(v_i,W)
\land H(p)\notin\mathcal{L}_w^{fwd}(v_i,T_{fwd})
\label{eq:nfwd_detect}
\end{equation}
As given in Equation~\eqref{eq:nfwd_detect},
$\mathcal{L}_w^{recv}(v_i,W)$ is the set of hashes
witnessed as received by $v_i$, and
$\mathcal{L}_w^{fwd}(v_i,T_{fwd})$ the set witnessed as
forwarded within $T_{fwd}$.


\subsubsection{Signed Alert Submission}
The signed duplication alert is given in
Equation~\eqref{eq:da_sign}.
\begin{equation}
\alpha_w = \textsc{ML-DSA-87.Sign}\!\left(sk_w,\;
H_{\text{SHA3-512}}\!\left(
H(p) \,\|\, dst \,\|\, dst' \,\|\, ts_w
\right)\right)
\label{eq:da_sign}
\end{equation}
As given in Equation~\eqref{eq:da_sign}, $sk_w$
is the witness vehicle's \ac{MLDSA} private key.
\ac{MLDSA} is required here rather than \ac{HMAC}
because $\alpha_w$ is written directly to the
blockchain and must be attributable to a specific
witness vehicle by any third-party verifier without
a pre-shared key. The signed non-forwarding alert
is given in Equation~\eqref{eq:nfa_sign}.
\begin{equation}
\beta_w = \textsc{ML-DSA-87.Sign}\!\left(sk_w,\;
H_{\text{SHA3-512}}\!\left(
H(p) \,\|\, v_i \,\|\, ts_w \,\|\, T_{fwd}
\right)\right)
\label{eq:nfa_sign}
\end{equation}
As given in Equation~\eqref{eq:nfa_sign}, $\beta_w$
binds the suspected node $v_i$ and the safety
deadline $T_{fwd}$ under the witness's \ac{MLDSA}
key, ensuring the non-forwarding alert is
non-repudiable and independently verifiable
on-chain.


\subsubsection{BFT Trust Penalty Threshold}
A single witness alert must never trigger a trust
penalty directly. The smart contract penalty trigger
is given in Equation~\eqref{eq:bft_penalty}.
\begin{equation}
\textsc{SC.PenalizeTrust}(v_i) \Leftarrow
\Bigl|\bigl\{
w : \exists\,(\alpha_w \text{ or } \beta_w)
\text{ s.t. }
\textsc{ML-DSA-87.Verify}(\cdot,pk_w,\cdot) = 1
\;\land\; \text{event}(v_i,p,w)\text{ confirmed}
\bigr\}\Bigr|
\geq 2f{+}1
\label{eq:bft_penalty}
\end{equation}

As given in Equation~\eqref{eq:bft_penalty}, $f$ is
the maximum number of colluding malicious witnesses
in the locality. The cardinality is over
\textbf{distinct witness vehicles} $w$ that have
each submitted at least one valid alert
(signed under their own $pk_w$) for the same
$(v_i, p)$ event --- not over the total number of
alert messages. A single witness submitting multiple
alerts for the same event is counted as one, preventing
any one vehicle from crossing the $2f{+}1$ threshold
unilaterally. The smart contract deduplicates by
$(w, v_i, H(p))$ before counting, and independently
verifies each alert under $pk_w$ to reject forged
submissions.




\subsection{Hybrid Cryptographic Integrity Layer}
\label{sec:crypto_layer}



Figure~\ref{fig:figure2} illustrates the complete
cryptographic packet pipeline with multi-hop V2V routing, showing
the sequence of ML-DSA-87 signing, randomized batch aggregation, STARK proof generation, RSU
endorsement, and blockchain commitment operations that together
constitute the hybrid cryptographic integrity layer.
\begin{figure}[H]
    \centering
    \includegraphics[width=1.1\textwidth]{Figure_2_Final.eps}
    \caption{Cryptographic Packet Pipeline Sequence Diagram with Multi-hop V2V Routing.}
    \label{fig:figure2}
\end{figure}
As shown in Figure~\ref{fig:figure2}, the pipeline proceeds
through three structural stages. In the first stage, each
intermediate Vehicle Relay at hop $i$ aggregates the accumulated
per-hop ML-DSA-87 signature and generates per-hop STARK proofs before
forwarding to the RSU, directly reflecting the multi-hop attack
model and its impact on packet delivery ratio and end-to-end
latency. In the second stage, the RSU endorsement block shows
neighboring RSUs submitting PBFT endorsements to the blockchain,
which commits the flow rule only upon collecting $f{+}1$ matching
endorsements, implementing the zero-trust two-tier chain structure.
In the third stage, red crossed boxes mark the cryptographic failure
points corresponding to each attack class: Attacks~1--2 (selective
time delay) are detected via timing STARK proof failure; Attacks~5--6
(active hidden forwarding) are detected via aggregate signature
verification failure and a failed copy verification at the
unauthorised destination; and Attacks~7--8 (passive relay copying) pass both aggregate
and copy verification and are detected primarily through
witness-based monitoring
(Equations~\eqref{eq:dup_alert_cond}
and~\eqref{eq:bft_penalty}), with hop-legitimacy
\ac{STARK} proof failure ($b_{hop} = 0$) serving as a
corroborating cryptographic indicator. \ac{STARK} proof
failure alone is insufficient as a primary detector
because a compromised node can generate a valid
primary-channel proof while covertly duplicating over
the shared radio medium.


\subsubsection{Per-Packet ML-DSA-87 Signing}

Each forwarding node $v_i$ with \ac{MLDSA} private key
$sk_i$ signs every forwarded packet to provide NIST
Security Level~5 post-quantum integrity, authenticity,
and non-repudiation. The per-packet signing operation
is given in Equation~\eqref{eq:mldsa_sign}.
\begin{equation}
\sigma_i = \textsc{ML-DSA-87.Sign}\!\left(sk_i,\;
H_{\text{SHA3-512}}\!\left(
\text{msg\_id} \,\|\, ts_i \,\|\,
\eta_i \,\|\, \text{nh}_i \,\|\, z_i
\right)\right)
\label{eq:mldsa_sign}
\end{equation}
As given in Equation~\eqref{eq:mldsa_sign}, $sk_i$ is
the \ac{MLDSA} private key of node $v_i$, $ts_i$ is
the forwarding timestamp, $\eta_i$ is a fresh nonce
preventing replay, $\text{nh}_i$ is the chosen next-hop
identity, and $z_i$ is the current \ac{RSU} zone
identifier preventing cross-zone replay. The message is
hashed under \ac{SHAKE} before signing, binding identity,
time, routing decision, and mobility context in a single
quantum-resistant signature of 4{,}595 bytes per
packet. \ac{MLDSA} is chosen over pairing-based schemes
because it achieves NIST Security Level~5 via lattice
hardness assumptions (Module Learning With Errors),
providing post-quantum security consistent with the
\ac{ZKP} layer described below.




\subsubsection{Aggregate Signature via Randomized
Batch Verification}

Since \ac{MLDSA} does not support algebraic aggregation
via bilinear pairings, the aggregate signature is
constructed using randomized batch verification, which
combines all per-hop signatures into a single joint
lattice check at the \ac{RSU}. As a packet traverses
$n$ hops $v_1, v_2, \ldots, v_n$, the accumulated
signature set $\boldsymbol{\sigma} =
\{\sigma_i\}_{i=1}^n$ is forwarded alongside the packet.
At the \ac{RSU}, a non-interactive random challenge
vector is derived from the full transcript. The challenge
derivation is given in Equation~\eqref{eq:batch_challenge}.
\begin{equation}
\mathbf{r} = (r_1, \ldots, r_n) =
H_{\text{SHA3-512}}\!\left(
\sigma_1 \,\|\, \cdots \,\|\, \sigma_n \,\|\,
m_1 \,\|\, \cdots \,\|\, m_n
\right)
\label{eq:batch_challenge}
\end{equation}
As given in Equation~\eqref{eq:batch_challenge},
$\mathbf{r}$ is a Fiat-Shamir random challenge derived
by hashing the concatenation of all $n$ signatures and
their corresponding signed messages under \ac{SHAKE},
making the challenge unpredictable to any adversary who
has not yet committed to all signatures. Each \ac{MLDSA}
signature $\sigma_i = (\tilde{c}_i, \mathbf{z}_i,
\mathbf{h}_i)$ consists of a commitment hash
$\tilde{c}_i$, response vector $\mathbf{z}_i \in
\mathbb{Z}_q^\ell$, and hint vector $\mathbf{h}_i$.
The batch verification condition is given in
Equation~\eqref{eq:batch_verify}.
\begin{equation}
\textsc{BatchVerify}\!\left(
\{\sigma_i\}, \{pk_i\}, \{m_i\}, \mathbf{r}
\right) =
\left[
\sum_{i=1}^{n} r_i \mathbf{A}_i \mathbf{z}_i
\stackrel{?}{=}
\sum_{i=1}^{n} r_i
\!\left(
\mathbf{t}_i \mathbf{c}_i +
\mathbf{w}'_i
\right) \pmod{q}
\right]
\label{eq:batch_verify}
\end{equation}
As given in Equation~\eqref{eq:batch_verify},
$\mathbf{A}_i = \textsc{ExpandA}(\rho_i)$ is the public
matrix expanded from the seed $\rho_i$ in public key
$pk_i = (\rho_i, \mathbf{t}_i)$, $\mathbf{c}_i =
H_{\text{SHA3-512}}(\mu_i \,\|\, \tilde{c}_i)$ is the
challenge polynomial recovered from the commitment hash
and message hash $\mu_i = H_{\text{SHA3-512}}
(\text{tr}_i \,\|\, m_i)$, and $\mathbf{w}'_i =
\textsc{UseHint}(\mathbf{h}_i,\,
\mathbf{A}_i \mathbf{z}_i - \mathbf{t}_i \mathbf{c}_i)$
is the hint-corrected approximation. A single check of
the combined lattice equation is sufficient to accept
or reject all $n$ signatures simultaneously; any
individual forgery perturbs the random linear combination
and causes the equality to fail with overwhelming
probability $1 - n/q$ over the random challenge $\mathbf{r}$.

Randomized batch verification via
Equation~\eqref{eq:batch_verify} reduces the cost
of verifying $n$ independent \ac{MLDSA} signatures
from $n$ separate lattice checks to approximately
one combined matrix-vector operation over the random
linear combination. The dominant cost is $n$ matrix
multiplications $\mathbf{A}_i \mathbf{z}_i$ and the
corresponding inner product accumulation, all of which
are independent and can be parallelised across
multi-core \ac{RSU} hardware. The asymptotic cost
remains $O(n)$ matrix operations; however, the
constant factor is substantially reduced relative
to $n$ independent verifications because the combined
check shares common reduction steps. This parallelism
ensures that batch verification of the $n$ per-hop
signatures completes within the 50\,ms processing
budget at \ac{RSU} edge nodes even for realistic
multi-hop path lengths in dense vehicular environments.

When $\textsc{BatchVerify}=0$, the batch rejection
identifies that at least one hop signature is invalid
but does not identify which specific node is responsible.
A per-signature linear fallback is executed immediately
after batch failure to isolate the attacker. The
attacker identification step is formalised in
Equation~\eqref{eq:batch_fallback}.
\begin{equation}
v_{atk} = \arg\min_{i \in \{1,\ldots,n\}}
\Bigl\{i :
\textsc{ML-DSA-87.Verify}(\sigma_i, pk_i, m_i) = 0
\Bigr\}
\label{eq:batch_fallback}
\end{equation}
As given in Equation~\eqref{eq:batch_fallback}, the
linear scan over individual signatures identifies the
first failing node $v_{atk}$, whose identity is then
passed to \texttt{BTMM} for trust update and
quarantine. The $O(n)$ per-signature fallback is
incurred only on attack detection events, not on
normal traffic, preserving the efficiency of
\textsc{BatchVerify} for the common benign case.




\subsubsection{STARK Proof for Timing Compliance}

Each forwarding node $v_i$ generates a post-quantum
zero-knowledge proof demonstrating that its forwarding
delay was within threshold $\Delta_{max}$ without
revealing raw timestamps. Unlike Groth16, which requires
a trusted setup and relies on elliptic curve pairings,
the \ac{STARK} scheme uses \ac{SHAKE} as its sole
cryptographic primitive, providing a transparent setup
and NIST Security Level~5 post-quantum resistance
consistent with the \ac{MLDSA} signing layer.

The timing compliance proof is given in
Equation~\eqref{eq:stark_delay}.
\begin{equation}
\pi_{delay,i} = \textsc{STARK.Prove}\!\left(
\{t^{recv}_i,\, t^{fwd}_i,\, \rho_i\};\;
C_{delay};\;
H_{\text{SHA3-512}}
\right)
\label{eq:stark_delay}
\end{equation}
As given in Equation~\eqref{eq:stark_delay}, the
private witnesses are the reception timestamp
$t^{recv}_i$, the forwarding timestamp $t^{fwd}_i$,
and blinding randomness $\rho_i$; the arithmetic
circuit $C_{delay}$ encodes the constraint
$t^{fwd}_i - t^{recv}_i \leq \Delta_{max}
\;\land\; H_{\text{SHA3-512}}(\rho_i) = c_i$,
where $c_i$ is the publicly committed blinding
commitment; and \ac{SHAKE} is the hash function
used throughout the FRI polynomial commitment layer.
The proof $\pi_{delay,i}$ is publicly verifiable
without exposing $t^{recv}_i$ or $t^{fwd}_i$ to
any verifier. Verification is given in
Equation~\eqref{eq:stark_delay_verify}.
\begin{equation}
b_{delay,i} = \textsc{STARK.Verify}\!\left(
\pi_{delay,i},\; C_{delay},\;
H_{\text{SHA3-512}}
\right) \in \{0, 1\}
\label{eq:stark_delay_verify}
\end{equation}
As given in Equation~\eqref{eq:stark_delay_verify},
$b_{delay,i} = 1$ confirms the delay constraint was
satisfied without revealing any witness values;
$b_{delay,i} = 0$ flags the node as having delayed
the packet beyond $\Delta_{max}$, contributing to
Signatures~S1 and~S2 detection.


\subsubsection{STARK Proof for Next-Hop Legitimacy}

Each forwarding node $v_i$ additionally proves that
its chosen next hop $v_{i+1}$ is authorised per the
current routing policy and produced a valid receipt
signature. The next-hop legitimacy proof is given in
Equation~\eqref{eq:stark_hop}.
\begin{equation}
\pi_{hop,i} = \textsc{STARK.Prove}\!\left(
\{\sigma^{rcpt}_{i+1},\, \mathcal{P}_i\};\;
C_{hop};\;
H_{\text{SHA3-512}}
\right)
\label{eq:stark_hop}
\end{equation}
As given in Equation~\eqref{eq:stark_hop}, the private
witnesses are the receipt signature
$\sigma^{rcpt}_{i+1}$ from the next hop and the
local copy of the authorised next-hop policy set
$\mathcal{P}_i$ committed to the blockchain at flow
setup. The arithmetic circuit $C_{hop}$ encodes the
constraint $\textsc{ML-DSA-87.Verify}(
\sigma^{rcpt}_{i+1}, pk_{i+1}) = 1 \;\land\;
v_{i+1} \in \mathcal{P}_i$, certifying both
delivery and routing compliance without revealing
the receipt signature to external verifiers.
Verification is given in
Equation~\eqref{eq:stark_hop_verify}.
\begin{equation}
b_{hop,i} = \textsc{STARK.Verify}\!\left(
\pi_{hop,i},\; C_{hop},\;
H_{\text{SHA3-512}}
\right) \in \{0, 1\}
\label{eq:stark_hop_verify}
\end{equation}
As given in Equation~\eqref{eq:stark_hop_verify},
$b_{hop,i} = 0$ is a corroborating cryptographic
indicator for passive hidden forwarding
(Signatures~S7, S8); the primary detection mechanism
for passive variants is witness-based monitoring
via Equations~\eqref{eq:dup_alert_cond}
and~\eqref{eq:bft_penalty}. For active hidden
forwarding (Signatures~S5, S6), $b_{hop,i} = 0$
is a contributing indicator alongside aggregate
verification failure. The combined proof
$\pi_i = (\pi_{delay,i}, \pi_{hop,i})$ is forwarded
alongside $\boldsymbol{\sigma}$ to the \ac{RSU}
and blockchain for verification.


\subsubsection{Lightweight Authentication (HMAC-SHA3-512)}

For resource-constrained \ac{OBU}s operating in
lightweight mode, \ac{MLDSA} is replaced with an
\ac{HMAC} under \ac{SHAKE}, using a pre-shared
symmetric key $k_i$ established during \ac{RSU}
association, since \ac{OBU}s do not participate in
aggregate signature construction or \ac{STARK} proof
generation. The lightweight authentication tag is
given in Equation~\eqref{eq:hmac_light}.
\begin{equation}
\tau_i = \textsc{HMAC-SHA3-512}\!\left(k_i,\;
\text{msg\_id} \,\|\, ts_i \,\|\, \eta_i
\right)
\label{eq:hmac_light}
\end{equation}
As given in Equation~\eqref{eq:hmac_light}, $\tau_i$
is a 64-byte tag providing integrity and authenticity
under the 512-bit \ac{SHAKE} primitive, consistent
with NIST Security Level~5. It does not provide
non-repudiation since the verifier shares the key
$k_i$; lightweight nodes therefore contribute only
to rule-based detection and depend on \ac{RSU}-level
full verification for mitigation.

The per-packet cryptographic overhead in each mode
is evaluated to characterise the cost of achieving
NIST Security Level~5 post-quantum assurance. The
full-mode overhead per packet at hop $i$ is given
in Equation~\eqref{eq:overhead_full}.
\begin{align}
O_{full}^{(i)} &=
\underbrace{|\sigma_i|}_{4{,}595\,\text{B}}
+ \underbrace{|\pi_{delay,i}|}
  _{\leq\,100\,\text{KB}}
+ \underbrace{|\pi_{hop,i}|}
  _{\leq\,100\,\text{KB}}
\;\approx\; 200\,\text{KB per hop}
\label{eq:overhead_full}
\end{align}
As given in Equation~\eqref{eq:overhead_full},
$|\sigma_i| = 4{,}595$\,bytes is the \ac{MLDSA}
signature size per hop, and $|\pi_{delay,i}|$,
$|\pi_{hop,i}| \leq 100$\,KB each are the
\ac{STARK} proof sizes for timing compliance and
next-hop legitimacy respectively under conservative
FRI blowup parameters. The total per-batch
verification overhead at the \ac{RSU}, amortised
over a batch of $B$ packets each traversing $n$
hops, is given in Equation~\eqref{eq:overhead_batch}.
\begin{equation}
O_{batch} = B \cdot \left(
n \cdot |\sigma_i|
+ |\pi_{delay}| + |\pi_{hop}|
\right)
+ \textsc{BatchVerify cost}
\label{eq:overhead_batch}
\end{equation}
As given in Equation~\eqref{eq:overhead_batch},
the overhead is amortised across the batch at the
\ac{RSU} edge node, which has sufficient compute
to perform the combined lattice check of
Equation~\eqref{eq:batch_verify} and parallel
\ac{STARK} verification within the 50\,ms
processing budget; this overhead is incurred
exclusively in full mode and is never imposed on
resource-constrained \ac{OBU}s. The lightweight
overhead is given in
Equation~\eqref{eq:overhead_light}.
\begin{equation}
O_{light} = |\tau_i| = 64\,\text{bytes
(HMAC-SHA3-512)}
\label{eq:overhead_light}
\end{equation}
As given in Equation~\eqref{eq:overhead_light},
the lightweight \ac{OBU} mode incurs only a
64-byte \ac{HMAC} tag, making it viable for
resource-constrained nodes on narrow vehicular
radio channels. The overhead contrast between full
and lightweight modes is deliberate: NIST Level~5
post-quantum security is provided where compute
is available (RSU edge), while minimal overhead
is maintained where it is not (OBU).




The full cryptographic integrity pipeline, executing
on \ac{RSU} edge nodes, is given in
Algorithm~\ref{alg:fcip}.

\begin{algorithm}[htb]
\caption{Full Cryptographic Integrity Pipeline
(\texttt{FCIP})}
\begin{algorithmic}[1]
\small
\Procedure{\texttt{FCIP}}{$p, v_i, sk_i, z_i,
  \Delta_{max}, \mathcal{P}_i$}
  \State $\eta_i \gets \texttt{RAND}()$
         \hfill\Comment{Fresh nonce}
  \State $\sigma_i \gets
         \texttt{ML-DSA-87.Sign}(sk_i,
         H_{\text{SHA3-512}}(
         \text{msg\_id}\|ts_i\|\eta_i\|
         \text{nh}_i\|z_i))$
         \hfill\Comment{Eq.~\eqref{eq:mldsa_sign}}
  \State $\boldsymbol{\sigma} \gets
         \boldsymbol{\sigma}^{prev}
         \cup \{\sigma_i\}$
         \hfill\Comment{Accumulate hop signatures}
  \State $\pi_{delay} \gets
         \texttt{STARK.Prove}(C_{delay},
         t^{recv}_i, t^{fwd}_i, \rho_i;
         H_{\text{SHA3-512}})$
         \hfill\Comment{Eq.~\eqref{eq:stark_delay}}
  \State $\sigma^{rcpt} \gets
         \texttt{AWAIT\_RECEIPT}(v_{i+1})$
  \State $\pi_{hop} \gets
         \texttt{STARK.Prove}(C_{hop},
         \sigma^{rcpt}, \mathcal{P}_i;
         H_{\text{SHA3-512}})$
         \hfill\Comment{Eq.~\eqref{eq:stark_hop}}
  \State \texttt{FORWARD}$(p,
         \boldsymbol{\sigma}, \pi_{delay},
         \pi_{hop})$
  \State \textbf{return}
         $(\boldsymbol{\sigma}, \pi_{delay},
         \pi_{hop})$
\EndProcedure
\end{algorithmic}
\label{alg:fcip}
\end{algorithm}

As shown in Algorithm~\ref{alg:fcip}, the procedure
signs each packet under \ac{MLDSA} at NIST Level~5
(Eq.~\ref{eq:mldsa_sign}), accumulates the signature
into the growing set $\boldsymbol{\sigma}$, and
generates two \ac{STARK} proofs — one for timing
compliance (Eq.~\ref{eq:stark_delay}) and one for
next-hop legitimacy (Eq.~\ref{eq:stark_hop}). At the
terminal \ac{RSU}, \textsc{BatchVerify}
(Eq.~\ref{eq:batch_verify}) is invoked over the full
accumulated set $\boldsymbol{\sigma}$ with the
Fiat-Shamir challenge (Eq.~\ref{eq:batch_challenge}).


\subsection{Federated LSTM Anomaly Detection}
\label{sec:fed_lstm}

The federated LSTM architecture with blockchain-verified Byzantine-robust
aggregation is depicted in Figure~\ref{fig:federated-lstm}.

\begin{figure}[H]
    \centering
    \includegraphics[width=1.1\textwidth]{Figure_3_Final.eps}
    \caption{Federated LSTM architecture with blockchain-verified aggregation.}
    \label{fig:federated-lstm}
\end{figure}

As shown in Figure~\ref{fig:federated-lstm}, each RSU trains a
local LSTM on input vector $x_t$ , uploads model weights with hash
commitment $H(W_{local})$ for integrity verification. The server
applies Krum filtering and broadcasts $W_{global}$ back to all RSUs.
The Controller (read-only) node (leftmost) enforces the zero-trust
constraint: it receives $W_{global}$ for policy output only and has
no write access to the aggregation layer.

\subsubsection{Local LSTM Model}
Each RSU $r_k$ trains a local LSTM model on its observed packet
sequences. The input feature vector at timestep $t$ for flow $f$
is formally defined in Equation~\eqref{eq:lstm_input}:
\begin{equation}
\mathbf{x}_t^{(f)} = \left[\delta_t,\, \lambda_{PI,t},\,
U_{TCAM,t},\, \mathbb{1}[\pi_{delay}=\bot],\,
\mathbb{1}[\pi_{hop}=\bot],\, \rho_t,\, \bar{v}_t \right]^T
\label{eq:lstm_input}
\end{equation}
As shown in Equation~\eqref{eq:lstm_input}, the feature vector
combines three network-level observables---the per-hop forwarding
delay $\delta_t$, the PACKET\_IN flood rate $\lambda_{PI,t}$, and the
TCAM utilisation $U_{TCAM,t}$---with two novel mobility-cryptographic
binary indicators encoding ZKP timing compliance failure
$\mathbb{1}[\pi_{delay}=\bot]$ and hop-legitimacy failure
$\mathbb{1}[\pi_{hop}=\bot]$, augmented by the current vehicle density
$\rho_t$ and mean speed $\bar{v}_t$ to provide the LSTM with explicit
mobility context at each timestep.
where the ZKP failure indicators $\mathbb{1}[\pi=\bot]$ are novel
mobility-cryptographic features that encode cryptographic evidence
of attacks directly into the LSTM input. The hidden state evolves
according to Equation\eqref{eq:lstm_hidden}:
\begin{equation}
\mathbf{h}_t = \textsc{LSTM}\!\left(\mathbf{x}_t, \mathbf{h}_{t-1};\,
\mathbf{W}_{local}^{(k)}\right)
\label{eq:lstm_hidden}
\end{equation}
As shown in Equation~\eqref{eq:lstm_hidden}, the hidden state
$\mathbf{h}_t$ is updated at each timestep as a function of the
current input vector $\mathbf{x}_t$ and the previous hidden state
$\mathbf{h}_{t-1}$ under the locally trained weight matrix
$\mathbf{W}_{local}^{(k)}$, enabling the model to capture temporal
dependencies in the packet sequence that are characteristic of
sustained attack behaviour rather than transient benign anomalies.
The anomaly score for timestep $t$ is the reconstruction error of an
autoencoder-mode LSTM,
as defined in Equation~\eqref{eq:anomaly_score}:
\begin{equation}
\mathcal{A}_t^{(k)} =
\left\|\mathbf{x}_t - \hat{\mathbf{x}}_t\right\|_2^2
\label{eq:anomaly_score}
\end{equation}
As shown in Equation ~\eqref{eq:anomaly_score}, the anomaly
score $\mathcal{A}_t^{(k)}$ is the squared $\ell_2$ reconstruction
error between the observed feature vector $\mathbf{x}_t$ and its
autoencoder reconstruction $\hat{\mathbf{x}}_t$; a high score
indicates that the current packet sequence deviates significantly
from the distribution of benign traffic seen during local training.
A detection event is triggered
according to the threshold rule defined in
Equation~\eqref{eq:lstm_detection}:
\begin{equation}
D_{LSTM}^{(k)}(t) =
\begin{cases}
1 & \mathcal{A}_t^{(k)} > \theta^{(k)} \\
0 & \text{otherwise}
\end{cases}
\label{eq:lstm_detection}
\end{equation}
As shown in Equation~\eqref{eq:lstm_detection}, the binary
detection flag $D_{LSTM}^{(k)}(t)$ is raised whenever the anomaly
score exceeds the RSU-specific threshold $\theta^{(k)}$, ensuring
that each RSU applies a locally calibrated decision boundary rather
than a global threshold that would fail to account for
RSU-specific traffic distributions.
where the RSU-specific threshold $\theta^{(k)}$ is computed from the
local training distribution,
as given in Equation~\eqref{eq:lstm_threshold}:
\begin{equation}
\theta^{(k)} = \mu_{\mathcal{A}}^{(k)} + z_\alpha \cdot \sigma_{\mathcal{A}}^{(k)}
\label{eq:lstm_threshold}
\end{equation}
As shown in Equation\eqref{eq:lstm_threshold}, the threshold
$\theta^{(k)}$ is set at $z_\alpha$ standard deviations above the
mean anomaly score $\mu_{\mathcal{A}}^{(k)}$ observed on benign
training traffic at RSU $r_k$, where $z_\alpha$ is chosen to achieve
the target false positive rate $\alpha$, allowing each RSU to
independently control its sensitivity without centralised threshold
tuning.
and $z_\alpha$ is the $z$-score corresponding to false positive rate
$\alpha$.



\subsubsection{Federated Aggregation with Byzantine Robustness}
Since vehicular nodes may be compromised, standard FedAvg is replaced
with a Byzantine-robust aggregation strategy. The global model weights
are updated
as formalised in Equation \eqref{eq:fed_robust}:
\begin{equation}
\mathbf{W}_{global}^{(t+1)} =
\frac{\sum_{k=1}^{K} n_k \cdot \mathbf{W}_{local}^{(k)}(t)
\cdot \mathbb{1}\!\left[d\!\left(\mathbf{W}_{local}^{(k)},
\tilde{\mathbf{W}}\right) < \gamma\right]}
{\sum_{k=1}^{K} n_k \cdot \mathbb{1}\!\left[d\!\left(\mathbf{W}_{local}^{(k)},
\tilde{\mathbf{W}}\right) < \gamma\right]}
\label{eq:fed_robust}
\end{equation}
As shown in Equation\eqref{eq:fed_robust}{, the global weight
update is a weighted average restricted exclusively to RSUs whose
submitted weights lie within Euclidean distance $\gamma$ of the
coordinate-wise median $\tilde{\mathbf{W}}$; RSUs whose updates
deviate beyond $\gamma$---indicating potential gradient poisoning by
a compromised node---are excluded from the aggregation by the
indicator function $\mathbb{1}[\cdot]$, so that no single
Byzantine RSU can shift the global model arbitrarily.}
where $\tilde{\mathbf{W}}$ is the coordinate-wise median of submitted
weights, $d(\cdot,\cdot)$ is Euclidean distance, and $\gamma$ is an
outlier rejection threshold. This Krum-inspired filtering excludes
model updates from RSUs whose weights deviate significantly from the
consensus---rejecting gradient poisoning by compromised RSU nodes.
Model hash submissions are written to the RSU-chain exclusively by
RSUs; the smart contract verifies the hash of submitted weights
before inclusion in aggregation, and the controller has no
participation in this validation process;
the blockchain model hash verification condition is defined in
Equation~\eqref{eq:bc_model_verify}:
\begin{equation}
\mathbb{1}_{BC}^{(k)} = \textsc{SC.VerifyModelHash}\!\left(
H(\mathbf{W}_{local}^{(k)}),\; \textsc{CommittedHash}^{(k)}\right)
\label{eq:bc_model_verify}
\end{equation}
As shown in Equation~\eqref{eq:bc_model_verify}\, the binary
indicator $\mathbb{1}_{BC}^{(k)}$ returns 1 only when the hash of
the submitted local weights $H(\mathbf{W}_{local}^{(k)})$ matches
the hash pre-committed to the blockchain by RSU $r_k$, ensuring that
no RSU can submit a weight update that differs from its committed
model and that the controller, which has no blockchain write access,
cannot interfere with the validation process.



The Byzantine-robust federated aggregation
procedure incorporating trust score gating is given in
Algorithm~\ref{alg:brfa_v2}.

\begin{algorithm}[htb]
\caption{Updated Byzantine-Robust Federated Aggregation
(\texttt{BRFA-v2})}
\begin{algorithmic}[1]
\small
\Procedure{\texttt{BRFA-v2}}{$\{\mathbf{W}^{(k)}\}_{k=1}^{K},
\{n_k\},\{T_{r_k}\},\gamma,T_{min}$}
  \State \textbf{Step 1: Trust gate}
  \State $\mathcal{K}_e\gets
         \{k:T_{r_k}\geq T_{min}\}$
         \hfill\Comment{Exclude low-trust RSUs first}
  \If{$|\mathcal{K}_e|<2f{+}1$}
    \State \textbf{abort}
           \hfill\Comment{Insufficient honest participants}
  \EndIf
  \State \textbf{Step 2: Hash verification}
  \For{$k\in\mathcal{K}_e$}
    \State $\mathbb{1}_{BC}^{(k)}\gets
           \texttt{SC.VerifyModelHash}(\mathbf{W}^{(k)})$
  \EndFor
  \State \textbf{Step 3: Krum geometric filter}
  \State $\tilde{\mathbf{W}}\gets
         \texttt{COORD\_MEDIAN}
         (\{\mathbf{W}^{(k)}\}_{k\in\mathcal{K}_e})$
  \For{$k\in\mathcal{K}_e$}
    \State $\omega_k\gets n_k
           \cdot\mathbb{1}_{BC}^{(k)}
           \cdot\mathbb{1}[d(\mathbf{W}^{(k)},
           \tilde{\mathbf{W}}){<}\gamma]$
  \EndFor
  \State \textbf{Step 4: Weighted aggregation}
  \State $\mathbf{W}_{global}\gets
         {\textstyle\sum_k}\omega_k\mathbf{W}^{(k)}
         /{\textstyle\sum_k}\omega_k$
  \State \texttt{BC.COMMIT}$(H(\mathbf{W}_{global}))$
  \State \textbf{return} $\mathbf{W}_{global}$
\EndProcedure
\end{algorithmic}
\label{alg:brfa_v2}
\end{algorithm}


As shown in Algorithm~\ref{alg:brfa_v2}, the procedure computes 
the coordinate-wise median of all submitted weights as a robust reference, 
verifies each submission's hash against the blockchain commitment 
(Eq.~\ref{eq:bc_model_verify}), and rejects any RSU whose weights 
deviate beyond threshold $\gamma$ from the median (Eq.~\ref{eq:fed_robust}
). The remaining verified weights are averaged to produce the global 
model, whose hash is committed to the blockchain for auditability.

\subsection{Blockchain Architecture Under Zero-Trust
Controller Assumption}
\label{sec:blockchain_redesign}
Figure \ref{fig:figure1} illustrates the two-tier blockchain architecture used in the proposed framework. 
\begin{figure}[H]
    \centering
    \includegraphics[width=1.1\textwidth]{Figure_4_Final.eps}
    \caption{Blockchain Architecture and Trust Management}
    \label{fig:figure1}
\end{figure}
A critical requirement arising from the zero-trust threat model is
that the SDN controller---which is assumed potentially
malicious---must be \emph{excluded from the blockchain write path}.
Permitting the controller to commit flow policies, log FlowMod
installations, or participate in model validation would reduce the
blockchain's integrity guarantee to that of the controller itself,
directly contradicting the zero-trust assumption.As given in the Figure~\ref{fig:figure1} the designed
architecture enforces three principles: (i)~all blockchain writes
originate exclusively from RSUs, which independently observe network
behavior; (ii)~flow-rule validity is established through
\emph{multi-RSU endorsement} rather than controller assertion; and
(iii)~the controller receives verified defense policies \emph{from}
the blockchain rather than writing to it.


\subsubsection{Multi-RSU Flow-Rule Endorsement}
Rather than the controller unilaterally committing the authorized
forwarding policy $\mathcal{P}(s,d)$, each FlowMod installation
must be endorsed by $f+1$ independent neighboring RSUs before
blockchain commitment, where $f$ is the maximum number of Byzantine
RSUs tolerated. Each endorsing RSU $r_j$ independently verifies
that the proposed flow rule is consistent with the current network
topology and established routing policy, then signs the endorsement;
The \ac{RSU} endorsement signature is given in
Equation~\eqref{eq:rsu_endorsement}.
\begin{equation}
\epsilon_j = \textsc{ML-DSA-87.Sign}\!\left(
sk_{r_j},\;
H_{\text{SHA3-512}}\!\left(
\texttt{FlowMod} \,\|\, ts \,\|\,
\mathcal{T}_{r_j}
\right)\right)
\label{eq:rsu_endorsement}
\end{equation}
As given in Equation~\eqref{eq:rsu_endorsement},
$\mathcal{T}_{r_j}$ is \ac{RSU} $r_j$'s
independently maintained topology view, binding
the endorsement to current network state. \ac{MLDSA}
is required because endorsements are written to the
blockchain and must be individually attributable to
specific \ac{RSU}s for the $f{+}1$ endorsement
count to be verifiable by any node.The flow rule is committed to the blockchain
only when $f+1$ valid endorsements are collected,
as formalised in Equation~\eqref{eq:endorsed_commit}:
\begin{equation}
C_{\mathcal{P}} = \textsc{BC.Commit}\!\left(
H(\text{FlowMod}) \,\|\,
\{\epsilon_j\}_{j=1}^{f+1} \,\|\, ts_{commit}\right)
\label{eq:endorsed_commit}
\end{equation}
As shown in Equation \eqref{eq:endorsed_commit}, the
blockchain commitment $C_{\mathcal{P}}$ is formed only when $f+1$
independent endorsement signatures are present alongside the FlowMod
hash and a commit timestamp, meaning a malicious controller issuing
unauthorized FlowMods will fail to collect $f+1$ honest endorsements
and therefore cannot produce a valid $C_{\mathcal{P}}$, directly
enabling Signature~S3 and~S5 detection without relying on the
controller's own reporting.
A malicious controller issuing unauthorized FlowMods will fail to
collect $f+1$ honest endorsements, preventing the malicious rule
from being recorded as legitimate on the blockchain and directly
enabling Signature~S3 and~S5 detection against a malicious
controller without relying on the controller's own reporting.

\subsubsection{RSU-Initiated Blockchain Writes}
All blockchain events---FlowMod installations, ZKP verification
outcomes, trust score updates, and detection alerts---are written
exclusively by RSUs based on independent local observations. An
RSU $r_k$ writes a detection event upon a rule engine or LSTM
detection flag;
The \ac{RSU} blockchain write operation is given
in Equation~\eqref{eq:rsu_write}.
\begin{equation}
\textsc{BC.Write}\!\left(r_k,\;
v_s \,\|\, S_i \,\|\, ts \,\|\,
\textsc{ML-DSA-87.Sign}\!\left(
sk_{r_k},\;
H_{\text{SHA3-512}}\!\left(
v_s \,\|\, S_i \,\|\, ts
\right)\right)\right)
\label{eq:rsu_write}
\end{equation}
As given in Equation~\eqref{eq:rsu_write}, the
\ac{RSU} signs its detection event under its
\ac{MLDSA} private key before writing to the
blockchain, providing non-repudiation for all
logged detection decisions independently of
controller involvement.

As shown in Equation \eqref{eq:rsu_write}, each blockchain record binds the identity of the suspected node $v_s$, the triggered signature index $S_i$, and the event timestamp $ts$ to an ML-DSA-87 signature generated by the reporting RSU $r_k$. Consequently, detection records are cryptographically attributable to a specific RSU, ensuring authenticity, non-repudiation, and resistance to forgery.
The controller has \emph{read-only} access to the blockchain,
receiving verified defense policies through smart contract outputs.
This inverts the trust relationship: the controller's behavior is
audited by the blockchain rather than the blockchain being
populated by the controller.

\subsubsection{Controller Behavior Auditing via RSU Logging}
Since the controller issues FlowMods to RSUs, each RSU independently
logs every received FlowMod to the local RSU-chain immediately upon
receipt, \emph{before} installation;
{the logging operation is formally defined in
Equation \eqref{eq:flowmod_log}:
\begin{equation}
\textsc{BC.LogFlowMod}\!\left(r_k,\;
H(\text{FlowMod}_{recv}) \,\|\, ts_{recv} \,\|\,
\epsilon_{endorsement}\right)
\label{eq:flowmod_log}
\end{equation}
As shown in Equation \eqref{eq:flowmod_log}, the log record
commits the hash of the received FlowMod, the reception timestamp
$ts_{recv}$, and the accompanying endorsement $\epsilon_{endorsement}$
to the RSU-chain before the rule is installed, creating a
tamper-proof pre-installation audit trail that makes it impossible
for a compromised controller to retroactively alter the record of
which FlowMods were issued and when.
A FlowMod arriving at RSU $r_k$ without a corresponding blockchain
commitment $C_{\mathcal{P}}$ from Eq.~\eqref{eq:endorsed_commit}
is immediately flagged as unauthorized;
the unauthorised FlowMod detection flag is formalised in
Equation \eqref{eq:unauth_flowmod}:

To support cross-\ac{RSU} packet identity verification
for Signatures~S6--S8, each \ac{RSU} additionally logs
every packet reception to the \ac{RSU}-chain immediately
upon receipt. The packet reception log entry is formalised
in Equation~\eqref{eq:packet_receipt_log}.
\begin{equation}
\textsc{BC.LogReceipt}\!\left(
H(p),\; dst,\; r_k,\; ts_{recv},\;
\textsc{ML-DSA-87.Verify}(\sigma_{copy},pk_s,m_{copy})
\right)
\label{eq:packet_receipt_log}
\end{equation}
As given in Equation~\eqref{eq:packet_receipt_log},
the log commits the packet hash $H(p)$, the observed
destination $dst$, the receiving \ac{RSU} $r_k$, the
reception timestamp, and the \ac{MLDSA} verification
result on the received copy. Any \ac{RSU} may then
query: $\textsc{BC.QueryReceipt}(H(p), d', W)$ to
determine whether $H(p)$ was received at destination
$d'$ within window $W$ under any \ac{RSU}'s
observation, enabling evaluation of cross-\ac{RSU}
duplication conditions without requiring direct
\ac{RSU}-to-\ac{RSU} communication channels.


\begin{equation}
f_{unauth}(r_k, t) =
\begin{cases}
1 & \text{FlowMod}_{recv} \notin
    \textsc{BC.Query}(\mathcal{C}_{\mathcal{P}}) \\
0 & \text{otherwise}
\end{cases}
\label{eq:unauth_flowmod}
\end{equation}
As given in Equation~\eqref{eq:unauth_flowmod},
$f_{unauth}$ evaluates to 1 whenever the received
\texttt{FlowMod} hash is absent from
$\mathcal{C}_{\mathcal{P}}$. Critically, this flag acts
as a \textbf{pre-installation enforcement gate}: an
\ac{RSU} that evaluates $f_{unauth}(r_k,t){=}1$ rejects
and does not install the \texttt{FlowMod} into its flow
table. The hash is still logged to the blockchain via
Eq.~\eqref{eq:flowmod_log} for audit purposes, but the
malicious rule never takes effect. This converts the
detection flag into a prevention mechanism, ensuring
that controller compromise cannot alter forwarding
behavior even transiently before detection is confirmed.
Because it relies solely on the tamper-proof commitment
record rather than controller self-reporting,
$f_{unauth}$ directly underlies the detection and
prevention of Attacks~1, 3, and~5, where the attacker
actively installs an unauthorized \texttt{FlowMod}. For
Attack~7, $f_{unauth}$ contributes only a corroborating
condition within the S7 signature
(Eq.~\ref{eq:sig_s7}); the primary detection mechanism
there is witness-based monitoring
(Eq.~\ref{eq:bft_penalty}), since the passive variant of
Attack~7 produces no \texttt{FlowMod}-level observable
that is absent from the blockchain.

\subsubsection{Two-Tier Chain Structure}
A local RSU-level permissioned chain (RSU-chain) operates with
low-latency PBFT consensus among neighboring RSUs, tolerating
$f < n/3$ Byzantine RSU nodes. The global chain anchors RSU-chain
Merkle roots periodically;
the anchor hash construction is formally defined in
Equation \eqref{eq:anchor_hash}:
\begin{equation}
H_{anchor}^{(r)} = H\!\left(H_{root}^{RSU} \,\|\,
ts_{anchor} \,\|\, H_{prev}^{global}\right)
\label{eq:anchor_hash}
\end{equation}
As shown in Equation \eqref{eq:anchor_hash}{, each anchor hash
$H_{anchor}^{(r)}$ is computed over the RSU-chain Merkle root
$H_{root}^{RSU}$, the anchor timestamp $ts_{anchor}$, and the
previous global chain hash $H_{prev}^{global}$, forming an
immutable cryptographic link between the local RSU-chain state and
the global chain that limits global consensus overhead to periodic
anchor events while enabling sub-second local logging latency for
real-time detection.}
This limits global consensus overhead to anchor events while
enabling sub-second local logging latency for real-time detection.

\subsubsection{Smart Contract Trust Management}
Each node $v$ maintains a blockchain-recorded trust score
$T_v \in [0,1]$, updated by RSUs via smart contract upon each
ZKP and aggregate signature verification outcome;
the trust score update rule is formally defined in
Equation \eqref{eq:trust_update}:

\begin{equation}
T_v^{(t+1)} =
\begin{cases}
\min(1,\; T_v^{(t)} + \Delta_r) &
\text{if } \textsc{BatchVerify} = 1
\;\land\; b_{\pi} = 1 \\
\max(0,\; T_v^{(t)} - \Delta_p) &
\text{if } \textsc{BatchVerify} = 0
\;\lor\; b_{\pi} = 0
\end{cases}
\label{eq:trust_update}
\end{equation}

As shown in Equation \eqref{eq:trust_update}, the trust score
$T_v^{(t+1)}$ is incremented by reward $\Delta_r$ when both randomized batch verification (Eq.~\ref{eq:batch_verify}) succeeds and both STARK proofs are valid ($b_\pi = 1$) and decremented by penalty $\Delta_p$ when either condition fails, with the $\min$ and $\max$ operators bounding the
score within $[0,1]$ to prevent unbounded accumulation of trust or
distrust from repeated interactions.
where $\Delta_r$ and $\Delta_p$ are the reward and penalty
increments respectively. When $T_v^{(t)} < T_{min}$, the smart
contract automatically emits a \textsc{Quarantine} event,
instructing all RSUs to drop flows originating from $v$;
the quarantine trigger condition is formalised in
Equation \eqref{eq:quarantine}:
\begin{equation}
\textsc{SC.Quarantine}(v) \Leftarrow T_v^{(t)} < T_{min}
\label{eq:quarantine}
\end{equation}
As shown in Equation \eqref{eq:quarantine}, the quarantine
is triggered automatically by the smart contract whenever the trust
score falls below the minimum threshold $T_{min}$, issuing a
network-wide instruction to all RSUs to drop flows originating from
$v$ without requiring any manual intervention or centralised
controller decision, thereby eliminating the single point of failure
inherent in controller-enforced isolation.

\subsubsection{Flow-Rule Commitment Scheme}
At flow setup, the \emph{endorsed} forwarding policy
$\mathcal{P}(s,d)$ is committed to the blockchain by the endorsing
RSUs, not the controller;
the flow-rule commitment is formally defined in
Equation \eqref{eq:policy_commit}:
\begin{equation}
C_{\mathcal{P}} = \textsc{Commit}\!\left(
H(\mathcal{P}(s,d)) \,\|\, ts_{setup} \,\|\,
\{\epsilon_j\}_{j=1}^{f+1}\right)
\label{eq:policy_commit}
\end{equation}
As shown in Equation \eqref{eq:policy_commit}{, the commitment
$C_{\mathcal{P}}$ binds the hash of the authorised forwarding policy
$\mathcal{P}(s,d)$ together with the flow setup timestamp
$ts_{setup}$ and the $f+1$ independent RSU endorsement signatures,
ensuring that the ZKP next-hop proofs defined in
Equation \eqref{eq:stark_hop}{ can reference this commitment as an
immutable ground truth, preventing the controller from secretly
modifying $\mathcal{P}(s,d)$ after flow installation without
detection by honest RSUs.}
ZKP next-hop proofs reference this RSU-endorsed commitment,
preventing the controller from secretly modifying $\mathcal{P}(s,d)$
after flow installation without detection by honest RSUs.

\subsubsection{Trust Architecture Summary}
Table~\ref{tab:trust_roles} summarizes the access rights of each
entity under the corrected zero-trust design.
\begin{table}[H]
\centering
\caption{Blockchain Access Rights Under Zero-Trust Architecture}
\label{tab:trust_roles}
\renewcommand{\arraystretch}{1.1}
\footnotesize
\begin{tabular}{|p{1.6cm}|p{1.4cm}|p{1.4cm}|p{3.0cm}|}
\hline
\textbf{Entity} & \textbf{BC Read} & \textbf{BC Write} &
\textbf{Role} \\
\hline
Controller & \checkmark & \texttimes &
Receives verified policies from SC; FlowMods audited by RSUs \\
\hline
RSUs & \checkmark & \checkmark &
Sole writers; endorse rules; log FlowMods; execute trust updates \\
\hline
Vehicles & \checkmark & \texttimes &
Generate ML-DSA-87 and ZKP proofs; trust scores managed by RSUs \\
\hline
Blockchain & --- & --- &
Ground truth for policies, trust scores, and detection events \\
\hline
\end{tabular}
\end{table}
As shown in Table \ref{tab:trust_roles}, the controller is
the only entity granted read-only access without any write
privileges, inverting the trust relationship of traditional
centralised SDN frameworks where the controller is the primary
authority; RSUs are the sole blockchain writers, ensuring that all
audit records, endorsements, and trust updates originate from
independently observable network nodes rather than from the entity
that is assumed potentially malicious.

The blockchain trust management and mitigation procedure, which
automates trust scoring and node isolation via smart contracts, is
formalized in Algorithm\ref{alg:btmm}.

\begin{algorithm}[htb]
\caption{Blockchain Trust Management and Mitigation
(\texttt{BTMM})}
\begin{algorithmic}[1]
\small
\Procedure{\texttt{BTMM}}{$node,\; attack\_plane$}
  \Comment{$node$ is $v_{atk}$ (data-plane) or
  $c_{atk}$ (control-plane)}
  \If{$attack\_plane = \textsc{ctrl-plane}$}
    % --- Controller mitigation path ---
    \State Submit delay evidence
           $f_{S1}(r,t)$ to blockchain
           \hfill\Comment{Eq.~\ref{eq:delay_evidence}}
    \State \textbf{if} $\mathcal{E}_{delay}(c_{atk},t)
           \geq f{+}1$ \textbf{then}
    \State \quad $T_{c_{atk}} \gets
           \max(0,T_{c_{atk}}-\Delta_p^{ctrl})$
           \hfill\Comment{Eq.~\ref{eq:ctrl_trust_update}}
    \State \quad \texttt{BC.LOG}$(c_{atk},
           \mathcal{E}_{delay},ts_{event})$
    \State \textbf{if} $T_{c_{atk}}
           < T_{min}^{ctrl}$ \textbf{then}
    \State \quad \texttt{SC.Revoke}$(c_{atk})$
           \hfill\Comment{Eq.~\ref{eq:sc_revoke}:
           triggers failover}
  \Else
    % --- Node/vehicle mitigation path ---
    \State $\mathbf{r} \gets
           H_{\text{SHA3-512}}(
           \boldsymbol{\sigma}\|m_1\|\cdots\|m_n)$
    \State $b_{batch} \gets
           \textsc{BatchVerify}(\boldsymbol{\sigma},
           \{pk_i\},\{m_i\},\mathbf{r})$
           \hfill\Comment{Eq.~\ref{eq:batch_verify}}
    \State $b_{\pi} \gets
           \textsc{STARK.Verify}(\pi_{delay},
           C_{delay},H_{\text{SHA3-512}})$
           $\land\;\textsc{STARK.Verify}(\pi_{hop},
           C_{hop},H_{\text{SHA3-512}})$
    \If{$b_{batch} \land b_{\pi}$}
      \State $T_{node} \gets
             \min(1,T_{node}+\Delta_r)$
             \hfill\Comment{Eq.~\ref{eq:trust_update}}
    \Else
      \State $T_{node} \gets
             \max(0,T_{node}-\Delta_p)$
             \hfill\Comment{Eq.~\ref{eq:trust_update}}
      \State \texttt{BC.LOG}$(node,
             \neg b_{batch},\neg b_{\pi},
             ts_{event})$
    \EndIf
    \If{$T_{node} < T_{min}$}
      \State \texttt{SC.Quarantine}$(node)$
             \hfill\Comment{Eq.~\ref{eq:quarantine}}
    \EndIf
  \EndIf
  \State \textbf{return} $T_{node}$ (or $T_{c_{atk}}$)
\EndProcedure
\end{algorithmic}
\label{alg:btmm}
\end{algorithm}

As shown in Algorithm\ref{alg:btmm}, the procedure verifies the
aggregate signature (Eq.\ref{eq:batch_verify}) and both ZKP proofs
for each forwarded packet. Successful verification earns a trust
reward while any failure incurs a larger penalty
(Eq.\ref{eq:trust_update}), with failure events logged immutably.
If a node's trust score falls below the minimum threshold, the
smart contract automatically triggers network-wide quarantine
(Eq.\ref{eq:quarantine}), enforced independently of the
controller.

\subsection{Multi-Controller Zero-Trust Architecture}
\label{sec:multi_ctrl}


A single-controller architecture creates an irrecoverable
single point of failure: controller compromise exposes the
entire network with no honest recovery path. The proposed
framework adopts a flat multi-controller architecture in
which a registered set of \ac{SDVN} controllers operates
without pre-designated primary or backup roles. Dynamic
\ac{RSU} assignment is governed by on-chain trust scores.



\subsubsection{Controller Trust Scoring}
Let $\mathcal{C}{=}\{c_1,\ldots,c_m\}$ denote the
registered controller set. The trusted controller set
at time $t$ is given in
Equation~\eqref{eq:trusted_ctrl_set}.
\begin{equation}
\mathcal{C}_{trusted}(t){=}
\{c_i\in\mathcal{C}:T_{c_i}(t)\geq T_{min}^{ctrl}\}
\label{eq:trusted_ctrl_set}
\end{equation}
As given in Equation~\eqref{eq:trusted_ctrl_set},
$T_{min}^{ctrl}$ is the minimum operational trust
threshold. Controllers below it are excluded from the
trusted set and cannot receive \ac{RSU} assignments.


The controller trust score update is driven by two
independent evidence channels. Conflict evidence
accumulates when RSUs receive \texttt{FlowMod}s absent
from the blockchain-committed endorsed policy
(Eq.~\ref{eq:unauth_flowmod}). Delay evidence
accumulates when independent RSUs assigned to $c_i$
detect timing anomalies attributable to $c_i$'s
\texttt{FlowMod}s via Signature~S1. The delay
evidence count for controller $c_i$ at time $t$ is
given in Equation~\eqref{eq:delay_evidence}.
\begin{equation}
\mathcal{E}_{delay}(c_i,t) =
\left|\left\{
r_k \in \mathcal{R}_{c_i}(t) :
f_{S1}(r_k,t) = 1
\right\}\right|
\label{eq:delay_evidence}
\end{equation}
As given in Equation~\eqref{eq:delay_evidence},
$\mathcal{R}_{c_i}(t)$ is the set of \ac{RSU}s
currently assigned to controller $c_i$
(Eq.~\ref{eq:rsu_ctrl_assign}). A threshold of
$f{+}1$ independent \ac{RSU}s is required before
delay evidence triggers a trust penalty, preserving
the \ac{BFT} structure and preventing a single
compromised \ac{RSU} from falsely accusing a
controller. The controller trust score update is
given in Equation~\eqref{eq:ctrl_trust_update}.


\begin{equation}
T_{c_i}^{(t+1)}{=}
\begin{cases}
\min(1,T_{c_i}^{(t)}{+}\Delta_r^{ctrl}) &
\text{conflict evidence}{<}f{+}1
\;\land\;
\mathcal{E}_{delay}(c_i,t){<}f{+}1
\\
\max(0,T_{c_i}^{(t)}{-}\Delta_p^{ctrl}) &
\text{conflict evidence}{\geq}f{+}1
\;\lor\;
\mathcal{E}_{delay}(c_i,t){\geq}f{+}1
\end{cases}
\label{eq:ctrl_trust_update}
\end{equation}

As given in Equation~\eqref{eq:ctrl_trust_update},
the trust score decrements when either conflict
evidence (unauthorized \texttt{FlowMod}s) or delay
evidence (Eq.~\ref{eq:delay_evidence}, $f{+}1$
independent S1 detections across assigned \ac{RSU}s)
reaches threshold. This closes the critical gap
where a timing-manipulating controller issuing
topologically valid \texttt{FlowMod}s accumulates
conflict evidence of zero and is incorrectly rewarded;
the delay evidence channel makes the controller trust
system responsive to timing behavior independently
of \texttt{FlowMod} content.




\subsubsection{Dynamic RSU-Controller Assignment}
The \ac{RSU}-to-controller assignment restricted to
trusted controllers is given in
Equation~\eqref{eq:rsu_ctrl_assign}.
\begin{equation}
c^*(k,t){=}
\underset{c_i\in\mathcal{C}_{trusted}(t)}
{\arg\min}\;d(r_k,c_i)
\label{eq:rsu_ctrl_assign}
\end{equation}
As given in Equation~\eqref{eq:rsu_ctrl_assign},
$d(r_k,c_i)$ is the geographic distance metric. Every
\ac{RSU} is always assigned to a controller with a
blockchain-verified behavioral history.



\subsubsection{Trusted Controller Failover via SC-Revoke}
The smart contract revocation trigger is given in
Equation~\eqref{eq:sc_revoke}.
\begin{equation}
\textsc{SC.Revoke}(c_i)\Leftarrow
T_{c_i}^{(t)}<T_{min}^{ctrl}
\label{eq:sc_revoke}
\end{equation}
As given in Equation~\eqref{eq:sc_revoke}, revocation
updates $\mathcal{C}_{trusted}$ on-chain and broadcasts
a \textsc{ControllerRevoked} event. Each affected \ac{RSU}
immediately re-executes the trusted failover given in
Equation~\eqref{eq:ctrl_failover}.
\begin{equation}
c^*(k,t^+){=}
\underset{c_j\in
\mathcal{C}_{trusted}(t)\setminus\{c_i\}}
{\arg\min}\;d(r_k,c_j)
\label{eq:ctrl_failover}
\end{equation}
As given in Equation~\eqref{eq:ctrl_failover}, failover
is always to a controller currently in the trusted set
with a verified behavioral history, never to a
pre-designated backup that could itself be compromised.



\subsubsection{Distributed Key Generation}
The controller is removed entirely from the key plane.
\ac{ZKP} proving and verification keys are generated via
a \ac{DKG} ceremony among \ac{RSU} ring members. The
blockchain-committed verification key is given in
Equation~\eqref{eq:vk_commit}.
\begin{equation}
\textsc{BC.Commit}\!\left(vk_{ZKP},
\{\text{Com}_j\}_{j=1}^{n_{RSU}},
ts_{setup}\right)
\label{eq:vk_commit}
\end{equation}
As given in Equation~\eqref{eq:vk_commit}, $vk_{ZKP}$
is the jointly generated verification key,
$\{\text{Com}_j\}$ are the per-\ac{RSU} \ac{DKG}
commitments, and $ts_{setup}$ is the ceremony timestamp.
A compromised controller gains zero key-plane influence.



\subsubsection{ZKP Proving Key Rotation}
When \ac{RSU} $r_j$ is revoked, all proving keys it
participated in generating must be rotated. The rotation
trigger is given in
Equation~\eqref{eq:key_rotation_trigger}.
\begin{equation}
\textsc{SC.RotateKeys}\Leftarrow
\exists\,r_j:T_{r_j}^{(t)}<T_{min}
\land r_j\in\mathcal{P}_{DKG}^{current}
\label{eq:key_rotation_trigger}
\end{equation}
As given in Equation~\eqref{eq:key_rotation_trigger},
$\mathcal{P}_{DKG}^{current}$ is the set of \ac{RSU}s
that participated in the current ceremony. The rotated
key commitment is given in
Equation~\eqref{eq:vk_commit_rotated}.
\begin{equation}
\textsc{BC.Commit}\!\left(vk_{ZKP}^{new},
\{\text{Com}_j\}_{j\in\mathcal{P}_{DKG}^{new}},
ts_{rotation}\right)
\label{eq:vk_commit_rotated}
\end{equation}
As given in Equation~\eqref{eq:vk_commit_rotated},
$\mathcal{P}_{DKG}^{new}{=}
\mathcal{P}_{DKG}^{current}\setminus\{r_j\}$.



% b...
% \hl{Figure1}

% \hl{Figure 2}


% \hl{figure 3} 


% \hl{Figure 4}


\section{Data Collection}
The attacks will be simulated in ns-3 and data will be collected using simulated enviroment.



\chapter{Perfomance Evaluvation and Benchmarking}
% ...
\section{Baseline Selection and Benchmarking Rationale}
\label{sec:baselines}

\subsection{External Baselines}
Three recent baselines are selected that directly address selective
time delay or hidden forwarding attacks in SDN or SDVN.

\textbf{B1 — TAP}~\cite{Arsalan2018}: Closest prior art for
selective packet timing delay detection (Variants~1--3).
Represents the best physics-based, mobility-aware timing approach;
comparison exposes its failure under hidden forwarding attacks and
compromised controller conditions.

\textbf{B2 — SFTO-Guard}~\cite{Tang2023SFTO}: Most recent dedicated
defense against slow-rate TCAM exhaustion (Variants~3--4). Comparison
isolates the performance degradation of static ML-based thresholds
under vehicular topology dynamics.

\textbf{B3 — FADE}~\cite{Zhang2021FADE}: State-of-the-art
flow-conservation-based forwarding anomaly detector covering packet
duplication and path deviation (Variants~5--8). Comparison directly
demonstrates the short-observation-window limitation
(Eq.~\eqref{eq:observation_window}) of flow-conservation approaches
under vehicular mobility.

\subsubsection{Internal Ablation Baselines}
Thirteen component-ablated configurations
(AB1-A through AB13-B) are defined in
Section~\ref{sec:ablation}. Each experiment
isolates one formulated design component by
replacing or removing it while holding all
others fixed; the resulting performance change
is attributed solely to that component.

\section{Performance Metrics}
\label{sec:metrics}

Twelve metrics are defined to evaluate all formulated
components of the proposed methodology, covering detection
quality, attack-class-specific effectiveness, operational
overhead, resilience of each security layer, and privacy
preservation. Metrics~M5, M8, M9, and~M12 are evaluated
exclusively in ablation studies since no state-of-the-art
baseline implements the corresponding components; Metric~M7
is reported for the proposed framework with baseline figures
cited from published work as context only.

% ----------------------------------------------------------------
% M1 — Detection Quality and Mobility Robustness
% ----------------------------------------------------------------

\textbf{M1 — Detection Quality and Mobility Robustness:}
The Matthews Correlation Coefficient is computed both
as an aggregate topline figure and stratified across
attack variants, detection modes, and mobility regimes.
The base \ac{MCC} formula is given in
Equation~\eqref{eq:mcc} (reproduced for completeness).
\begin{equation}
\text{MCC} =
\frac{TP \cdot TN - FP \cdot FN}
{\sqrt{(TP{+}FP)(TP{+}FN)(TN{+}FP)(TN{+}FN)}}
\label{eq:mcc}
\end{equation}



As given in Equation~\eqref{eq:mcc}, MCC produces a
scalar in $[-1,1]$ reliable under the severe class
imbalance of attack detection datasets.
The per-variant, per-mode score is given in
Equation~\eqref{eq:mcc_variant}.
\begin{equation}
\text{MCC}_{s,m} = \text{MCC evaluated over variant }
s \in \{S1,\ldots,S8\},\;
m \in \{\text{OBU},\text{RSU}\}
\label{eq:mcc_variant}
\end{equation}
As given in Equation~\eqref{eq:mcc_variant}, computing
a separate score per variant and per mode prevents
strong variants from masking weak ones in aggregate
figures and exposes any mode-specific detection gap.
To validate the mobility-adjusted thresholds
(Eq.~\ref{eq:mobility_baseline}) against the static
thresholds used by all three baselines, the mobility-
stratified score is given in
Equation~\eqref{eq:mcc_mobility}.
\begin{equation}
\text{MCC}_{s}[\rho_b, \bar{v}_b] =
\text{MCC}_{s}\text{ computed when }
\rho(t) \in \rho_b \;\land\;
\bar{v}(t) \in \bar{v}_b
\label{eq:mcc_mobility}
\end{equation}
As given in Equation~\eqref{eq:mcc_mobility},
$\rho_b \in \{\text{low, medium, high}\}$ and
$\bar{v}_b \in \{\text{low, high}\}$ are vehicle
density and speed bins derived from the SUMO mobility
traces. A drop in $\text{MCC}_{s}$ at high-density,
high-speed conditions indicates that the mobility
adjustment fails, while a stable score across all bins
validates Equation~\eqref{eq:mobility_baseline}.
The \ac{FPR} is evaluated as a hard constraint rather
than a separate top-level metric: hyperparameters are
rejected if the resulting FPR exceeds 1\% in either mode.


% ----------------------------------------------------------------
% M2 — Safety-Critical Threshold Violation Rate (TVR)
% ----------------------------------------------------------------

\textbf{M2 — Safety-Critical Threshold Violation Rate (TVR):}
Definitively evaluates Selective Time Delay Attack
containment (Variants~1--4). Unlike L$_{e2e}$, TVR
is not diluted by best-effort traffic or legitimate
handoff latency. The TVR is given in
Equation~\eqref{eq:tvr}.
\begin{equation}
\text{TVR} =
\frac{\left|\left\{
p \in \mathcal{P}_{crit} :
\delta_p(v,r,t) > \Delta_{max}
\right\}\right|}{|\mathcal{P}_{crit}|}
\label{eq:tvr}
\end{equation}
As given in Equation~\eqref{eq:tvr},
$\mathcal{P}_{crit} = \{p :
\text{Priority}(p) = \textsc{High}\}$ is the set
of safety-critical packets, and $\delta_p(v,r,t)$
is the per-hop forwarding delay anchored to the
distributed time reference $T_{ref}(t)$
(Eq.~\ref{eq:delay_updated}), so that handoff-induced
latency is normalised out before the threshold
comparison. Packets that are never forwarded within
the observation window are assigned $\delta_p =
\infty > \Delta_{max}$, capturing delay-equivalent
non-delivery within the same formula
(Eq.~\ref{eq:nfwd_detect}). Under no attack,
TVR $\approx 0$; under Variants~1--4, TVR rises with
attack intensity and returns toward 0 post-quarantine.
Compared against baselines B1--B3, none of which
enforce per-packet delay thresholds on safety-critical
traffic.


% ----------------------------------------------------------------
% M3 — Unauthorized Copy Rate (UCR)
% ----------------------------------------------------------------

\textbf{M3 — Unauthorized Copy Rate (UCR):}
Definitively evaluates Hidden Forwarding Attack
containment (Variants~5--8), including passive variants
where \ac{PDR} remains at 100\% and attack-induced
overhead is zero, making UCR the only metric that
exposes Variants~7--8. The UCR is given in
Equation~\eqref{eq:ucr}.
\begin{equation}
\text{UCR} =
\frac{\left|\left\{
p \in \mathcal{P}_{total} :
\exists\, d' \notin \mathcal{P}(s_p, d_p),\;
p \in \mathcal{R}(d', W)
\right\}\right|}{|\mathcal{P}_{total}|}
\label{eq:ucr}
\end{equation}
As given in Equation~\eqref{eq:ucr},
$\mathcal{P}(s_p, d_p)$ is the blockchain-committed
authorised forwarding policy (Eq.~\ref{eq:policy_commit}),
and $\mathcal{R}(d', W)$ is the set of packet hashes
received at $d'$ within window $W$. UCR $= 0$ exactly
under no attack; UCR $> 0$ under Variants~5--8;
and UCR returns to 0 post-quarantine. For active
variants (5, 6), the copy carries modified content
confirmed by \textsc{ML-DSA-87.Verify}$= 0$; for
passive variants (7, 8), the copy is unmodified
but $d' \notin \mathcal{P}(s_p, d_p)$ remains
detectable via the blockchain-committed policy.
Compared against baselines B1--B3, none of which
monitor unauthorized reception at off-path destinations.


% ----------------------------------------------------------------
% M4 — Node/Flow Mitigation Latency
% ----------------------------------------------------------------

\textbf{M4 — Node/Flow Mitigation Latency
($L_{mit}$):}
Measures the elapsed time from attack onset to
successful smart contract quarantine of the offending
node, directly evaluating the responsiveness of the
trust management pipeline
(Eq.~\ref{eq:trust_update}, \ref{eq:quarantine}).
The mitigation latency is given in
Equation~\eqref{eq:l_mit}.
\begin{equation}
L_{mit} = t_{quarantine} - t_{onset}
\label{eq:l_mit}
\end{equation}
As given in Equation~\eqref{eq:l_mit},
$t_{onset}$ is the timestamp of the first
attack-conforming packet and $t_{quarantine}$ is the
timestamp at which \textsc{SC.Quarantine}$(v)$ is
executed by the smart contract. Target: $L_{mit}
\leq 100$\,ms. Compared against baseline mitigation mechanisms:
B1 (TAP) reports emergency packet duplication
reduction via the Controller-Defaulter-List rather
than an explicit quarantine-latency
figure~\cite{Arsalan2018}; B2 (SFTO-Guard) reports
eviction correction times under rule-count
prediction~\cite{Tang2023SFTO}; B3 (FADE) reports no
explicit node isolation mechanism, blocking attack
flows via flow-rule removal
only~\cite{Zhang2021FADE}.


% ----------------------------------------------------------------
% M5 — Controller Failover Latency (ABLATION ONLY)
% ----------------------------------------------------------------

\textbf{M5 — Controller Failover Latency
($L_{failover}$) [Ablation only]:}
Evaluates the multi-controller zero-trust architecture
(Eq.~\ref{eq:sc_revoke}, \ref{eq:ctrl_failover}).
No state-of-the-art baseline implements trusted
dynamic controller failover; all three baselines assume
a permanently available trusted controller.
The failover latency is given in
Equation~\eqref{eq:l_failover}.
\begin{equation}
L_{failover} =
\underset{k \in \mathcal{K}_{affected}}{\max}
\left(
t_{reassign}^{(k)} - t_{revoke}
\right)
\label{eq:l_failover}
\end{equation}
As given in Equation~\eqref{eq:l_failover},
$t_{revoke}$ is the timestamp at which
\textsc{SC.Revoke}$(c_i)$ fires, and
$t_{reassign}^{(k)}$ is the timestamp at which
\ac{RSU} $r_k$ completes reassignment to its new
trusted controller (Eq.~\ref{eq:ctrl_failover}).
The worst-case (maximum) over all affected RSUs
$\mathcal{K}_{affected}$ is reported to capture
stragglers under geographic dispersion. Target:
$L_{failover} \leq 100$\,ms. Evaluated across
0, 1, and 2 compromised controller scenarios.


% ----------------------------------------------------------------
% M6 — Overall System Latency (L_e2e)
% ----------------------------------------------------------------

\textbf{M6 — Overall System Latency ($L_{e2e}$):}
Measures the mean per-packet latency accumulated
across all hops including the framework's own
cryptographic, detection, and blockchain processing
overhead, benchmarked against the 100\,ms
safety-critical bound. The end-to-end latency is
given in Equation~\eqref{eq:l_e2e}.
\begin{equation}
L_{e2e} =
\frac{1}{|\mathcal{P}_{total}|}
\sum_{p \in \mathcal{P}_{total}}
\left(
t_{recv}^{final}(p) - t_{send}(p)
\right)
\label{eq:l_e2e}
\end{equation}
As given in Equation~\eqref{eq:l_e2e},
$t_{send}(p)$ is the transmission timestamp at the
source and $t_{recv}^{final}(p)$ is the reception
timestamp at the final authorized destination,
both anchored to $T_{ref}(t)$
(Eq.~\ref{eq:delay_updated}). Evaluated under
attack and post-mitigation and compared against
all three baselines to quantify the latency cost
of the framework's security mechanisms relative
to simpler defenses.


% ----------------------------------------------------------------
% M7 — Security Processing and Consensus Overhead
%        (Ablation/context only)
% ----------------------------------------------------------------

\textbf{M7 — Security Processing and Consensus
Overhead [Ablation/context only]:}
Combines the cryptographic communication overhead
per packet, the RSU verification processing time
as a function of batch size, and the blockchain
consensus latency for FlowMod endorsement, jointly
validating the $O(n)$ parallelism claim
(Eq.~\ref{eq:batch_verify}) and the endorsement
design (Eq.~\ref{eq:endorsed_commit}).
The per-packet cryptographic communication overhead
is given in Equation~\eqref{eq:o_crypto}.
\begin{equation}
O_{crypto}^{(n)} =
n \cdot |\sigma_i|
+ |\pi_{delay}| + |\pi_{hop}|
\label{eq:o_crypto}
\end{equation}
As given in Equation~\eqref{eq:o_crypto},
$n$ is the number of hops, $|\sigma_i| =
4{,}595$\,bytes is the \ac{MLDSA} signature
size, and $|\pi_{delay}|$, $|\pi_{hop}|
\leq 100$\,KB each are the \ac{STARK} proof
sizes (Eq.~\ref{eq:overhead_full}).
The \ac{RSU} batch verification processing time
as a function of batch size $B$ is given in
Equation~\eqref{eq:t_verify}.
\begin{equation}
T_{verify}(B) =
T_{batch}(B) + T_{STARK}
\label{eq:t_verify}
\end{equation}
As given in Equation~\eqref{eq:t_verify},
$T_{batch}(B)$ is the wall-clock time for
\textsc{BatchVerify} over $B$ packets
(Eq.~\ref{eq:batch_verify}) and $T_{STARK}$
is the \textsc{STARK.Verify} latency per packet.
The blockchain endorsement consensus latency
is given in Equation~\eqref{eq:t_consensus}.
\begin{equation}
T_{consensus} =
t_{commit} - t_{FlowMod\_recv}
\label{eq:t_consensus}
\end{equation}
As given in Equation~\eqref{eq:t_consensus},
$t_{commit}$ is the timestamp at which
$C_\mathcal{P}$ (Eq.~\ref{eq:endorsed_commit})
is written to the blockchain after $f{+}1$
endorsements are collected, and
$t_{FlowMod\_recv}$ is the arrival timestamp
of the FlowMod at the first endorsing \ac{RSU}.
All three sub-metrics are reported as a function
of network load; \ac{FTISCON} and \ac{BlockREV}
figures cited as literature context for M7 only.


% ----------------------------------------------------------------
% M8 — Federated Model Poisoning Resistance
%        (ABLATION ONLY)
% ----------------------------------------------------------------

\textbf{M8 — Federated Model Poisoning Resistance
[Ablation only]:}
Evaluates the Byzantine-robust aggregation layer
(Algorithm~\ref{alg:brfa_v2}) by measuring
degradation of the global model's detection
quality as the fraction of malicious RSUs
submitting poisoned gradients increases from 0
to the PBFT Byzantine bound $f < n/3$.
The poisoning degradation is given in
Equation~\eqref{eq:delta_poison}.
\begin{equation}
\Delta_{poison}(\rho_{mal}) =
\text{MCC}_{clean} -
\text{MCC}(\rho_{mal})
\label{eq:delta_poison}
\end{equation}
As given in Equation~\eqref{eq:delta_poison},
$\text{MCC}_{clean}$ is the global model's MCC
under zero malicious RSUs, and
$\text{MCC}(\rho_{mal})$ is the MCC when a
fraction $\rho_{mal}$ of the 64 RSUs submit
gradient-poisoned updates. $\Delta_{poison}$
is evaluated for BRFA-v2 versus naive FedAvg to
isolate the contribution of trust-gating
(Step~1), hash verification (Step~2), and Krum
filtering (Step~3) in Algorithm~\ref{alg:brfa_v2}.
Target: $\Delta_{poison}(\rho_{mal}) \approx 0$
for $\rho_{mal} < 1/3$.


% ----------------------------------------------------------------
% M9 — Distributed Time Reference Robustness
%        (ABLATION ONLY)
% ----------------------------------------------------------------

\textbf{M9 — Distributed Time Reference Robustness
[Ablation only]:}
Validates the formal guarantee of
Equation~\eqref{eq:time_consensus} that $T_{ref}(t)$
remains within honest clock range when up to
$f < n_{RSU}/2$ RSU clocks are compromised.
The time reference deviation under $f_{bad}$
compromised clocks is given in
Equation~\eqref{eq:eps_ref}.
\begin{equation}
\varepsilon_{ref}(f_{bad}) =
\left|
T_{ref}(t) - T_{ground}(t)
\right|
\text{ with } f_{bad}
\text{ RSU clocks offset by } \delta_{attack}
\label{eq:eps_ref}
\end{equation}
As given in Equation~\eqref{eq:eps_ref},
$T_{ground}(t)$ is the \ac{GPS} ground truth
timestamp available from the simulation,
$f_{bad}$ is the number of RSU clocks
artificially offset by a fixed attack value
$\delta_{attack}$, and $\varepsilon_{ref}$
is the absolute deviation of the computed
median from ground truth. $\varepsilon_{ref}$
is swept over $f_{bad} \in \{0, 1,
\lfloor n_{RSU}/4 \rfloor,
\lfloor n_{RSU}/2 \rfloor - 1\}$ to
confirm the bound fails exactly at
$f_{bad} = n_{RSU}/2$ and holds for all
smaller values, validating
Equation~\eqref{eq:time_consensus}.


% ----------------------------------------------------------------
% M10 — Privacy Leakage / Raw Data Exposure
% ----------------------------------------------------------------

\textbf{M10 — Privacy Leakage Score ($L_{priv}$):}
Directly evaluates the privacy preservation
claim — Gap~4 in the literature review —
by quantifying the volume of raw re-identifying
data categories transmitted outside the local
RSU trust boundary, comparing the proposed
framework against baselines that centralize
raw metadata. The privacy leakage score is
given in Equation~\eqref{eq:l_priv}.
\begin{equation}
L_{priv} =
\sum_{d \in \mathcal{D}_{raw}}
w_d \cdot
\mathbb{1}\!\left[
d \text{ transmitted beyond local trust boundary}
\right]
\label{eq:l_priv}
\end{equation}
As given in Equation~\eqref{eq:l_priv},
$\mathcal{D}_{raw} = \{$location history,
vehicle identity, trajectory, raw packet
headers, raw flow metadata$\}$ is the set
of re-identifying data categories, and
$w_d \in [0,1]$ is the privacy-sensitivity
weight of each category (location history
and vehicle identity are assigned $w_d = 1$
as highest sensitivity; flow metadata
$w_d = 0.5$). This is an architectural
evaluation computed from each baseline's
published design rather than a runtime
output: MOBIGUARD exposes only model weight
deltas, hash commitments, and proof
outcomes ($L_{priv} = 0$ for the five
raw categories); HSA centralizes raw packet headers ($L_{priv} = w_{headers}
+ w_{trajectory}$); TAP centralizes only the
vehicle ID of a detected attacker at the
controller via the Controller-Defaulter-List
($L_{priv} = w_{vehicle\_identity}$).


% ----------------------------------------------------------------
% M11 — Unauthorized FlowMod Containment Rate (UFCR)
% ----------------------------------------------------------------

\textbf{M11 — Unauthorized FlowMod Containment
Rate (UFCR):}
Directly evaluates the multi-RSU blockchain
endorsement mechanism
(Eqs.~\ref{eq:rsu_endorsement},
\ref{eq:endorsed_commit},
\ref{eq:unauth_flowmod}) against control-plane
attacks (Variants~1, 3, 5, 7). UFCR measures
the fraction of unauthorized FlowMod commands
successfully blocked before RSU installation
among all FlowMods flagged as unauthorized by
$f_{unauth}$ (Eq.~\ref{eq:unauth_flowmod}).
The UFCR is given in
Equation~\eqref{eq:ufcr}.
\begin{equation}
\text{UFCR} =
\frac{\left|\left\{
FM \in \mathcal{FM}_{unauth} :
FM \text{ blocked before installation}
\right\}\right|}{|\mathcal{FM}_{unauth}|}
\label{eq:ufcr}
\end{equation}
As given in Equation~\eqref{eq:ufcr},
$\mathcal{FM}_{unauth} = \{FM :
f_{unauth}(r_k, t) = 1\}$ is the set of
FlowMods arriving without a valid
blockchain commitment $C_\mathcal{P}$
(Eq.~\ref{eq:endorsed_commit}). Under no
attack, $|\mathcal{FM}_{unauth}| = 0$ so
UFCR is undefined; under control-plane
Variants~1, 3, 5, 7, target UFCR $= 1.0$,
meaning all unauthorised FlowMods are
blocked. FTISCON and BlockREV report
partial containment figures in their
published evaluations and are cited as
context; baselines B1--B3 report UFCR $= 0$
since none verify FlowMods before installation.


% ----------------------------------------------------------------
% M12 — Witness Alert Precision and Recall (WAP-R)
%         (ABLATION ONLY)
% ----------------------------------------------------------------

\textbf{M12 — Witness Alert Precision and Recall
(WAP-R) [Ablation only]:}
Evaluates the witness-based passive hidden
forwarding detection mechanism
(Eqs.~\ref{eq:dup_alert_cond},
\ref{eq:bft_penalty}) specifically against
Variants~7 and~8, which are undetectable by
ZKP alone and where UCR captures the outcome
but not the witness mechanism's accuracy.
Witness alert precision is given in
Equation~\eqref{eq:wap}.
\begin{equation}
P_W =
\frac{TP_W}{TP_W + FP_W}
\label{eq:wap}
\end{equation}
As given in Equation~\eqref{eq:wap},
$TP_W$ is the count of $2f{+}1$-threshold
alert events that correctly identify a true
passive hidden forwarding instance, and
$FP_W$ is the count of threshold events
triggered by a false witness coalition.
Witness alert recall is given in
Equation~\eqref{eq:war}.
\begin{equation}
R_W =
\frac{TP_W}{TP_W + FN_W}
\label{eq:war}
\end{equation}
As given in Equation~\eqref{eq:war},
$FN_W$ is the count of true passive
forwarding instances for which fewer than
$2f{+}1$ valid alerts were submitted
within window $W$ — representing cases
where witness density was insufficient.
$P_W$ validates the false-accusation
resistance of the $2f{+}1$ threshold
(Eq.~\ref{eq:bft_penalty}), and $R_W$
validates the feasibility claim that
$f{=}1$ (requiring 3 witnesses) is
achievable under mean vehicle density of
$\approx 3.1$ vehicles per RSU cell. No
baseline has a witness mechanism;
WAP-R is reported as ablation context only.


% \subsection{Primary Performance Metrics}
% Four metrics are selected to evaluate both detection quality and
% mitigation effectiveness.

% \textbf{M1 — Matthews Correlation Coefficient (MCC)}: Primary
% detection quality metric, selected over accuracy and F1 because it
% remains reliable under the class imbalance inherent in attack
% detection datasets where benign packets vastly outnumber attack
% packets; the MCC is formally defined in
% Equation \ref{eq:mcc}:
% \begin{equation}
% \text{MCC} = \frac{TP \cdot TN - FP \cdot FN}
% {\sqrt{(TP+FP)(TP+FN)(TN+FP)(TN+FN)}}
% \label{eq:mcc}
% \end{equation}
% As shown in Equation \eqref{eq:mcc}, the MCC is computed from
% all four cells of the confusion matrix: true positives $TP$, true
% negatives $TN$, false positives $FP$, and false negatives
% $FN$, producing a single scalar in $[-1, 1]$ where $+1$ indicates
% perfect classification, $0$ indicates performance no better than
% random guessing, and $-1$ indicates total misclassification; unlike
% accuracy, which can reach high values simply by predicting the
% majority benign class, and unlike F1, which ignores true negatives
% entirely, MCC remains informative even when attack packets constitute
% a small minority of the total traffic.


% \textbf{M2 — Per-Variant Detection Rate (DR)}: True positive rate
% evaluated independently for each of the eight attack variants,
% ensuring no variant is masked by aggregate accuracy figures.

% \textbf{M3 — False Positive Rate (FPR)}: Critical in vehicular
% safety contexts where false alarms trigger mitigation actions that
% themselves degrade network performance; evaluated separately for
% lightweight and full modes.

% \textbf{M4 — Mitigation Latency} ($L_{mit}$): Time from attack
% onset to smart contract quarantine execution
% (Eq.~\eqref{eq:quarantine}), measured against the 100\,ms
% safety-critical latency bound; directly evaluates mitigation
% effectiveness rather than detection alone.

% \textbf{M5 — Packet Delivery Ratio (PDR)}: Fraction of
% safety-critical packets successfully delivered to the final
% destination across the full multihop path, evaluated under active
% attack with and without MOBIGUARD. PDR directly captures the combined
% impact of both selective time delay (packets dropped after timeout)
% and hidden forwarding (packets corrupted or rerouted) on end-to-end
% delivery across all intermediate hops.

% \textbf{M6 — End-to-End Latency} ($L_{e2e}$): Mean per-packet
% latency accumulated across all hops from source to final destination,
% compared against the 100\,ms safety-critical bound. Evaluated under
% attack and post-mitigation to quantify the latency restoration
% effectiveness of the proposed framework relative to the three external
% baselines.

% \textbf{M7 — Safety-Critical Threshold Violation Rate (TVR):}
% Existing metrics L$_{e2e}$ (M6) measure mean latency over all
% packets, diluting extreme delays in safety-critical traffic with
% normal traffic and conflating legitimate handoff latency with
% attack-induced delay. TVR instead measures the fraction of
% safety-critical packets whose per-hop forwarding delay
% definitively exceeds the safety threshold $\Delta_{max}$ ---
% the precise condition constituting a successful Selective Time
% Delay Attack. The TVR is given in Equation~\eqref{eq:tvr}.
% \begin{equation}
% \text{TVR} = \frac{\left|\left\{
% p \in \mathcal{P}_{crit} :
% \delta_p(v,r,t) > \Delta_{max}
% \right\}\right|}{|\mathcal{P}_{crit}|}
% \label{eq:tvr}
% \end{equation}
% As given in Equation~\eqref{eq:tvr},
% $\mathcal{P}_{crit} = \{p : \text{Priority}(p) =
% \textsc{High}\}$ is the set of safety-critical packets,
% and $\delta_p(v,r,t)$ is the per-hop forwarding delay
% anchored to the distributed time reference $T_{ref}(t)$
% (Eq.~\ref{eq:delay_updated}), ensuring that
% handoff-induced delay is normalised out before the
% threshold comparison. Under no attack, TVR $\approx 0$;
% under Attack Variants~1--4, TVR rises directly
% proportional to attack intensity; and post-quarantine,
% TVR returns toward 0. TVR is evaluated per attack variant
% and per attack percentage, and compared against all three
% external baselines (B1--B3) to quantify their failure
% to suppress threshold violations under mobility.

% \textbf{M8 — Unauthorized Copy Rate (UCR):}
% \ac{PDR} (M5) measures delivery to the authorized
% destination, which remains at 100\% under passive hidden
% forwarding (Attack Variants~7--8) because the original
% packet path is unaffected --- the attack is entirely
% invisible to PDR. UCR instead directly measures the
% fraction of transmitted packets for which at least one
% copy was received at an unauthorized destination,
% definitively capturing all four hidden forwarding variants.
% The UCR is given in Equation~\eqref{eq:ucr}.
% \begin{equation}
% \text{UCR} = \frac{\left|\left\{
% p \in \mathcal{P}_{total} :
% \exists\, d' \notin \mathcal{P}(s_p, d_p),\;
% p \in \mathcal{R}(d', W)
% \right\}\right|}{|\mathcal{P}_{total}|}
% \label{eq:ucr}
% \end{equation}
% As given in Equation~\eqref{eq:ucr},
% $\mathcal{P}(s_p, d_p)$ is the blockchain-committed
% authorised forwarding policy for packet $p$ from
% source $s_p$ to destination $d_p$
% (Eq.~\ref{eq:policy_commit}),
% $\mathcal{R}(d', W)$ is the set of packet hashes
% received at node $d'$ within observation window $W$,
% and $\mathcal{P}_{total}$ is the total set of
% transmitted packets. Under no attack, UCR $= 0$ exactly;
% under Attack Variants~5--8, UCR $> 0$; and
% post-quarantine, UCR returns to 0. For active hidden
% forwarding variants (5, 6), the unauthorized copy
% carries modified content confirmed by
% $\textsc{ML-DSA-87.Verify} = 0$; for passive
% variants (7, 8), the content is unmodified but the
% unauthorized destination $d' \notin \mathcal{P}(s_p, d_p)$
% remains detectable via the blockchain-committed policy,
% making UCR the only metric in the set that
% definitively exposes passive hidden forwarding.
% UCR is evaluated per attack variant and per attack
% percentage, and compared against all three external
% baselines (B1--B3).

\section{Ablation Study}
\label{sec:ablation}

Thirteen ablation experiments are designed to isolate the
contribution of each individual formulated component of the
proposed framework. In each experiment, one component is
replaced or removed while all others remain active; the
resulting performance change is attributed solely to that
component. The x variable in every experiment is categorical
— configuration type — rather than a continuous parameter
(parameter sensitivity is addressed separately). All
applicable metrics from Section~\ref{sec:metrics} are
evaluated per experiment.

% ----------------------------------------------------------------
% AB1 — Dual-Mode Architecture
% ----------------------------------------------------------------
\subsection{AB1 — Dual-Mode Detection Architecture}
\label{sec:ab1}

\textbf{Component isolated:}
The dual-mode architecture pairing a lightweight rule-based
engine on \ac{OBU}s (Algorithm~\ref{alg:lrad_obu}) with a
full federated \ac{LSTM} detector on \ac{RSU}s
(Algorithm~\ref{alg:lrad_rsu}).

\textbf{Configurations compared:}
\begin{itemize}
    \item \textbf{AB1-A (Rule-only):} \texttt{LRAD-OBU}
    evaluates S1, S2-partial, S3, S4 locally; no escalation
    to RSU; no \ac{LSTM} inference.
    \item \textbf{AB1-B (LSTM-only):} All packets forwarded
    directly to RSU; \texttt{LRAD-OBU} disabled; full
    \ac{LSTM} inference without rule pre-filtering.
    \item \textbf{AB1-C (Full dual-mode):} Complete proposed
    framework — \texttt{LRAD-OBU} pre-filters, escalates to
    \texttt{LRAD-RSU} and BRFA-v2 \ac{LSTM}.
\end{itemize}

\textbf{x variable:} Configuration
$\in \{\text{AB1-A}, \text{AB1-B}, \text{AB1-C}\}$ (categorical).

\textbf{y metrics evaluated:} M1 (per-variant, per-mode MCC),
M2 (TVR for Variants~1--4), M3 (UCR for Variants~5--8),
M4 (mitigation latency). AB1-A is expected to miss
mobility-masked variants; AB1-B is expected to incur higher
latency due to the absence of OBU pre-filtering.


% ----------------------------------------------------------------
% AB2 — Federated vs Centralised LSTM
% ----------------------------------------------------------------
\subsection{AB2 — Federated vs Centralised LSTM}
\label{sec:ab2}

\textbf{Component isolated:}
The federated learning design in which each of the 64
\ac{RSU}s trains a local model and shares only weight
updates (Algorithm~\ref{alg:brfa_v2}), versus centralising
all raw flow data at a single training server.

\textbf{Configurations compared:}
\begin{itemize}
    \item \textbf{AB2-A (Centralised):} All 64 RSUs
    transmit raw per-packet feature logs to a central
    training server; one global \ac{LSTM} is trained
    on the pooled dataset.
    \item \textbf{AB2-B (Federated — proposed):}
    Each \ac{RSU} trains locally; only weight updates
    with hash commitments are transmitted; BRFA-v2
    aggregation applies.
\end{itemize}

\textbf{x variable:} Training architecture
$\in \{\text{AB2-A}, \text{AB2-B}\}$ (categorical).

\textbf{y metrics evaluated:} M1 (MCC — confirms federated
matches centralised detection quality), M3 (UCR), M10
(privacy leakage score — AB2-A exposes raw flow metadata,
AB2-B does not). M10 is the primary differentiating metric
for this experiment, as it quantifies the privacy cost of
centralisation directly.


% ----------------------------------------------------------------
% AB3 — ZKP Failure Features in LSTM Input
% ----------------------------------------------------------------
\subsection{AB3 — ZKP Failure Indicators as LSTM Features}
\label{sec:ab3}

\textbf{Component isolated:}
The two novel binary \ac{ZKP} failure indicators
$\mathbb{1}[\pi_{delay}{=}\bot]$ and
$\mathbb{1}[\pi_{hop}{=}\bot]$ in the \ac{LSTM} input
vector (Eq.~\ref{eq:lstm_input}), which feed cryptographic
evidence of attacks directly into the detection layer.

\textbf{Configurations compared:}
\begin{itemize}
    \item \textbf{AB3-A (5-feature baseline):}
    Input vector contains $[\delta_t, \lambda_{PI},
    U_{TCAM}, \rho_t, \bar{v}_t]$ only — no
    cryptographic feedback.
    \item \textbf{AB3-B (7-feature proposed):}
    Full input vector including both \ac{ZKP}
    failure indicators (Eq.~\ref{eq:lstm_input}).
\end{itemize}

\textbf{x variable:} Feature configuration
$\in \{\text{AB3-A}, \text{AB3-B}\}$ (categorical).

\textbf{y metrics evaluated:} M1 (per-variant MCC,
particularly for Variants~5--8 where ZKP failures are the
primary cryptographic signal), M3 (UCR). The improvement
from AB3-A to AB3-B isolates the contribution of linking
the cryptographic layer to the learning layer.


% ----------------------------------------------------------------
% AB4 — Individual STARK Proof Contributions
% ----------------------------------------------------------------
\subsection{AB4 — Individual STARK Proof Contributions}
\label{sec:ab4}

\textbf{Component isolated:}
The two \ac{STARK} proof components — timing compliance
$\pi_{delay}$ (Eq.~\ref{eq:stark_delay}) and next-hop
legitimacy $\pi_{hop}$ (Eq.~\ref{eq:stark_hop}) —
evaluated in isolation and in combination to identify
which proof drives detection of which attack class.

\textbf{Configurations compared:}
\begin{itemize}
    \item \textbf{AB4-A (No ZKP):} ML-DSA-87 signatures
    and batch verification retained; both \ac{STARK}
    proofs removed entirely.
    \item \textbf{AB4-B ($\pi_{delay}$ only):}
    Timing compliance proof active; hop legitimacy
    proof removed.
    \item \textbf{AB4-C ($\pi_{hop}$ only):}
    Hop legitimacy proof active; timing compliance
    proof removed.
    \item \textbf{AB4-D (Both proofs — proposed):}
    Complete \ac{STARK} layer active.
\end{itemize}

\textbf{x variable:} Proof configuration
$\in \{\text{AB4-A}, \text{AB4-B},
\text{AB4-C}, \text{AB4-D}\}$ (categorical).

\textbf{y metrics evaluated:} M1 (per-variant MCC),
M2 (TVR — expected sensitive to $\pi_{delay}$),
M3 (UCR — expected sensitive to $\pi_{hop}$),
M7 ($T_{verify}$ — overhead reduction without each proof).


% ----------------------------------------------------------------
% AB5 — Byzantine-Robust Aggregation vs FedAvg
% ----------------------------------------------------------------
\subsection{AB5 — Byzantine-Robust Aggregation}
\label{sec:ab5}

\textbf{Component isolated:}
The four-step BRFA-v2 procedure (trust gate, hash
verification, Krum filtering, weighted aggregation —
Algorithm~\ref{alg:brfa_v2}) versus standard FedAvg
under a fixed adversarial RSU fraction of
$\rho_{mal} = 0.25$ (16 of 64 RSUs submit poisoned
gradients), representing a realistic worst-case within
the PBFT bound.

\textbf{Configurations compared:}
\begin{itemize}
    \item \textbf{AB5-A (Naive FedAvg):} Standard
    weighted averaging with no trust gate, no hash
    verification, no Krum filtering.
    \item \textbf{AB5-B (BRFA-v2 — proposed):}
    Full four-step Byzantine-robust procedure.
\end{itemize}

\textbf{x variable:} Aggregation algorithm
$\in \{\text{AB5-A}, \text{AB5-B}\}$ (categorical).

\textbf{y metrics evaluated:} M8 ($\Delta_{poison}$
at $\rho_{mal} = 0.25$), M1 (MCC — does the global
model's detection quality hold under poisoning?).


% ----------------------------------------------------------------
% AB6 — Witness Monitoring Mechanism
% ----------------------------------------------------------------
\subsection{AB6 — Witness-Based Forwarding Verification}
\label{sec:ab6}

\textbf{Component isolated:}
The witness monitoring mechanism for passive hidden
forwarding detection
(Eqs.~\ref{eq:dup_alert_cond},~\ref{eq:bft_penalty}),
which is the primary detector for Variants~7 and~8
where \ac{ZKP} corroboration alone is insufficient.
Eq.~\ref{eq:nfwd_detect} detects delay-equivalent
non-delivery and is formally classified as a severe
case of Selective Time Delay Variants~1--4; it does
not target hidden forwarding and is excluded from
this ablation's scope.


\textbf{Configurations compared:}
\begin{itemize}
    \item \textbf{AB6-A (No witness monitoring):}
    Passive hidden forwarding detection relies solely
    on $b_{hop} = 0$ as the corroborating indicator;
    no witness alert submission or BFT threshold.
    \item \textbf{AB6-B (Full witness mechanism —
    proposed):} Witness vehicles log $(H(p), dst)$,
    submit $\alpha_w$ or $\beta_w$ directly on-chain,
    and \textsc{SC.PenalizeTrust} fires at $2f{+}1$
    verified alerts.
\end{itemize}

\textbf{x variable:} Witness mechanism
$\in \{\text{AB6-A}, \text{AB6-B}\}$ (categorical).

\textbf{y metrics evaluated:} M3 (UCR for Variants~7--8),
M12 (WAP-R precision and recall — defined only for
AB6-B; AB6-A reports UCR only). The UCR gap between
AB6-A and AB6-B directly quantifies the contribution
of witness monitoring to passive attack containment.


% ----------------------------------------------------------------
% AB7 — Blockchain Trust Scoring and Quarantine
% ----------------------------------------------------------------
\subsection{AB7 — Blockchain Trust Scoring and
Automated Quarantine}
\label{sec:ab7}

\textbf{Component isolated:}
The smart contract trust score update and automated
quarantine pipeline (Eqs.~\ref{eq:trust_update},
\ref{eq:quarantine}). Without this component,
attacks are detected but the attacking node is never
isolated, leaving the network vulnerable to continued
attacks.

\textbf{Configurations compared:}
\begin{itemize}
    \item \textbf{AB7-A (Detection-only):}
    \texttt{LRAD-RSU} and \ac{LSTM} detection active;
    trust score update and \textsc{SC.Quarantine}
    disabled — detection events are logged but no
    node is isolated.
    \item \textbf{AB7-B (Full pipeline — proposed):}
    Complete trust update and quarantine pipeline active.
\end{itemize}

\textbf{x variable:} Quarantine mechanism
$\in \{\text{AB7-A}, \text{AB7-B}\}$ (categorical).

\textbf{y metrics evaluated:} M1 (MCC — sustained
attacks after initial detection reduce ongoing MCC in
AB7-A), M4 (mitigation latency — undefined for AB7-A
since no quarantine occurs), M6 ($L_{e2e}$ — without
quarantine, attack overhead accumulates over time).


% ----------------------------------------------------------------
% AB8 — Multi-RSU FlowMod Endorsement
% ----------------------------------------------------------------
\subsection{AB8 — Multi-RSU FlowMod Endorsement}
\label{sec:ab8}

\textbf{Component isolated:}
The $f{+}1$ independent \ac{RSU} endorsement requirement
before blockchain FlowMod commitment
(Eqs.~\ref{eq:rsu_endorsement},
\ref{eq:endorsed_commit}, \ref{eq:unauth_flowmod}).
Without endorsement, the controller can unilaterally
commit arbitrary FlowMods.

\textbf{Configurations compared:}
\begin{itemize}
    \item \textbf{AB8-A (No endorsement):}
    FlowMods committed to blockchain on controller
    assertion alone — no RSU endorsement signatures
    collected; $f_{unauth}$ cannot fire.
    \item \textbf{AB8-B (Full endorsement —
    proposed):} $f{+}1$ RSU endorsements required;
    unauthorized FlowMods blocked before installation.
\end{itemize}

\textbf{x variable:} Endorsement mechanism
$\in \{\text{AB8-A}, \text{AB8-B}\}$ (categorical).

\textbf{y metrics evaluated:} M11 (UFCR — zero in
AB8-A by definition; target 1.0 in AB8-B), M4
(mitigation latency for control-plane attacks),
M7 ($T_{consensus}$).


% ----------------------------------------------------------------
% AB9 — Multi-Controller vs Single Controller
% ----------------------------------------------------------------
\subsection{AB9 — Multi-Controller Zero-Trust
Architecture}
\label{sec:ab9}

\textbf{Component isolated:}
The flat multi-controller set with SC-Revoke trusted
failover (Eqs.~\ref{eq:trusted_ctrl_set},
\ref{eq:sc_revoke}, \ref{eq:ctrl_failover}) versus
a single centralised controller, directly testing the
single-point-of-failure elimination claim.

\textbf{Configurations compared:}
\begin{itemize}
    \item \textbf{AB9-A (Single controller):}
    One fixed controller; no trust scoring; no failover.
    When the controller is compromised, the framework
    operates with no control-plane oversight.
    \item \textbf{AB9-B (Multi-controller —
    proposed):} Flat registered controller set with
    on-chain trust scoring, SC-Revoke, and dynamic
    RSU reassignment.
\end{itemize}

\textbf{x variable:} Controller architecture
$\in \{\text{AB9-A}, \text{AB9-B}\}$ (categorical),
evaluated under one compromised controller.

\textbf{y metrics evaluated:} M1 (MCC — does detection
degrade when the controller is compromised?), M4
(mitigation latency), M5 (controller failover latency —
undefined for AB9-A since no failover exists), M11
(UFCR — AB9-A cannot block unauthorized FlowMods from
a compromised single controller).


% ----------------------------------------------------------------
% AB10 — DKG vs Controller-Managed Key Generation
% ----------------------------------------------------------------
\subsection{AB10 — Distributed Key Generation vs
Controller-Managed Keys}
\label{sec:ab10}

\textbf{Component isolated:}
The \ac{DKG} ceremony among \ac{RSU} ring members
(Eq.~\ref{eq:vk_commit}) that removes the controller
from the key plane entirely, versus a design where
the controller issues \ac{ZKP} proving keys.

\textbf{Configurations compared:}
\begin{itemize}
    \item \textbf{AB10-A (Controller key authority):}
    A compromised controller issues false proving keys
    that allow arbitrary behaviour to be proven
    compliant — no \ac{DKG} ceremony.
    \item \textbf{AB10-B (\ac{DKG} — proposed):}
    Keys generated jointly among honest \ac{RSU}s;
    blockchain-committed $vk_{ZKP}$ is the only
    accepted verification key.
\end{itemize}

\textbf{x variable:} Key generation architecture
$\in \{\text{AB10-A}, \text{AB10-B}\}$ (categorical),
evaluated under a compromised controller.

\textbf{y metrics evaluated:} M3 (UCR — can a
controller with false proving keys suppress valid
$\pi_{hop}$ proofs and enable undetected hidden
forwarding?), M11 (UFCR — can false keys allow
unauthorized FlowMods to appear compliant?).


% ----------------------------------------------------------------
% AB11 — ZKP Key Rotation on RSU Revocation
% ----------------------------------------------------------------
\subsection{AB11 — ZKP Proving Key Rotation on
RSU Revocation}
\label{sec:ab11}

\textbf{Component isolated:}
The key rotation trigger mechanism
(Eqs.~\ref{eq:key_rotation_trigger},
\ref{eq:vk_commit_rotated}) that initiates a fresh
\ac{DKG} ceremony when a revoked \ac{RSU} was a
\ac{DKG} participant, preventing retained key
material from being exploited post-revocation.

\textbf{Configurations compared:}
\begin{itemize}
    \item \textbf{AB11-A (No rotation):} Revoked
    \ac{RSU} retains its \ac{DKG} key material;
    existing $vk_{ZKP}$ remains in use after
    revocation.
    \item \textbf{AB11-B (With rotation —
    proposed):} \textsc{SC.RotateKeys} triggers
    immediately on revocation; fresh \ac{DKG}
    among remaining honest \ac{RSU}s; new
    $vk_{ZKP}^{new}$ committed to blockchain.
\end{itemize}

\textbf{x variable:} Key rotation policy
$\in \{\text{AB11-A}, \text{AB11-B}\}$ (categorical),
evaluated at $t = \{1, 2, 3, 4, 5\}$ time steps
after RSU revocation.

\textbf{y metrics evaluated:} M3 (UCR — does the
revoked RSU exploit retained keys to produce valid
$\pi_{hop}$ proofs enabling hidden forwarding in
AB11-A?), M1 (MCC for hidden forwarding variants
S5--S8 — does detection quality hold post-revocation
in both configurations?).


% ----------------------------------------------------------------
% AB12 — Blockchain Write Restriction
% ----------------------------------------------------------------
\subsection{AB12 — Blockchain Write Restriction}
\label{sec:ab12}

\textbf{Component isolated:}
The architectural constraint that only \ac{RSU}s
may write to the blockchain
(Eqs.~\ref{eq:rsu_write}, \ref{eq:flowmod_log}),
directly testing the zero-trust principle that the
controller is excluded from the audit trail.

\textbf{Configurations compared:}
\begin{itemize}
    \item \textbf{AB12-A (Controller write access):}
    A compromised controller is permitted to write
    to the blockchain — it can retroactively insert
    false endorsement records and manipulate the
    audit trail.
    \item \textbf{AB12-B (RSU-only writes —
    proposed):} Controller has read-only access;
    all writes originate from \ac{RSU}s via
    \ac{PBFT} endorsement.
\end{itemize}

\textbf{x variable:} Write access policy
$\in \{\text{AB12-A}, \text{AB12-B}\}$ (categorical),
evaluated under a compromised controller.

\textbf{y metrics evaluated:} M11 (UFCR — can a
controller with write access retroactively legitimize
unauthorized FlowMods?), M1 (MCC — does the integrity
of detection event logs hold when a compromised
controller can alter the blockchain?), M4 (mitigation
latency — can a controller with write access cancel
quarantine instructions?).


% ----------------------------------------------------------------
% AB13 — Mobility-Adjusted vs Static Thresholds
% ----------------------------------------------------------------
\subsection{AB13 — Mobility-Adjusted vs Static
Detection Thresholds}
\label{sec:ab13}

\textbf{Component isolated:}
The mobility-adjusted baseline delay estimator
$\bar{\delta}_r(t) = \delta_0 + \alpha_\rho \cdot
\rho(t) + \alpha_v \cdot \bar{v}(t)^{-1}$
(Eq.~\ref{eq:mobility_baseline}) and \ac{EWMA}
variance estimator (Eq.~\ref{eq:ewma_variance}),
which together constitute the paper's direct
response to Gap~1 (mobility-induced vulnerability).

\textbf{Configurations compared:}
\begin{itemize}
    \item \textbf{AB13-A (Static threshold):}
    Delay baseline replaced with a fixed constant
    $\bar{\delta}_{static}$; variance replaced
    with a fixed $\sigma_{static}$, reproducing
    the threshold design of all three baselines
    (TAP, SFTO-Guard, FADE).
    \item \textbf{AB13-B (Mobility-adjusted —
    proposed):} Full $\bar{\delta}_r(t)$ and
    $\sigma_r(t)$ estimators active.
\end{itemize}

\textbf{x variable:} Threshold design
$\in \{\text{AB13-A}, \text{AB13-B}\}$ (categorical).

\textbf{y metrics evaluated:} M1
($\text{MCC}_s[\rho_b, \bar{v}_b]$ —
Eq.~\ref{eq:mcc_mobility}, the mobility-stratified
score which directly tests whether detection quality
holds across density and speed regimes), M2 (TVR —
do legitimate handoff delays cause false threshold
violations in AB13-A?). This experiment is the
direct ablation analog of the mobility-adjustment
robustness metric M1 and together they constitute
the complete evaluation of the paper's Gap~1 claim.

\section{Benchmarking Against State-of-the-Art}
\label{sec:benchmarking}

Five benchmarking experiments are designed to evaluate
the proposed framework against baselines B1 (TAP),
B2 (SFTO-Guard), and B3 (FADE) at the system level.
All eight attack variants operate simultaneously in
every experiment unless otherwise stated. Metrics
applicable only in ablation studies (M5, M8, M9, M12)
are excluded from all benchmarking experiments.
Default settings are held fixed for any variable not
under variation in that experiment.

% ================================================================
% EXPERIMENT 1 — Attack Penetration x Attack Intensity
% ================================================================

\subsection{Experiment~1 — Attack Penetration Under
Varying Attack Intensity}
\label{sec:exp1}

\textbf{Motivation:} Directly answers RQ1 by evaluating
how detection quality, threshold violations, and
mitigation responsiveness change as the number of
attacking nodes increases under three distinct levels
of attack intensity, capturing both the breadth
(how many attackers) and the severity (how hard each
attacks) of the threat.

\textbf{x variable — Attack Penetration $p$:}
$p \in \{0\%, 20\%, 40\%, 60\%, 80\%, 100\%\}$,
six equidistant data points. At each value of $p$,
$\lfloor 0.01p \times 264 \rfloor$ attacker nodes are
allocated among the 200 vehicles and 64 \ac{RSU}s
plane-specifically per attack variant; controller
compromise follows the threshold rule defined in
Section~\ref{sec:simulation} automatically.

\textbf{Second dimension — Attack Intensity $\mathcal{I}$:}
Three discrete levels produce three lines per metric
graph. Attack intensity is defined separately for each
attack class.

For Selective Time Delay (Variants~1--4), intensity is
the induced per-hop delay as a multiple of $\Delta_{max}$,
given in Equation~\eqref{eq:intensity_td}.
\begin{equation}
\mathcal{I}_{TD} \in
\left\{
1.1\cdot\Delta_{max},\;
2\cdot\Delta_{max},\;
4\cdot\Delta_{max}
\right\}
\label{eq:intensity_td}
\end{equation}
As given in Equation~\eqref{eq:intensity_td},
$1.1\cdot\Delta_{max}$ represents a sub-threshold
stealth injection just above the detection boundary,
$2\cdot\Delta_{max}$ represents a moderate injection,
and $4\cdot\Delta_{max}$ represents an aggressive
injection causing severe desynchronization.

For Hidden Forwarding (Variants~5--8), intensity is
the fraction of eligible flows that are secretly
forwarded, given in Equation~\eqref{eq:intensity_hf}.
\begin{equation}
\mathcal{I}_{HF} \in
\left\{25\%,\; 50\%,\; 100\%\right\}
\;\text{of eligible flows}
\label{eq:intensity_hf}
\end{equation}
As given in Equation~\eqref{eq:intensity_hf},
$25\%$ represents stealthy selective duplication,
$50\%$ moderate exposure, and $100\%$ full
forwarding of all eligible flows.

\textbf{Graph format:} Three lines per metric (one
per intensity level), each with six equidistant
data points at $p \in \{0\%, 20\%, 40\%, 60\%,
80\%, 100\%\}$.

\textbf{y metrics evaluated:}
M1 (MCC), M2 (TVR), M3 (UCR), M4 ($L_{mit}$),
M6 ($L_{e2e}$), M10 ($L_{priv}$), M11 (UFCR).

\textbf{Baselines:} B1 (TAP), B2 (SFTO-Guard),
B3 (FADE). Each baseline is expected to degrade
significantly above $p = 20\%$ and under controller
compromise, which is automatically activated at
$p \geq 33\%$ per the threshold rule.


% ================================================================
% EXPERIMENT 2 — Vehicle Mobility Speed
% ================================================================

\subsection{Experiment~2 — Detection Quality Under
Varying Vehicle Mobility}
\label{sec:exp2}

\textbf{Motivation:} Directly answers RQ2 by
evaluating whether the mobility-adjusted threshold
design maintains detection quality across the full
speed range from near-stationary urban crawl to
high-speed highway, the exact regime where
static-threshold baselines are expected to fail
due to handoff-masking of delay signatures.

\textbf{x variable — Mean Vehicle Speed $\bar{v}$:}
$\bar{v} \in \{10, 60, 100, 140\}$\,km/h,
four data points corresponding to urban crawl,
urban arterial, highway, and high-speed highway.
The x variable is given in
Equation~\eqref{eq:speed_x}.
\begin{equation}
\bar{v}_{SUMO} = \frac{\bar{v}_{km/h}}{3.6}
\;\text{m/s} \;\in
\{2.78,\; 16.67,\; 27.78,\; 38.89\}\;\text{m/s}
\label{eq:speed_x}
\end{equation}
As given in Equation~\eqref{eq:speed_x}, all speeds
are converted from km/h to m/s for \ac{SUMO}
configuration. To ensure all vehicles travel at
exactly the target speed without inter-vehicle speed
variation, the following \ac{SUMO} \texttt{vType}
parameters are set: \texttt{speedFactor="1.0"} and
\texttt{speedDev="0.0"}. Both the \texttt{vType}
\texttt{maxSpeed} attribute and the road/edge
\texttt{speed} attribute are set to the same m/s
value, ensuring every vehicle traverses every road
segment at precisely $\bar{v}_{SUMO}$.

\textbf{Experimental rationale:} At $\bar{v} = 10$\,km/h,
handoff latency is minimal (long zone residence time
of $\approx 90$--$220$\,s), providing a near-ideal
detection baseline. At $\bar{v} = 140$\,km/h,
zone residence time drops to $\approx 7.7$\,s,
maximizing handoff-induced jitter and the
Bhattacharyya distance overlap
(Eq.~\ref{eq:bhattacharyya}). The mobility-adjusted
baseline $\bar{\delta}_r(t)$ (Eq.~\ref{eq:mobility_baseline})
should compensate for this jitter; static-threshold
baselines cannot.

\textbf{y metrics evaluated:}
M1 ($\text{MCC}_s[\rho_b, \bar{v}_b]$ —
Eq.~\ref{eq:mcc_mobility}), M2 (TVR), M3 (UCR),
M4 ($L_{mit}$), M6 ($L_{e2e}$), M11 (UFCR).


% ================================================================
% EXPERIMENT 3 — Network Scalability
% ================================================================

\subsection{Experiment~3 — System Performance Under
Network Scalability}
\label{sec:exp3}

\textbf{Motivation:} Directly answers RQ3 and RQ4
by evaluating whether detection quality, mitigation
responsiveness, and cryptographic overhead remain
within safety-critical bounds as the network scales
from 100 to 400 vehicles over the same fixed
geographic area and \ac{RSU} infrastructure.

\textbf{x variable — Number of Vehicles $N_v$:}
$N_v \in \{100, 200, 300, 400\}$, four equidistant
data points. The 64 \ac{RSU}s on the 8$\times$8
grid and 4 controllers remain fixed across all
points; only the vehicle count scales, increasing
vehicle density from $\approx 23$ to $\approx 90$
vehicles per km$^2$ over the
2061\,m$\times$2137\,m simulation area.
The vehicle density is given in
Equation~\eqref{eq:density_x}.
\begin{equation}
\rho(N_v) =
\frac{N_v}{A_{map}}
\;\text{veh/m}^2,
\quad A_{map} = 2061 \times 2137\;\text{m}^2
\label{eq:density_x}
\end{equation}
As given in Equation~\eqref{eq:density_x},
$\rho(N_v)$ increases proportionally with $N_v$,
raising both legitimate FlowMod installation rates
and the volume of detection events written to the
blockchain, stressing the BRFA-v2 aggregation load,
the blockchain consensus pipeline, and the
\textsc{BatchVerify} processing budget simultaneously.
Vehicle type ratios are maintained proportionally
across all four points.

\textbf{y metrics evaluated:}
M1 (MCC), M2 (TVR), M3 (UCR), M4 ($L_{mit}$),
M6 ($L_{e2e}$), M7 ($T_{verify}(B)$ and
$T_{consensus}$ — the two overhead sub-metrics
most sensitive to network scale), M11 (UFCR).


% ================================================================
% EXPERIMENT 4 — Attack Observable Evidence Index (AOEI)
% ================================================================

\subsection{Experiment~4 — Detection Under Varying
Attack Observable Evidence}
\label{sec:exp4}

\textbf{Motivation:} Directly answers RQ5 by
evaluating whether the framework detects attacks at
the low-evidence extreme — the "low-rate stealthy"
operating regime explicitly stated in the threat
model — where every baseline fails because
statistical evidence in the network is deliberately
minimized. This experiment captures two defining
attack characteristics simultaneously: the
\textit{selectivity} of the attacker (how broadly
it targets eligible traffic) and the
\textit{copy destination divergence} (how far the
unauthorized destination is from the authorized
path in hop count, directly affecting witness
monitoring coverage).

\textbf{x variable — Attack Observable Evidence
Index (AOEI):}
AOEI is a composite index combining both attack
characteristics on a unified scale. The AOEI is
given in Equation~\eqref{eq:aoei}.
\begin{equation}
\text{AOEI} =
\frac{1}{2}
\left(
\frac{n_{targeted}}{n_{eligible}}
+
\frac{1}{1 + d_{div}/d_{max}}
\right)
\label{eq:aoei}
\end{equation}
As given in Equation~\eqref{eq:aoei},
$n_{targeted}/n_{eligible}$ is the
\textbf{selectivity ratio}: the fraction of
eligible high-priority packets (for Variants~1--4)
or eligible flows (for Variants~5--8) that the
attacker actively targets. $d_{div}$ is the
\textbf{copy destination divergence}: the
topological hop distance between the unauthorized
destination $d'$ and the authorized next-hop on
the legitimate forwarding path. $d_{max}$ is the
maximum hop distance across the simulated network
topology. The second term $1/(1 + d_{div}/d_{max})$
is a normalized proximity score: when $d' $ is
adjacent to the authorized path ($d_{div} = 0$),
this term equals 1 (maximum witness observability);
when $d'$ is at maximum network divergence, this
term approaches 0 (minimum witness observability).
AOEI therefore equals 1.0 when the attack leaves
maximum observable evidence (high selectivity,
copy destination within witness range) and
approaches 0 when the attack is maximally
stealthy (low selectivity, copy destination
outside witness range).

The four equidistant data points are given in
Equation~\eqref{eq:aoei_points}.
\begin{equation}
\text{AOEI} \in
\{0.25,\; 0.50,\; 0.75,\; 1.0\}
\label{eq:aoei_points}
\end{equation}
As given in Equation~\eqref{eq:aoei_points}, each
data point is realized in simulation by jointly
setting the selectivity ratio and $d_{div}$:
\begin{itemize}
    \item \textbf{AOEI = 0.25:}
    $n_{targeted}/n_{eligible} = 0.10$,
    $d_{div} = 0.8 \cdot d_{max}$.
    Maximally stealthy — 10\% of eligible
    traffic targeted, unauthorized copy sent
    near maximum divergence from authorized path.

    \item \textbf{AOEI = 0.50:}
    $n_{targeted}/n_{eligible} = 0.25$,
    $d_{div} = 0.5 \cdot d_{max}$.
    Low-moderate evidence — 25\% targeting,
    mid-range divergence.

    \item \textbf{AOEI = 0.75:}
    $n_{targeted}/n_{eligible} = 0.75$,
    $d_{div} = 0.2 \cdot d_{max}$.
    High evidence — 75\% targeting, copy
    destination relatively close to authorized path.

    \item \textbf{AOEI = 1.0:}
    $n_{targeted}/n_{eligible} = 1.0$,
    $d_{div} = 0$ (copy destination is adjacent
    to authorized path).
    Maximum evidence — full targeting, unauthorized
    copy within direct witness radio range.
\end{itemize}

\textbf{Note:} For Variants~1--4 (Selective Time
Delay), $d_{div}$ does not apply since there is
no copy destination; only the selectivity ratio
component contributes, and AOEI is computed as
$n_{targeted}/n_{eligible}$ only. This is made
explicit in results reporting.

\textbf{y metrics evaluated:}
M1 (MCC), M2 (TVR), M3 (UCR), M4 ($L_{mit}$),
M12 (WAP-R — most sensitive to $d_{div}$ dimension
since witness recall degrades as $d_{div}$ increases).


% ================================================================
% EXPERIMENT 5 — Per-Attack Variant Table
% ================================================================

\subsection{Experiment~5 — Per-Attack Variant
Performance Under Default Settings}
\label{sec:exp5}

\textbf{Motivation:} Provides a single-table
system-level comparison of per-variant detection
performance under default settings, giving
reviewers the attack-class-specific breakdown
that Experiments~1--4 aggregate away. This
experiment produces no graph — results are
presented as a table.

\textbf{Setup:} Default settings throughout.
All 8 attack variants evaluated independently
in separate simulation runs (one variant active
at a time) to produce per-variant scores.

\textbf{Table structure:} One row per attack
variant S1--S8. Columns: M1 (MCC per variant,
all frameworks), M2 (TVR — reported for S1--S4
only; N/A for S5--S8), M3 (UCR — reported for
S5--S8 only; N/A for S1--S4), M4 ($L_{mit}$,
all frameworks), M6 ($L_{e2e}$, all frameworks).
Each baseline (B1--B3) reports values only for
the attack variants it was designed to address;
cells are marked \textit{not evaluated} for
out-of-scope variants, directly exposing the
coverage gaps that Gap~2 in the literature
review identified.





\section{Simulation settings}
\label{sec:simulation}

Prior to all simulation runs, a one-time \ac{DKG}
ceremony (Eq.~\ref{eq:vk_commit}) is executed among
all 64 RSUs to generate the \ac{STARK} proving and
verification keys. No controller participates in the
ceremony. The resulting verification key $vk_{ZKP}$
is committed to the blockchain and treated as a fixed
pre-simulation assumption throughout all runs.
Key rotation (Eq.~\ref{eq:key_rotation_trigger}) is
triggered only when an RSU revocation event occurs
during simulation and is evaluated as part of the
multi-controller trust management assessment in
Section~\ref{sec:multi_ctrl}.

\subsection{Network and mobility environment}
We selected a dense urban area in Los Angeles, California for simulation by importing mobility traces using SUMO and OpenStreetMap, while the network was simulated using Network Simulator 3 (NS-3.35). The SDVN architecture follows a Vehicle$\rightarrow$RSU$\rightarrow$Controller hierarchy: vehicles and RSUs communicate over IEEE 802.11p DSRC for the data plane, while RSU-to-controller links use Ethernet (CSMA, 1000 Mbps) for the control plane. Four controllers were set in a $2\times2$ grid, rather than a single centralized controller, so that no individual controller compromise can disable the entire control plane - directly addressing the single-point-of-failure limitation of centralized SDVN designs motivating this work.  We simulate 200 vehicles (100 cars, 25 buses, 25 lorries, 25 vans, 25 trucks) and 64 RSUs deployed on an 8$\times$8 grid with 260 m horizontal and 270 m vertical spacing over a 2061 m $\times$ 2137 m mobility trace, with a maximum vehicle speed of 150 kmph. To prevent unnatural vehicle clustering and ensure realistic spatial distribution, SUMO trip generation was configured with a minimum route distance of 1000 m and a high random-routing factor, forcing vehicles to take diverse, long-distance routes across the entire map \cite{sumo_randomtrips}. The DSRC transmission power is set to 41 dBm with 0 dB transmit/receive gain, giving a maximum reliable link range of 270 m under the Cost231 propagation loss model (constant-speed propagation delay); this value was selected to remain within the standard DSRC power class (23--44 dBm)\cite{FCC03-324} while providing inter-RSU coverage consistently, so that every grid cell remains within at least one RSU's reliable range. Data transmission and routing decisions occur at 1 Hz, with UDP packets of 750 bytes transmitted per flow. This shares the routing-update and data-transmission clock, set below the SAE J2735 Basic Safety Message rate of 10 Hz \cite{sae_j2735} to limit DSRC channel load and routing-recomputation overhead at 200 vehicles and 64 RSUs, while remaining frequent enough to detect Selective Time Delay and Hidden Forwarding signatures within their 9--22 s observation window.

\begin{figure}[H]
    \centering
    \includegraphics[width=0.95\textwidth]{sumo.eps}
    \caption{SUMO simulation environment}
    \label{fig:Sumo}
\end{figure}

Figure \ref{fig:Sumo} shows the SUMO simulation window running the generated road network, with the live network parameter panel displaying key metrics such as 200 loaded vehicles.

\begin{figure}[H]
    \centering
    \includegraphics[width=0.75\textwidth]{netanim.eps}
    \caption{NetAnim visualization with RSU placement}
    \label{fig:Netanim}
\end{figure}
Figure \ref{fig:Netanim} shows the NetAnim animator displaying the simulated network topology, with vehicle nodes (green), Roadside Units (RSUs, yellow), and four SDVN controllers (CTRL-1 to CTRL-4, purple) positioned across the deployment area according to the imported SUMO mobility trace, illustrating the spatial distribution of nodes used for evaluating the proposed routing scheme.

\subsection{Proposed dual-mode detection architecture}
The dual-mode detection architecture parameters are established as follows:

\begin{itemize}
    \item \textbf{Lightweight rule engine (OBU):} Hop-delay detection threshold $= 50$\,ms; forwarding-ratio deviation threshold $= 20$\%.
    \item \textbf{Signature coverage:} Signatures~S1 and S2-partial are evaluated locally
at the \ac{OBU} using mobility-adjusted thresholds;
Signatures~S3 and S4 are evaluated at the \ac{RSU}
using \ac{RSU}-observable TCAM metrics (FlowMod rate,
PACKET\_IN rate per source); Signatures~S2-full and
S5--S8 are evaluated at the \ac{RSU} via the
cryptographic pipeline, since
they require \ac{ZKP} timing ($\pi_{\text{delay}}$)
and hop-legitimacy ($\pi_{\text{hop}}$) proof inputs
unavailable on resource-constrained \ac{OBU}s
(Section~\ref{sec:rule_based}).

     \item \textbf{Mobility baseline parameters
    $\delta_0$, $\alpha_\rho$, $\alpha_v$
    (Eq.~\ref{eq:mobility_baseline}):}
    Calibrated from benign-only ($0\%$ attack)
    SUMO traces via least-squares regression over
    $\rho(t)$ and $\bar{v}(t)$. Robustness
    evaluated by sweeping $\pm\{10\%, 20\%, 30\%\}$
    of calibrated values and measuring change in FPR.

    \item \textbf{EWMA forgetting factor $\beta$
    (Eq.~\ref{eq:ewma_variance}):}
    $\beta$ \textbf{[tbd: $\{0.7, 0.8, 0.9,
    0.95\}$]}, selected for fastest stable
    convergence of $\sigma_r^2(t)$ within the
    9\,s minimum RSU zone residence window derived
    in Section~\ref{sec:mobility_amplification}.
    \item \textbf{S3 \ac{TCAM} exhaustion detection
    (Eq.~\ref{eq:rule_s3}):}
    Primary signal is $f_{unauth}$ (blockchain endorsement
    check, Eq.~\ref{eq:unauth_flowmod}); no rate threshold
    required. Detection rate $\approx 100\%$ at $0\%$
    \ac{FPR} confirmed on validation runs; rate-based
    monitoring (Eq.~\ref{eq:density_normalized_rate})
    gives $\approx 3\%$ detection at the same \ac{FPR}
    target under realistic mobility and is excluded from
    the detection path.
    
    \item \textbf{S4 \ac{TCAM} utilisation threshold
    $U_{thresh}$ (Eq.~\ref{eq:rule_s4}):}
    $U_{thresh}$ \textbf{[tbd: $\{0.5, 0.6, 0.7,
    0.8\}$]}, selected for best \ac{MCC} at $\leq 1\%$
    \ac{FPR} on the benign validation split. Benign
    utilisation confirmed $\leq 29\%$ under all simulated
    vehicle densities; attacked \ac{RSU}s saturate to
    $\approx 100\%$. Rate-based S4 gives $\approx 2\%$
    detection at the same \ac{FPR} target and is excluded.
    Attacker identified as $v_{atk} = \arg\max_{v'}
    \lambda_{PI}(v',r,t)$ concurrently with detection.
    

    \item \textbf{Standard deviation multiplier
    $k$ (Eq.~\ref{eq:rule_s1}):}
    $k$ \textbf{[tbd: $\{1, 2, 3\}$]}, selected
    for best MCC at $\leq 1\%$ FPR on the
    validation split; $k{=}3$ (three-sigma rule)
    is the initial candidate~\cite{pukelsheim1994three}.

    \item \textbf{Escalation to LSTM detector:} Immediate escalation occurs when the lightweight anomaly score $\geq 0.5$. This threshold is selected to balance the trade-off between detection sensitivity and OBU computational overhead. During bootstrapping (initial 10\,s window), a periodic trigger is used; thereafter, the LSTM input buffer shifts by one step each 1\,s data cycle for continuous re-evaluation.
    \item \textbf{LSTM input window and stride:} A window length of 10\,s and a stride of 5\,s (50\% overlap) are selected based on the RSU observation window of 9--22\,s derived in Section~\ref{sec:mobility_amplification}. This 10\,s window captures at least one full vehicle transit episode while remaining within the minimum observation bound. The final selection will be validated via a sensitivity analysis evaluated over windows of $\{5, 10, 15, 20\}$\,s on the validation split.
    \item \textbf{Attack configurations:} 8 variants (4 Selective Time Delay, 4 Hidden Forwarding) are modeled at attack percentages of 0\%, 20\%, 40\%, 60\%, 80\%, and 100\%; each run lasts 300\,s. To ensure statistical soundness, 5 fixed pseudorandom seeds per configuration are utilized, seeded as \texttt{srand($s_i$)} with $s_i \in \{1,2,3,4,5\}$, each producing a distinct SUMO mobility realisation.

    Attack percentage denotes the fraction of total
non-controller network nodes (200 vehicles $+$
64 RSUs $= 264$ nodes) that are attackers.
At attack percentage $p$,
$\lfloor 0.01p \times 264 \rfloor$ attacker nodes
are allocated among vehicles and RSUs according
to the attack variant: control-plane variants
(Attacks~1, 3, 5, 7) assign attackers among RSUs;
data-plane variants (Attacks~2, 4, 6, 8) assign
attackers among vehicles and RSUs.
Controller attacker count is governed separately:
$p < 33\%$ assigns 1 attacker controller;
$33\% \leq p < 66\%$ assigns 2;
$p \geq 66\%$ assigns 3; $p = 100\%$ assigns all 4.

\end{itemize}

\subsection{Federated LSTM anomaly detector}
\label{sec:fed_lstm}
The Federated LSTM parameters are structured as follows:

\begin{itemize}
   \item \textbf{Data collection:}
    Per-packet features
    $[\delta_t,\, \lambda_{PI,t},\, U_{TCAM,t},\,
    \mathbb{1}[\pi_{delay}{=}\bot],\,
    \mathbb{1}[\pi_{hop}{=}\bot],\,
    \rho_t,\, \bar{v}_t]$
    (Eq.~\ref{eq:lstm_input}) logged per \ac{RSU}
    from NS-3 runs across all attack variants and
    percentages, yielding $6 \times 8 \times 5 = 240$
    labelled run-instances.
    $\mathbb{1}[\pi_{delay}{=}\bot]$ and
    $\mathbb{1}[\pi_{hop}{=}\bot]$ are derived from
    \ac{STARK} proof verification outcomes
    (Eqs.~\ref{eq:stark_delay_verify},
    \ref{eq:stark_hop_verify}) logged by the
    cryptographic layer during each run.
    \item \textbf{Preprocessing:} Z-score normalisation is applied per RSU using benign-only (0\% attack) statistics across a sliding window of 10\,s with a 5\,s stride.
    \item \textbf{Train/validation/test split:} A 70\% / 15\% / 15\% split is enforced, partitioned by seed and stratified by attack variant and percentage.
    \item \textbf{LSTM architecture:} The model consists of 2 stacked LSTM layers (64 and 32 units respectively) followed by a fully connected sigmoid output layer.
    \item \textbf{Training hyperparameters \& Grid Search Space:} To mitigate client drift under non-IID traffic distributions and optimize federated convergence across the 64 RSUs, the hyperparameter configuration is determined via an explicit grid search on the validation split. The search space targets the optimization of the Matthews Correlation Coefficient (MCC) under a strict constraint of $\le$1\% False Positive Rate (FPR):
    \begin{itemize}
        \item \textit{Learning rate (Adam):} $\eta \in \{10^{-4}, 10^{-3}, 10^{-2}\}$
        \item \textit{Batch size:} $B \in \{32, 64, 128\}$
        \item \textit{Local epochs per FedAvg round ($E$):} $E \in \{1, 3, 5\}$ (tuned to balance local underfitting within the short 9--22\,s RSU window against weight drift).
        \item \textit{Global aggregation rounds ($R$):} $R \in \{50, 100, 150\}$, monitored via global loss convergence curves.
    \end{itemize}
    \item \textbf{LSTM detection threshold $\theta^{(k)}$ (Eq.~\ref{eq:lstm_threshold}):}
        $z = 2.326$ ($z_{0.99}$, one-tailed Gaussian approximation~\cite{casella2002statistical}), targeting $\leq 1\%$ FPR per RSU; validated on the validation split for $\geq 95\%$ precision jointly with $\leq 1\%$ FPR. Sensitivity to the Gaussian approximation evaluated by comparing empirical FPR against the 1\% target on the held-out test set.
    \item \textbf{BRFA-v2 trust gate and hash verification         (Algorithm~\ref{alg:brfa_v2}, Steps~1--2):} Trust gate applies $T_{min}$
            (Eq.~\ref{eq:quarantine}) as the RSU eligibility threshold before Krum filtering; aggregation aborts if eligible RSUs $< 2f{+}1$. Hash verification confirms $H(\mathbf{W}_{local}^{(k)})$ matches the blockchain-committed hash
            (Eq.~\ref{eq:bc_model_verify}) for each eligible RSU before the Krum step.
            $T_{min}$ sweep: [tbd: $\{0.3, 0.5, 0.7\}$] as defined in the permissioned blockchain subsection above.
    \item \textbf{Krum outlier rejection threshold $\gamma$ (Eq.~\ref{eq:fed_robust}):}
        $\gamma$ \textbf{[tbd: $\{1.5,\, 2.0,\, 2.5,\,3.0\}$]} (multiples of coordinate-wise standard deviation of submitted weight updates across RSUs), selected
        for best M8 ($\Delta_{poison}$ at $\rho_{mal} = 0.25$) on the validation split under ablation AB5 (Section~\ref{sec:ab5}).


\end{itemize}

\subsection{Hybrid cryptographic integrity layer}
The cryptographic parameters are designed to support strict post-quantum resilience and evaluate real-time vehicular execution bounds:

\begin{itemize}
    \item \textbf{Signature scheme:} CRYSTALS-Dilithium5 (NIST ML-DSA-87) asymmetric signatures ($4,595$\,bytes per signature) are deployed to guarantee quantum-safe entity authentication and data origin integrity at the network layer.
    
    \item \textbf{Batch size $B$
    (Eq.~\ref{eq:overhead_batch}):}
    $B$ \textbf{[tbd: $\{10, 15, 20\}$\,pkts/batch]},
    selected for lowest $L_{e2e}$ within the
    $\leq 50$\,ms RSU processing budget under
    ML-DSA-87 signature size of 4{,}595\,bytes;
    confirmed empirically against measured RSU
    CPU utilisation.
   

    \item \textbf{ZKP scheme:} Post-quantum Zero-Knowledge Proofs for timing
    compliance ($\pi_{\text{delay}}$) and next-hop legitimacy ($\pi_{\text{hop}}$)
    are modelled using SHA3-512 as the underlying hash primitive (NIST Security
    Level~5). Proof generation overhead is modelled as $\leq$10\,ms per packet;
    proof size is modelled as $\leq$100\,KB per
  verification cycle (FRI-STARK, conservative
  blowup; consistent with Eqs.~\ref{eq:stark_delay},
    \ref{eq:stark_hop}).
 Proof outcome
    is determined by the constraint $t_{\text{fwd}} - t_{\text{recv}} \leq
    \Delta_{\max}$.
    
\item \textbf{Time reference commit interval
    $T_{sync}$ (Eq.~\ref{eq:time_consensus}):}
    $T_{sync}$ \textbf{[tbd:
    $\{0.5, 1.0, 2.0\}$\,s]}, co-swept with
    blockchain block interval to assess joint
    impact on $L_{e2e}$. RSU clocks assumed
    GPS-synchronised (sub-millisecond accuracy).
    All Signature~S1 and~S2 comparisons are
    evaluated against $T_{ref}(t)$
    (Eq.~\ref{eq:delay_updated}).

    \item \textbf{Lightweight mode:} HMAC-SHA3-512 ($64$\,bytes) is selectively deployed for highly resource-constrained OBUs to bypass post-quantum asymmetric signature overhead during non-critical broadcast periods.
    \item \textbf{Endorsement threshold:} Verified over $f+1$ RSU endorsements per FlowMod modification, where the fault-tolerance parameter $f$ is swept systematically against varying adversarial node density configurations.
\end{itemize}

\subsection{Permissioned blockchain audit trail}
The permissioned ledger maintains a highly compressed transactional history with parameters optimized via network sensitivity analysis:

\begin{itemize}
    \item \textbf{Block interval sensitivity:} Evaluated across a parametric range of $\{0.5, 1.0, 2.0\}$\,s to assess the trade-off between blockchain propagation delay and end-to-end attack mitigation latency.
    \item \textbf{Block throughput capacity:} Maximum transactions per block scaled up to 50; sensitivity bounds are adjusted based on the peak network event rates captured under 100\% attack intensities.
    
    \item \textbf{Vehicle/RSU trust reward $\Delta_r$,
    penalty $\Delta_p$
    (Eq.~\ref{eq:trust_update}):}
    Both \textbf{[tbd]}, swept independently with
    constraint $\Delta_p > \Delta_r$ enforced to
    penalise faster than reward, consistent with
    asymmetric security requirements. Selected for
    minimum false quarantine rate and maximum
    attack detection speed on the validation split.

    \item \textbf{Vehicle/RSU quarantine threshold
    $T_{min}$ (Eq.~\ref{eq:quarantine}):}
    $T_{min}$ \textbf{[tbd: $\{0.3, 0.5, 0.7\}$]},
    swept independently, selected for best
    quarantine accuracy on the validation split.



    
    \item \textbf{Consensus \& Cluster Bounds:} Local RSU-chains execute PBFT consensus within regional clusters. To satisfy the Byzantine fault-tolerance condition ($f < n/3$), the endorsement threshold parameter $f$ is systematically varied against simulated adversarial node densities.

\end{itemize}

\subsection{Witness-Based Forwarding Verification}
\begin{itemize}
    \item \textbf{Byzantine witness fault parameter
    $f$ (Eq.~\ref{eq:bft_penalty}):}
    $f$ \textbf{[tbd: $\{1, 2, 3\}$]}, selected
    for best MCC at $\leq 1\%$ FPR. Under mean
    vehicle density of $\approx 3.1$ vehicles per
    RSU cell, $f{=}1$ (requiring $2f{+}1{=}3$
    witnesses) is the maximum achievable value
    under average density. Sensitivity to vehicle
    density is evaluated across all attack
    percentages.

    \item \textbf{Forwarding deadline $T_{fwd}$
    (Eq.~\ref{eq:nfwd_detect}):}
    Fixed at $\Delta_{max} = 50$\,ms, equal to
    the hop-delay detection threshold, so that
    delay-equivalent non-delivery and Signature~S2
    are triggered at the same boundary.

    \item \textbf{Observation window $W$
    (Eq.~\ref{eq:dup_alert_cond}):}
    $W$ \textbf{[tbd: $\{5, 10, 15\}$\,s]},
    selected for best MCC, upper-bounded by the
    9\,s minimum RSU zone residence time from
    Section~\ref{sec:mobility_amplification}.
\end{itemize}


\subsection{Multi-Controller Trust Management}
\begin{itemize}
    \item \textbf{Controller trust threshold
    $T_{min}^{ctrl}$ (Eq.~\ref{eq:sc_revoke}):}
    $T_{min}^{ctrl}$ \textbf{[tbd:
    $\{0.3, 0.5, 0.7\}$]}, swept independently
    of $T_{min}$, selected for lowest false
    controller revocation rate.

    \item \textbf{Controller trust reward
    $\Delta_r^{ctrl}$ and penalty
    $\Delta_p^{ctrl}$
    (Eq.~\ref{eq:ctrl_trust_update}):}
    Both \textbf{[tbd]}, swept independently
    of $\Delta_r$ and $\Delta_p$, with
    $\Delta_p^{ctrl} > \Delta_r^{ctrl}$ enforced.
    Selected for minimum time-to-revocation of a
    compromised controller while minimising false
    revocations.

    \item \textbf{Failover latency:}
    Measured from \textsc{SC.Revoke}
    (Eq.~\ref{eq:sc_revoke}) to completion of RSU
    reassignment (Eq.~\ref{eq:ctrl_failover})
    across all affected RSUs, evaluated against
    the 100\,ms safety-critical bound.
\end{itemize}

Simulation settings are summarized in Table \ref{simulation_table}.

\begin{longtable}{|m{5.0cm}|m{9.0cm}|}
\caption{Simulation settings.}
\label{simulation_table}\\
\hline
\textbf{Parameter} & \textbf{Value} \\
\hline
\endfirsthead

\multicolumn{2}{c}%
{{\bfseries Table \thetable\ continued from previous page}} \\
\hline
\textbf{Parameter} & \textbf{Value} \\
\hline
\endhead

\hline
\multicolumn{2}{r}{{Continued on next page}} \\
\endfoot

\hline
\endlastfoot
            ZKP key setup (DKG) &
One-time ceremony among 64 RSUs pre-simulation
 (Eq.~\ref{eq:vk_commit}); no controller
 participation; key rotation evaluated on
 RSU revocation
 (Eq.~\ref{eq:key_rotation_trigger}) \\
          \hline
            Network simulation tool & Network Simulator version 3.35 \\
            \hline
            Mobility trace extraction & SUMO, OpenStreetMap \\
            \hline
            Map & Los Angeles, California, USA (2061 m $\times$ 2137 m) \\
            \hline
            Architecture & SDVN (Vehicle $\rightarrow$ RSU $\rightarrow$ Controller, no LTE) \\
            \hline
            Number of nodes & 268 (200 vehicles, 64 RSUs, 4 controllers) \\
            \hline
            Vehicle types & 100 cars, 25 buses, 25 lorries, 25 vans, 25 trucks \\
            \hline
            RSU placement & 8$\times$8 grid, 260 m horizontal and 270 m vertical spacing \\
            \hline
            Maximum vehicle speed & 150 kmph \\
            \hline
            Data plane & IEEE 802.11p DSRC (5.9 GHz, 10 MHz channel) \\
            \hline
            Control plane & Ethernet / CSMA (1000 Mbps) \\
            \hline
            Transmission power, gain & 41 dBm, 0 dB Tx/Rx gain \\
            \hline
            Max reliable DSRC range & 270 m (Cost231 propagation loss) \\
            \hline
            Propagation models & Cost231 (loss, urban), Constant-speed (delay) \\
            \hline
            Data transmission / routing frequency & 1 Hz \\
            \hline
            Flow packet type, size & UDP, 750 bytes \\
            \hline
            Attack variants evaluated & 8 (4 Selective Time Delay, 4 Hidden Forwarding; control \& data plane) \\
            \hline
            Attacker allocation (vehicles/RSUs) &
              $\lfloor 0.01p \times 264 \rfloor$ nodes;
              plane-specific per attack variant \\
                \hline
             Attacker allocation (controllers) &
               $<33\%$: 1; $33$--$66\%$: 2;
              $\geq 66\%$: 3; $100\%$: 4 \\
            
            \hline
            Attack percentages, seeds & 0\%, 20\%, 40\%, 60\%, 80\%, 100\%; 5 fixed pseudorandom seeds each (300 s/run) \\
            \hline
            Detection architecture & Dual-mode: lightweight rule-based (OBU) + federated LSTM (RSU edge); escalation at score $\geq$0.5 (bootstrapping: periodic at 10\,s; steady-state: sliding buffer, 1 step per 1\,s cycle) \\
            \hline
            $\delta_0$, $\alpha_\rho$, $\alpha_v$ &
  Calibrated from benign SUMO traces;
  robustness swept $\pm\{10,20,30\}\%$ \\
  \hline
 $\beta$ (EWMA forgetting factor) &
   [tbd: $\{0.7, 0.8, 0.9, 0.95\}$] \\
   \hline
 $k$ (S1 threshold multiplier) &
   [tbd: $\{1, 2, 3\}$]; initial candidate $k{=}3$ \\

            \hline
            Rule engine thresholds & Hop-delay $= 50$\,ms; forwarding-ratio deviation $= 20$\% \\
            \hline
            LSTM architecture & 2-layer LSTM (64, 32 units) + FC sigmoid output \\
            \hline
            Sequence window, stride & 10\,s window, 5\,s stride (50\% overlap) \\
            \hline
            Train/val/test split & 70\% / 15\% / 15\% (by seed, stratified by attack class \& \%) \\
            \hline
            LSTM detection threshold $\theta^{(k)}$ (Eq.~\ref{eq:lstm_threshold}) &
            $z = 2.326$ ($z_{0.99}$, one-tailed Gaussian); $\theta^{(k)} = \mu_{\mathcal{A}}^{(k)} + 2.326 \cdot \sigma_{\mathcal{A}}^{(k)}$;
            validated for $\geq$95\% precision and $\leq$1\% FPR jointly on validation split \\            
            \hline
              $\gamma$ (Krum outlier rejection threshold, Eq.~\ref{eq:fed_robust}) & [tbd: $\{1.5, 2.0, 2.5, 3.0\}$] multiples of coordinate-wise std.\ of weight updates;   selected for best poisoning resistance (M8) under 25\% malicious RSU fraction (AB5) \\
            \hline
              BRFA-v2 trust gate threshold (Step~1, Alg.~\ref{alg:brfa_v2}) &              $T_{min}$ [tbd: $\{0.3, 0.5, 0.7\}$] (shared with blockchain quarantine threshold, Eq.~\ref{eq:quarantine}); abort if eligible RSUs $< 2f{+}1$ \\
             \hline
              BRFA-v2 hash verification (Step~2, Alg.~\ref{alg:brfa_v2}) &              SHA3-512 hash equality check (Eq.~\ref{eq:bc_model_verify}); weights failing
              verification excluded before Krum filtering \\
              \hline
            Cryptographic layer & ML-DSA-87 (FIPS\,204) signatures (4,595\,B/pkt),
            HMAC-SHA3-512 lightweight mode (64\,B/pkt),
            batch verification $B$ [tbd: $\{10, 15, 20\}$\,pkts/batch];
            selected for lowest $L_{e2e}$ within 50\,ms,  ZKP proof size  $\leq 100$\,KB per proof   (FRI-STARK, SHA3-512; conservative blowup
             consistent with Eqs.~\ref{eq:stark_delay},
            \ref{eq:stark_hop}) \\

            \hline
            $T_{sync}$ (time reference commit interval) & [tbd: $\{0.5, 1.0, 2.0\}$\,s]; co-swept with block interval to assess joint impact on $L_{e2e}$; RSU clocks GPS-synchronised (sub-ms accuracy) \\
            \hline
            Blockchain audit trail & Permissioned, PBFT consensus, 268 peers (4 controllers + 64 RSUs), 1\,s block interval, $\leq$50 tx/block \\
            \hline
            $\Delta_r$, $\Delta_p$ (vehicle/RSU trust reward/penalty) & [tbd]; $\Delta_p > \Delta_r$ enforced \\
            \hline
            $T_{min}$ (vehicle/RSU quarantine threshold) & [tbd: $\{0.3, 0.5, 0.7\}$] \\
            \hline
            $f$ (Byzantine witness fault parameter) & [tbd: $\{1, 2, 3\}$]; $f{=}1$ expected under mean vehicle density \\
            \hline
            $T_{fwd}$ (forwarding deadline) & 50\,ms ($= \Delta_{max}$) \\
            \hline
            $W$ (witness observation window) & [tbd: $\{5, 10, 15\}$\,s]; $\leq 9$\,s zone residence bound \\
                
    
             \hline
            $T_{min}^{ctrl}$ (controller trust threshold,
            Eq.~\ref{eq:sc_revoke}) &
            [tbd: $\{0.3, 0.5, 0.7\}$]; swept
            independently of $T_{min}$; selected for
            lowest false controller revocation rate \\
            \hline
            $\Delta_r^{ctrl}$, $\Delta_p^{ctrl}$
            (controller trust reward/penalty,
            Eq.~\ref{eq:ctrl_trust_update}) &
            [tbd]; swept independently of $\Delta_r$,
            $\Delta_p$; $\Delta_p^{ctrl} >
            \Delta_r^{ctrl}$ enforced \\
            \hline
            Failover latency target
            (Eq.~\ref{eq:ctrl_failover}) &
            $\leq 100$\,ms from \textsc{SC.Revoke}
            to RSU reassignment completion \\
            \hline
            Controller compromise scenarios &
            0, 1, 2 compromised controllers evaluated;
            PBFT tolerates $f < n/3$ ($\leq 1$
            Byzantine controller) \\
             \hline
           TVR target (M2, Eq.~\ref{eq:tvr}) &
            TVR $\approx 0$ under no attack;
            evaluated per variant and attack percentage;
            post-quarantine return to $\approx 0$
            confirms mitigation effectiveness \\

          \hline
          UCR target (M3, Eq.~\ref{eq:ucr}) &
            UCR $= 0$ under no attack;
            UCR $> 0$ under Attacks~5--8;
            post-quarantine return to 0 confirms
            mitigation; only metric capturing passive
            hidden forwarding (Attacks~7--8) \\
        
  \hline
  $U_{thresh}$ (S4 TCAM utilisation threshold,
  Eq.~\ref{eq:rule_s4}) &
  [tbd: $\{0.5, 0.6, 0.7, 0.8\}$]; benign
  utilisation confirmed $\leq 29\%$; selected for
  best MCC at $\leq 1\%$ FPR on validation split \\
  \hline
  Rate-based estimator $\hat{\lambda}_a$
  (Eq.~\ref{eq:density_normalized_rate}) &
  Excluded from S3 and S4 detection paths;
  retained as theoretical motivation showing
  $\approx 3\%$ / $2\%$ DR under realistic mobility
  for S3 / S4 respectively at the 99th-percentile
  rate threshold \\

\end{longtable}

%...
\section{Simulation Results}
The simulation results for the Selective Time Delay Attack on the data 
plane under NS-3 are presented in Figure~\ref{fig:sim_attack2}.

\begin{figure}[H]
    \centering
    \includegraphics[width=0.95\textwidth]{S.Delay.eps}
    \caption{NS-3 Simulation Output: Selective Time Delay Attack 
    (Data Plane) with Signature Detection.}
    \label{fig:sim_attack2}
\end{figure}

As shown in Figure~\ref{fig:sim_attack2}, the test network simulation 
demonstrates the framework under a data plane selective time delay 
attack, where a malicious RSU injects an 80\,ms buffering delay. 
Signature~S2 is triggered when the hop delay exceeds the 50\,ms 
threshold, achieving an MCC of 1.0 and DR of 100\% at 0\% false 
positive rate, with an average mitigation latency of 50\,ms. The 
attack reduces the packet delivery ratio to 48.89\%, confirming 
the impact of the delay injection on end-to-end packet delivery.

The simulation results for the Passive Hidden Forwarding Attack on the 
data plane under the NS-3 test network are presented 
in Figure~\ref{fig:sim_attack8}.

\begin{figure}[H]
    \centering
    \includegraphics[width=0.95\textwidth]{Hidden_passive.eps}
    \caption{NS-3 Test Network Simulation: Passive Hidden Forwarding 
    Attack(Data Plane) }
    \label{fig:sim_attack8}
\end{figure}

As shown in Figure~\ref{fig:sim_attack8}, the test network simulation 
demonstrates the framework under a passive hidden forwarding attack, 
where a malicious RSU (Node~3) intercepts each packet, forwards the 
original to the legitimate destination while simultaneously sending a 
hidden duplicate to an unauthorized vehicle (Node~2), confirming 
successful eavesdropping. Despite the attack, the packet delivery 
ratio remains at 100\% since the original packet path is unaffected, 
with an average latency of 2.43\,ms. Variant~7 achieves an MCC of 
0 with one false negative, reflecting the passive nature of the 
attack where no observable network degradation occurs, with an 
average mitigation latency of 50\,ms.

\subsection{Federated LSTM Validation Results}
\label{sec:lstm_validation}

The federated LSTM detector (Section~\ref{sec:fed_lstm}) was trained and
evaluated on the full $6\times8\times5=240$ run-instance dataset (1,399,040
feature rows across 64 RSUs and all 8 attack variants plus benign baseline runs, following
the 70/15/15 seed-partitioned split, 10\,s/5\,s windowing, and hyperparameter
grid search specified above. Detection quality is reported on the held-out
test split (seed 5, never used in training or validation).

\begin{table}[H]
\centering
\caption{Federated LSTM Detection Quality (M1--M3) per Attack Variant, Held-Out Test Split}
\label{tab:lstm_detection}
\renewcommand{\arraystretch}{1.1}
\footnotesize
\begin{tabular}{|p{4.2cm}|p{1.6cm}|p{1.6cm}|p{1.6cm}|}
\hline
\textbf{Attack} & \textbf{M1 MCC} & \textbf{M2 DR (\%)} & \textbf{M3 FPR (\%)} \\
\hline
A1 CP-Selective Delay & 0.720 & 83.26 & 5.11 \\
\hline
A2 DP-Selective Delay & 0.893 & 92.87 & 2.84 \\
\hline
A3 CP-TCAM & 0.259 & 90.85 & 52.95 \\
\hline
A4 DP-TCAM & 0.434 & 9HF-
7.74 & 44.51 \\
\hline
A5 Active HF Attack (Control plane) & 0.544 & 70.10 & 6.22 \\
\hline
A6 Active HF Attack (Data Plane) & 0.505 & 63.99 & 6.65 \\
\hline
A7 Passive HF Attack (Control plane) & 0.513 & 68.09 & 6.03 \\
\hline
A8 Passive HF Attack (Data plane) & 0.487 & 64.71 & 6.08 \\
\hline

\textbf{Overall (A1--A8)} & \textbf{0.526} & \textbf{84.52} & \textbf{15.84} \\
\hline
\end{tabular}
\end{table}

\begin{figure}[H]
    \centering
    \includegraphics[width=0.95\textwidth]{Figure_LSTM_DetectionQuality.eps}
    \caption{Federated LSTM detection quality (M1 MCC, M2 DR, M3 FPR) per attack variant.}
    \label{fig:lstm_detection}
\end{figure}

As shown in Table~\ref{tab:lstm_detection} and Figure~\ref{fig:lstm_detection},
detection quality separates into two regimes: Selective Time Delay and
Hidden Forwarding attacks (A1, A2, A5--A8) combine moderate-to-high MCC
(0.49--0.89) with a low false-positive rate (3--7\%), while the TCAM
exhaustion attacks (A3, A4) achieve very high detection rates (91--98\%) at
a higher false-positive cost (45--53\%) -- a consequence of the ground-truth
labelling any RSU holding even one malicious TCAM entry as attack-positive
for its entire observation window, since installed entries persist. A
consistent pattern holds within both attack families: the data-plane variant
outperforms its control-plane sibling on MCC (A2: 0.893 vs.\ A1: 0.720; A4:
0.434 vs.\ A3: 0.259), likely because data-plane attacks act more directly
on observed traffic, leaving a cleaner per-RSU signal, whereas control-plane
attacks act indirectly through a compromised controller or flow rules,
potentially diluting the signal across additional intermediate steps.

\begin{figure}[H]
    \centering
    \includegraphics[width=0.8\textwidth]{Figure_LSTM_Convergence.eps}
    \caption{Federated training convergence (BRFA-v2): global validation loss vs.\ aggregation round.}
    \label{fig:lstm_convergence}
\end{figure}

Figure~\ref{fig:lstm_convergence} shows the global validation loss
(sample-weighted mean benign reconstruction error across all 64 RSUs)
converging over the $R\in\{50,100,150\}$ round grid. The steep loss
reduction in the first $\sim$20 rounds is consistent with the global model
moving away from its random initialisation toward the shared benign-traffic
manifold, with the long tail to round 150 plausibly reflecting fine-tuning
around per-RSU variation in local mobility and traffic conditions. $R=150$
was selected as the smallest round within the convergence tolerance of the
final loss.

\begin{figure}[H]
    \centering
    \includegraphics[width=0.8\textwidth]{Figure_LSTM_M8_Poisoning.eps}
    \caption{M8 federated model poisoning resistance: BRFA-v2 vs.\ naive FedAvg.}
    \label{fig:lstm_poisoning}
\end{figure}

Figure~\ref{fig:lstm_poisoning} evaluates M8 ($\Delta_{poison}$,
Eq.~\ref{eq:delta_poison}) as the malicious RSU fraction $\rho_{mal}$
increases toward the PBFT bound $f<n/3$. At $\rho_{mal}=0.3$, BRFA-v2 holds
$\Delta_{poison}=0.023$ versus naive FedAvg's $0.081$ -- a $3.5\times$
improvement. Naive FedAvg shows a sharp degradation between $\rho_{mal}=0.1$
and $0.2$ (0.014 $\to$ 0.069, a $5\times$ jump) that BRFA-v2 does not
exhibit: FedAvg averages every submitted update unconditionally, so once
enough poisoned gradients are present they begin to dominate the mean,
producing a tipping point. BRFA-v2's Krum step instead rejects updates that
are statistically distant from the coordinate-wise median \emph{before}
aggregation, so poisoned updates are excluded rather than diluting the
average, giving the smoother, near-linear degradation observed.

\begin{figure}[H]
    \centering
    \includegraphics[width=0.85\textwidth]{Figure_LSTM_MobilityStratifiedMCC.eps}
    \caption{Mobility-stratified MCC (Eq.~\ref{eq:mcc_mobility}) per attack variant, across vehicle density ($\rho$) and speed ($\bar{v}$) bins. Low-$\rho$/low-$\bar{v}$ and low-$\rho$/high-$\bar{v}$ cells are unpopulated in the evaluated trace (see text) and marked n/a.}
    \label{fig:lstm_mobility}
\end{figure}

Figure~\ref{fig:lstm_mobility} validates the mobility-adjusted threshold
(Eq.~\ref{eq:mobility_baseline}): MCC holds or improves at higher density
and speed for most attacks (e.g.\ A1: 0.65 $\to$ 0.82), with A5 and A6 as
the two exceptions that degrade under denser, faster traffic. A plausible
explanation is that for A1, denser and faster traffic increases the volume
of legitimate high-priority packets the mobility-adjusted baseline is
estimated from, sharpening the delay-anomaly signal the selective-delay
attack has to hide within, while A5 and A6 (active hidden forwarding)
degrade under the same conditions because denser traffic may give a hidden
duplicate more legitimate concurrent forwarding activity to blend into,
diluting the count-based duplication signal these signatures rely on. The
low-density bins are unpopulated in the evaluated SUMO trace low density
and low speed do not co-occur in this mobility model, and the
low-density/high-speed stratum contains no attack-positive test windows 
reflecting the evaluated scenario rather than a detection failure.


\chapter{Timeline and Resource Required}

\section{Timeline}

\begin{ganttchart}[
    canvas/.style={draw=none, fill=white},
    hgrid style/.style={draw=black!5, line width=.75pt},
    vgrid={*1{draw=black!5, line width=.75pt}},
    title/.style={draw=none, fill=gray!20},
    title label font=\bfseries\footnotesize,
    title height=1.3,
    bar/.style={draw=none, fill=blue!60},
    bar height=.65,
    bar label font=\scriptsize\color{black!80},
    x unit=0.9cm,
    y unit chart=0.55cm,
    time slot format=isodate-yearmonth,
    time slot unit=month
]{2025-12}{2026-08}

%-------------------------------------------------
% Titles
%-------------------------------------------------
\gantttitle{Research Project Timeline: SDVN Security}{9} \\
\gantttitlecalendar{year, month} \\

%-------------------------------------------------
% Tasks (To-Do List)
%-------------------------------------------------
\ganttbar{Identifying the Attacks}{2025-12}{2025-12} \\

\ganttbar{Literature Review}{2025-12}{2026-01} \\

\ganttbar{Simulation Environment Setup}{2026-01}{2026-01} \\
\ganttbar{Attack Modeling \& Simulation}{2026-01}{2026-02} \\
\ganttbar{Baseline Dataset Collection}{2026-02}{2026-02} \\

\ganttbar{LSTM-based Attack Detection Model}{2026-02}{2026-03} \\
\ganttbar{Federated Learning Integration}{2026-04}{2026-04} \\
\ganttbar{Blockchain \& Cryptographic Trust Layer}{2026-05}{2026-05} \\

\ganttbar{System Integration}{2026-06}{2026-06} \\
\ganttbar{Performance Evaluation}{2026-07}{2026-07} \\
\ganttbar{Thesis Writing \& Finalization}{2026-07}{2026-08} \\


\end{ganttchart}



\section{Resource Required}
To execute the high-fidelity simulations and train the Federated LSTM models, high-performance computing resources are required.

Since this research will be done using leexisting university computing infrastructure, there is no direct financial cost.

\chapter{Conclusion}

Software Defined Vehicular Networks (SDVNs) play a significant role in modern intelligent transport systems. They offer flexible, efficient, and adaptive communication among vehicles. Although the programmable and centrally controlled aspects of SDVNs make them very potent, they also introduce a significant security issue. Among all the other attacks, Selective Time Delay Attack and Hidden Forwarding Attack are considered extremely harmful. They target communication that is critical for safety without properly resulting in network malfunctions. Hence, traditional security solutions fail to identify such attacks. 

The proposed research study addresses all these problems with a comprehensive security framework that consists of Blockchain, cryptography techniques, and Federated LSTM Anomaly Detection. Contrary to all other proposed systems, which always work on a Centralized Architecture, the new framework will operate with a Zero-Trust Approach. This would make it possible for the system to identify as well as authenticate malicious behavior from any element in the network, including controllers, RSUs, or vehicles, without relying on a central point of trust.

The proposed approach aims at providing a secured environment by protecting privacy while at the same time providing a low-latency service for vehicular safety applications. The proposed approach will be demonstrated on a realistic simulation environment for SDVN in order to test the performance on a variety of attack scenarios. 

Finally, it can be concluded that this paper attempts at offering a privacy-preserving security platform that is also dependable for next-generation Software Defined Vehicular Networks. With the ability to counter Selective Time Delay Attacks as well as Hidden Forwarding attacks within the same framework, it will enable the development of safer vehicles on the road.


% References
\renewcommand{\bibname}{References}
\bibliographystyle{ieeetr}
\addcontentsline{toc}{chapter}{References} % Add to table of contents
\bibliography{bibliography} 
%\printbibliography

\end{document} 